\documentclass[11pt,a4paper]{article}
\usepackage[margin=29mm]{geometry}
\usepackage{amsmath,amssymb,amsthm,mathtools,microtype}
\usepackage{enumitem,booktabs}
\usepackage{xcolor}
\usepackage{graphicx}
\usepackage[colorlinks=true,linkcolor=blue!50!black,citecolor=blue!50!black,urlcolor=blue!50!black]{hyperref}
\newtheorem{theorem}{Theorem}[section]
\newtheorem{proposition}[theorem]{Proposition}
\newtheorem{lemma}[theorem]{Lemma}
\newtheorem{corollary}[theorem]{Corollary}
\theoremstyle{remark}\newtheorem{remark}[theorem]{Remark}
\newcommand{\sn}{\operatorname{sn}}
\newcommand{\cn}{\operatorname{cn}}
\newcommand{\dn}{\operatorname{dn}}
\newcommand{\atanh}{\operatorname{atanh}}
\newcommand{\atan}{\operatorname{atan2}}
\newcommand{\ind}{\mathbf{1}}
\newcommand{\ee}{\mathbf e}
\newcommand{\ep}{\mathbf e_{\perp}}
\title{Scalar inversion and geodesic counts\\by winding on $\mathrm{SE}(2)$}
\author{Igor Moiseev\thanks{\texttt{moiseev.igor@gmail.com}}}
\date{September 2026}
\begin{document}
\maketitle
\begin{abstract}
We count normal geodesics of the standard sub-Riemannian problem on
$\mathrm{SE}(2)$, retaining those that continue past their cut time and
distinguishing them by heading winding. The rotating endpoint map admits
an exact scalar inverse: the modulus and amplitudes are algebraic
functions of one coordinate, and an elliptic level equation selects the
actual periods. At strict headings every rotating fiber is finite, with
locally bounded cardinality even across the target seam.
An exact separatrix surface partitions the zero-winding fiber into one
oscillating source, one separatrix source, or a finite rotating fiber;
the last has odd cardinality at regular targets. Fixed-heading convexity
gives sharp radial existence intervals, and attained angular thresholds
classify the exterior rays supporting each positive period.
At $|\tan(\theta/2)|=3/4$, a computer-assisted proof identifies a unique
first threshold fold and an explicit radial count window with ordered
post-cut travel times. A separate target has exactly nine regular
rotating sources. Along a specified target segment the complete count
is $7/8/9/10/11$, with two ordinary fold values and strict time ordering.
For general $V>1,W>0$, an analytical comparison orders momenta within
a signed period family and clocks across a nondegenerate connected
sublevel component. Under simultaneous endpoint and effective-period
scaling, an explicit limiting modulus determines eventual counts and
clock offsets for two selected families. On the critical surface a
separate expansion classifies the counts unless both leading coefficients
vanish. Certified examples realize all four nondegenerate outcomes and
show that joint initial existence need not imply persistence.
On the half-heading ray the rotating count is exact at every positive
radius and grows as the inverse square of a vanishing heading, with an
explicit coefficient and bounded remainder. At fixed strict heading,
the zero-winding count along a spatial ray has three large-distance
regimes, proportional to $R$, to $R/\log R$, or equal to one. Every
nonzero winding class has exactly one oscillating source at every radius,
so the full fiber is countably infinite. Headings zero and $\pi$ have
separate exact catalogues, including an infinite rotating fiber on the
nonzero vertical axis. The minimizing conclusions use the known optimal
synthesis. General eventual scale thresholds remain unquantified.
\end{abstract}
\noindent\textit{Keywords:} sub-Riemannian geometry, group of motions of the
plane, exponential map, geodesic counting, Jacobi elliptic functions.\\
\textit{MSC 2020:} 53C17, 49K15, 33E05.

\section{Problem, scope, and main statement}

We declare the left-invariant vector fields
\[
 X_1=\cos\theta\,\partial_x+\sin\theta\,\partial_y,
 \qquad X_2=\partial_\theta
\]
orthonormal, and start trajectories at $(0,0,0)$. Thus
$\dot x=u_1\cos\theta$, $\dot y=u_1\sin\theta$, $\dot\theta=u_2$,
with unit-speed normalization $u_1^2+u_2^2=1$.
The normal Hamiltonian system reduces to a pendulum. Its rotating stratum
is denoted by $C_2$, following \cite{MS,SC}. A source means an initial
normal covector of Hamiltonian energy $1/2$ together with a positive
travel time; $C_2$ is specified when restricting to rotating sources.
We count distinct sources, without identifying reflection pairs or
discarding trajectories after their cut time. This convention is also
equivalent to counting the corresponding constant-speed parametrized
normal geodesics on a fixed time interval.

Throughout, elliptic integrals use the \emph{parameter} $m=k^2$:
\[
 K(m)=\int_0^{\pi/2}\frac{d\phi}{\sqrt{1-m\sin^2\phi}},\quad
 E(m)=\int_0^{\pi/2}\sqrt{1-m\sin^2\phi}\,d\phi.
\]
We write $L(m)=4(K(m)-E(m))$ for the longitudinal drift per period.
The corresponding incomplete integrals are $F(\phi\mid m)$ and
$E(\phi\mid m)$, continued as real integrals in $\phi$.
Set $\ee(a)=(\sin a,-\cos a)$ and $\ep(a)=(\cos a,\sin a)$.

\begin{theorem}[Far-field rotating count]\label{thm:main}
Fix $a_0,b_0\in(-\pi/2,\pi/2)$ with $\theta=a_0-b_0\ne0$, and put
\[
 q_R=(R\ee(a_0),\theta),\qquad
 m_*=\max(\sin^2a_0,\sin^2b_0),\qquad L_*=L(m_*).
\]
There is $R_0=R_0(a_0,b_0)$ such that, for $R>R_0$, the number
$\nu_2(q_R)$ of positive-time $C_2$ sources is finite and
\begin{equation}\label{eq:main}
 \nu_2(q_R)=\frac{4R}{L_*}+O(1),\qquad R\longrightarrow\infty.
\end{equation}
The same statement holds for $(-R\ee(a_0),\theta)$ by spatial negation.
The constants depend on the fixed angles. No effective value of $R_0$
or uniform estimate toward the excluded angle boundaries is asserted.
\end{theorem}

The proof distinguishes $\sin^2a_0\ne\sin^2b_0$ from $b_0=-a_0$.
In the first case an exact eventual floor formula is given in
Theorem~\ref{thm:generic}. In the second, omitting the zero-transverse
fiber would lose half of the coefficient in \eqref{eq:main}.
The finite quantity in Theorem~\ref{thm:main} is the rotating count.
The next theorem describes its place in the full fiber. Neither theorem
includes $\theta=0$.

\begin{theorem}[Normal fibers by heading winding]\label{thm:winding}
For the angles and targets of Theorem~\ref{thm:main}, for all sufficiently
large $R$, the full positive-time normal fiber has the following form.
With continuous heading $\widetilde\theta(0)=0$, define
$\ell\in\mathbb Z$ by $\widetilde\theta(t)=\theta+2\pi\ell$.
The class $\ell=0$ consists precisely of the finitely many $C_2$ sources,
with count \eqref{eq:main}. Each class $\ell\ne0$ contains exactly one
source, belonging to $C_1$. There are no $C_3,C_4,C_5$ sources.
The full fiber is countably infinite; every $C_1$ source is regular.
Finiteness of the class $\ell=0$ and the description of the classes
$\ell\ne0$ hold in fact at every radius $R>0$
(Corollary~\ref{cor:seam-finite}, Propositions~\ref{prop:uniform-finite}
and~\ref{prop:global-winding}); the count \eqref{eq:main} and the
exclusion of $C_3$ are the eventual clauses.

If $a_0+b_0\ne0$, the unique minimizing source is strictly pre-cut and
belongs to actual $C_2$ cell zero. If $a_0+b_0=0$, exactly two sources
minimize, both in actual $C_2$ cell zero at their first Maxwell time.
The same conclusions hold for spatially negated targets.
\end{theorem}

The minimizing conclusions use the published optimal synthesis.
The proof in Section~\ref{sec:winding} includes the oscillating covering
and seam arguments and the exclusions of the remaining regimes.
The eventual threshold for the combined theorem remains non-effective;
regularity of every rotating source is not asserted.

\begin{theorem}[All fixed spatial directions at strict heading]\label{thm:directions}
Fix a unit vector $\mathbf e=(e_x,e_y)$ and $0<|\theta|<\pi$.
For $q_R=(R\mathbf e,\theta)$ put
\[
 \delta_0=e_y,\qquad \delta_1=e_y\cos\theta-e_x\sin\theta,\qquad
 \Delta=\delta_0\delta_1 .
\]
For sufficiently large $R$, the zero-winding fiber is
\[
\begin{array}{ccl}
\text{condition}&\text{stratum}&\text{number of sources}\\[2pt]
\Delta>0&C_2&4R/L_*+O(1),\\
\Delta=0&C_2&R/[2(\log(4R/|\sin\theta|)-1)]+O(1),\\
\Delta<0&C_1&1,
\end{array}
\]
where $L_*=L(1-\min(\delta_0^2,\delta_1^2))$ in the first row.
In every row there are no $C_3,C_4,C_5$ sources eventually.
Exactly one source minimizes, except on the interior line
$e_x\cos(\theta/2)+e_y\sin(\theta/2)=0$, where exactly two minimize.
The interior and boundary thresholds are non-effective. For the exterior,
an explicit sufficient threshold after spatial sign normalization is
\eqref{eq:exterior-radius}.

For every $X\ne0$ and $0<|\theta|<\pi$, every nonzero winding class
contains exactly one source, in $C_1$, and it is regular. This last
assertion holds at every radius, including horizontal directions.
\end{theorem}

The interior statement is Theorem~\ref{thm:winding}; the nonhorizontal
boundary and exterior proofs are in Sections~\ref{sec:boundary}
and \ref{sec:exterior}. Section~\ref{sec:horizontal} gives the horizontal
completion and the invariant sign formulation.
Proposition~\ref{prop:global-winding} proves the all-radius winding
assertion. The remainder constants and thresholds in these fixed-angle
statements are not uniform when the angles vary with $R$.
When the terminal angle approaches the rotating boundary together with
the radius, the count passes through exponentially small transitions;
that regime is treated in a separate note in preparation \cite{TF}.

The remaining headings have exact catalogues at every radius:
Section~\ref{sec:pi} treats heading $\pi$ and pure orientation,
and Section~\ref{sec:zero} treats heading zero, including its coordinate
axes. These exact exceptional-heading results are separate from the
eventual strict-heading table. The identity has no positive-time
winding-zero source.

\paragraph{Relation to established results.}
The elliptic endpoint parametrization and reflection symmetries are
standard inputs from Moiseev--Sachkov \cite[Sections 3--5]{MS}.
Sachkov's conjugate-time analysis \cite[Section 2]{SC} concerns local
optimality; the complete conjugate-time set of the rotating geodesics is
determined in \cite{CT}. The optimal synthesis \cite[Theorems 3.1--3.2]{SS} gives one
or two minimizing trajectories on its stated target strata; its source
domain is restricted by the cut time. It therefore does not count all
positive-time sources considered here.
Duits--Boscain--Rossi--Sachkov \cite[Theorems 5--6 and 11]{DBRS}
already reduce the cuspless boundary-value problem to a scalar zero,
using spatial arclength and a source domain ending at the first cusp.
That inverse concerns their cuspless domain, whereas the present
components retain all positive-time rotating sources, including post-cut
cells. Scalar reduction itself is therefore an established method.
Lerario--Rizzi \cite[Example 2, equation (2), and Theorems 2, 7--9]{LR}
derive quantitative counts for contact Carnot groups and lower bounds
under local nilpotent approximation. In particular, the Heisenberg count
already has inverse-square growth in the horizontal displacement at fixed
nonzero vertical coordinate. The exponent two in
Theorem~\ref{thm:ray-small-heading-count} thus has an established model
precedent. Here the limit is taken in the fixed $\mathrm{SE}(2)$ geometry at fixed
positive radius, and the theorem gives its explicit radius-dependent
coefficient and bounded remainder. This limit, and the large-distance
limit in \eqref{eq:main}, differ from their local nilpotent scaling.
Wang--Ku--Bravo-Doddoli \cite[Theorems A--B]{WKB} classify periodic
geodesics and metric lines. Those statements address recurrence and
global minimizing behavior, rather than counts of all sources of a
specified endpoint.

The contributions of this paper are the complete post-cut scalar
components of the rotating inverse (Proposition~\ref{prop:inverse}) and
their global eventual control (Theorems~\ref{thm:generic} and
\ref{thm:boundary-count}), the seam and directional completion
(Sections~\ref{sec:seam} and \ref{sec:horizontal}), the fixed-cell scale
(Proposition~\ref{prop:scale}), and the exact exceptional-heading counts
(Sections~\ref{sec:pi} and \ref{sec:zero}). They also include finiteness
of every nonzero strict-heading rotating fiber, uniformly on compact
target families (Corollaries~\ref{cor:full-finite} and
\ref{cor:seam-finite}, Proposition~\ref{prop:uniform-finite}), the
all-radius half-heading ray count (Theorem~\ref{thm:ray-all-radii}),
and its small-heading asymptotic (Theorem~\ref{thm:ray-small-heading-count}).
The exact rotating image and odd regular-fiber parity are established in
Theorem~\ref{thm:rotating-image-parity}. The degree method is classical
\cite[Section 4]{Milnor}; the all-period image barrier and the properness
needed for this application are proved in Section~\ref{sec:rotating-image}.
Corollary~\ref{cor:zero-winding-regimes} combines that barrier with
the classical synthesis \cite{SS} to classify the entire zero-winding
fiber by pendulum regime at every radius; it asserts no new optimal
synthesis or exact arbitrary-target rotating count.
Section~\ref{sec:radial-existence} derives strict fixed-heading convexity
and the resulting sharp radial existence classification from the same
hyperbolic graph. The half-heading onset is already in
\cite[(2.13), Lemma 2.3]{SS}; these geometric consequences are
distinguished from the all-time source-count results.
Section~\ref{sec:exterior-sector} excludes all positive rotating periods
in an explicit exterior sector, giving an exact count there by the
classical pre-cut uniqueness theorem. Proposition~\ref{prop:exterior-three}
uses a complete interval certificate to count three regular rotating
sources at a wider exterior target, including two post-cut sources.
Theorem~\ref{thm:exterior-thresholds} characterizes the sharp angular
support of each positive period by an attained minimum, with period one
determining the all-radius uniqueness sector. It identifies conjugacy
at the threshold without asserting a fold/cusp type or an explicit value
of the minimum. These post-cut support statements use the classical
endpoint formulas and minimizing synthesis as inputs.
Corollary~\ref{cor:angular-threshold-certificate} encloses the first
threshold at $w=3/4$ in an interval of width $3\cdot10^{-15}$ using
a complete interval certificate and separate coverage audit.
Theorem~\ref{thm:threshold-fold-window} identifies a unique attaining
fold for this heading and its local radial count window.
Corollary~\ref{cor:effective-fold-window} supplies an explicit slope
range and certified radial endpoints at its upper cap.
Theorem~\ref{thm:window-time-order} then orders the two post-cut clocks
throughout that window, using classical first variation and a separate
momentum-sign certificate. Propositions~\ref{prop:period-action} and
\ref{prop:period-action-example} give a formal period-action identity
and a complete nine-source rotating fiber with strictly ordered period
groups, including a certified connected comparison between periods one
and two. Proposition~\ref{prop:nine-source-interval} continues that count
along an explicit target segment with exactly two ordinary fold values.
Corollary~\ref{thm:line-time-order} orders all its zero-winding travel times
throughout that segment. Theorem~\ref{thm:general-pair-order} supplies a
general analytical within-family momentum comparison for $V>1,W>0$ and
orders the clocks across a nondegenerate connected sublevel component.
Neither result compares arbitrary branches in different signed period
families. Section~\ref{sec:dilation} studies simultaneous endpoint and
period scaling. Its uniform source limit and strict-threshold theorem
give eventual counts and clock offsets for two selected signed families;
a critical-value expansion treats the equality surface except where
both unfolding coefficients vanish. Certified examples show multiple
critical transitions and failure of joint initial existence to guarantee
persistence. The source formulas and elliptic derivative identities are
classical \cite{MS,SC,DLMF}; the scaling is a formal period interpolation,
not the nilpotent approximation used in the contact counting results
of \cite{LR}. No new special-function identity is claimed.
The heading-zero and heading-$\pi$ source
conditions, and the intervals containing the roots of $f_1$, are
established in \cite[Propositions 5.3--5.6 and Lemma 5.3]{MS}.
The minimizing zero-heading branch and heading-$\pi$ radial equation
occur in \cite[Lemmas 2.9--2.14]{SS}. The additional arguments here
invert the modulus at a fixed target, order positive-cell branches,
identify overlaps and count all-time sources. Source reversal and
the Hamiltonian equations used for pure orientation are also known
inputs.
In particular, the all-period Maxwell roots and their ray condition are
already in \cite[Lemma 5.3, Proposition 5.6 and Theorem 5.2]{MS}.
The first ray threshold is Sachkov's curve $R_1^1(\theta)$, and its
pair of minimizing sources is part of the optimal synthesis
\cite[Lemma 2.5, Table 2 and Theorems 3.1--3.3]{SS}.
Theorem~\ref{thm:ray-all-radii} uses those known Maxwell branches to
invert the radius in every actual period and count the full rotating
fiber. Corollary~\ref{cor:ray-cut-status} identifies the established
minimizing part inside this all-time count.

\paragraph{Notation.}
The following symbols recur; each is also introduced where it first
enters.
\begin{itemize}[nosep,leftmargin=*]
\item $\operatorname{Exp}(\lambda,t)$ is the exponential map, the
 endpoint at time $t$ of the normal geodesic with initial covector
 $\lambda$; a source is a pair $(\lambda,t)$ with $t>0$. A source is
 \emph{regular} if the differential of $\operatorname{Exp}$ is
 nonsingular there, that is, $t$ is not a conjugate time. The cut time
 of a source is $t_{\rm cut}$ \cite{SS}; a source is strictly pre-cut if
 $t<t_{\rm cut}$ and post-cut if $t>t_{\rm cut}$.
\item $\nu_2(q)$ is the number of positive-time $C_2$ sources of $q$,
 and $N_0(q)$ the number of sources of $q$, in all pendulum strata,
 whose heading winding (Theorem~\ref{thm:winding}) is zero. By
 Corollary~\ref{cor:zero-winding-regimes}, $N_0=\nu_2$ throughout
 the exact rotating image and $N_0=1$ on and outside its separatrix boundary.
\item The \emph{cell} of a rotating source is its period index
 $n=\lfloor t/(4kK)\rfloor$ in \eqref{eq:bound}; ``actual cell''
 stresses this integer as opposed to a chart label. For oscillating
 sources the cell is the integer $n$ in $p=t/2=2nK+r$, $|r|<K$
 (Lemma~\ref{lem:c1cells}). In Section~\ref{sec:zero} a rotating source
 carries the label $n$ of $p=t/(2k)=2nK+F(A\mid k^2)$
 (Lemma~\ref{lem:zero-scalar}), which is its period index at the
 central sources $p=2nK$.
\item $s_1=\operatorname{sgn}\cos(\gamma/2)$ and
 $s_2=\operatorname{sgn}c$ are the sign parameters of \cite{MS,SC} in
 the pendulum coordinates $(\gamma,c)$ of Section~\ref{sec:winding};
 the sign $\varepsilon$ of the rotating decomposition
 \eqref{eq:decomposition} corresponds to $s_2$.
\item $f_1(p)=\cn p\,(\mathcal E(p)-p)-\dn p\sn p$ is the function of
 \cite[(5.12)]{MS}; its first positive root $p_1$ (written $p_1^1$ in
 Section~\ref{sec:zero}, where $p_n^1$ is the $n$-th root) gives the
 rotating cut time $2kp_1$ \cite{SS}. It first enters in the proof of
 Theorem~\ref{thm:winding}.
\item $a_X\in(-\pi/2,\pi/2)$ is the angle with $X=R\,\ee(a_X)$, where
 defined (Proposition~\ref{prop:cellbounds}); for the targets of
 Theorem~\ref{thm:main} it is $a_0$.
\item $\eta$ is the scaled cell index $n/R$ in the Remark following
 Proposition~\ref{prop:cellbounds}, and the amplitude
 $\operatorname{am}\tau$ of the midpoint phase $\tau$ in
 Sections~\ref{sec:winding} and \ref{sec:zero}.
\item Travel times are $t$, or $T$ where they are the unknown of a
 discrete equation: $T=|\Theta_\ell|$ in Lemma~\ref{lem:orientation};
 $T_n(R)=2(2n+1)K$ in Theorem~\ref{thm:pi} and
 Proposition~\ref{prop:pi-times}, and $T=4nK$ in
 Theorem~\ref{thm:zero-axes}, are the travel times of the oscillating
 sources of windings $n,-n-1$, respectively $\pm n$; in the proof of
 Theorem~\ref{thm:zero-axes}, $T_n=4nkK$ is the travel time of the
 central rotating pair of cell $n$ on the vertical axis.
\end{itemize}

\section{Exact endpoint decomposition and scalar inverse}
\label{sec:decomposition}

Write a rotating source as $(\psi,k,t,\varepsilon)$, where $0<k<1$,
$\psi\in\mathbb R/(4K(k^2)\mathbb Z)$, $t>0$, and
$\varepsilon\in\{1,-1\}$. With $v=\psi+t/k$, define
\begin{equation}\label{eq:ab}
 a=\arcsin(k\sn(\psi\mid m)),\quad
 b=\arcsin(k\sn(v\mid m)),\quad m=k^2.
\end{equation}
Both angles lie strictly between $-\pi/2$ and $\pi/2$.
The published rotating endpoint formulas \cite[Section 3.3]{MS}
give, by regrouping their two spatial components,
\begin{equation}\label{eq:decomposition}
 \begin{gathered}
 \theta=a-b,\qquad X=\varepsilon\{A\ee(a)+B\ep(a)\},\\
 A=\int_\psi^v m\sn^2(u\mid m)\,du>0,\qquad
 B=k(\cn\psi-\cn v).
 \end{gathered}
\end{equation}
Here and below Jacobi functions without a displayed parameter use $m$.
Indeed $\dn\psi=\cos a$, $k\sn\psi=\sin a$, and the longitudinal
integral equals $v-\psi-\mathcal E(v)+\mathcal E(\psi)$, where
$\mathcal E$ is Jacobi's epsilon function, $\mathcal E'(u)=\dn^2u$
\cite[Section 22.16]{DLMF}. This verifies the decomposition directly.
The heading difference is in $(-\pi,\pi)$ and represents the group heading.
In particular
\begin{equation}\label{eq:bound}
 |B|\le2k<2,\qquad
 A=nL(m)+Q_0,\quad 0\le Q_0<L(m),\quad
 n=\left\lfloor\frac{t}{4kK(m)}\right\rfloor.
\end{equation}
For nonzero heading, $0<Q_0<L(m)$. The phase circle here is $4K$,
even when scalar conjugate-point expressions have a $2K$ symmetry
\cite{CT}.
Changing $\varepsilon$ negates $X$ at unchanged heading.


\begin{proposition}[Scalar inverse off the seam]\label{prop:inverse}
Fix $X=(x,y)\ne0$ and $0<|\theta|<\pi$. For the sign $\varepsilon=1$,
put $z=\tan a$ and define
\begin{align}
 H&=1+z^2,& C&=x+yz,& D&=z\cos\theta-\sin\theta,\notag\\
 d&=D^2-z^2,& U&=C^2+d,&
 A&=\frac{xz-y}{\sqrt H},\quad B=\frac C{\sqrt H},\label{eq:zdata}\\
 f&=\frac U{2C\sqrt H},&
 g&=\frac{d-C^2}{2C\sqrt H},&
 m&=\frac{z^2+(U/(2C))^2}{H}.\notag
\end{align}
Its admissible domain is exactly
\begin{equation}\label{eq:zdomain}
 C\ne0,\quad \cos\theta+z\sin\theta>0,\quad xz-y>0,\quad 0<m<1.
\end{equation}
At such a point let
\begin{align*}
 \phi_0&=\atan(z/\sqrt H,f)\pmod{2\pi},&
 \phi_1&=\atan(D/\sqrt H,g)\pmod{2\pi},\\
 \psi&=F(\phi_0\mid m),&v_0&=F(\phi_1\mid m),
\end{align*}
using representatives in $[0,2\pi)$ for the amplitudes. Set
\begin{equation}\label{eq:residual}
 \begin{split}
 w&=\ind_{\{v_0<\psi\}},\qquad \delta=v_0-\psi+4K(m)w,\\
 Q_0&=\delta-E(\phi_1\mid m)+E(\phi_0\mid m)-4E(m)w,\\
 N(z)&=\frac{A-Q_0}{L(m)}.
 \end{split}
\end{equation}
Every intersection $N(z)=n\in\mathbb Z_{\ge0}$ gives exactly one source
\[
 (\psi,\sqrt m,\sqrt m\,\delta+4n\sqrt m K(m),1).
\]
These exhaust the $B\ne0$ sources of that sign. For $\varepsilon=-1$,
apply the same construction to $-X$.
\end{proposition}

\begin{proof}
The transverse and heading equations imply
$f-g=B$ and $f^2-g^2=\sin^2(a-\theta)-\sin^2a$ for
$f=k\cn\psi$, $g=k\cn v$. Solving gives \eqref{eq:zdata}.
Conversely those formulas satisfy
$f^2+\sin^2a=g^2+\sin^2(a-\theta)=m$.
The second inequality in \eqref{eq:zdomain} says
$\cos(a-\theta)>0$ and selects the required principal terminal angle.
Thus both signed sine--cosine pairs determine unique phases modulo $4K$.
Nonzero heading excludes coincident phases, so $0<\delta<4K$ and
$0<Q_0<L$. The remaining longitudinal equation is exactly $N=n$.
Its reconstructed time is positive and has the stated actual period
index. No division by either cosine occurred, so their zeros are included.
Local phase lifts remove the artificial discontinuities of the chosen
amplitude representatives.
\end{proof}

The admissible domain also has a useful algebraic description:
\begin{equation}\label{eq:factor}
 m-1=\frac{(U-2C)(U+2C)}{4C^2H}.
\end{equation}
Both numerator factors are quadratic in $z$. Their roots and the linear
boundaries in \eqref{eq:zdomain} partition the domain exactly. The
condition $m=0$ cannot occur under the other hypotheses: it would force
$z=0$ and $x^2+\sin^2\theta=0$.

\begin{proposition}[Regularity and folds in this chart]\label{prop:fold}
At a scalar integer intersection with $B\ne0$, the full exponential map
is regular if $N_z\ne0$. If $N_z=0$, it has corank one and is a fold
if and only if $N_{zz}\ne0$.
\end{proposition}
\begin{proof}
The reconstruction in Proposition~\ref{prop:inverse}, now with
independent variables $(a,B,\theta)$, gives local analytic source
coordinates. Write $Z=nL+Q_0$. In the oriented spatial basis
$(\ee,\ep)$, the determinant of
$X=Z\ee+B\ep$ at fixed $\theta$ is
\[
 J=Z_a-B-ZZ_B.
\]
The $B$ column is $Z_B\ee+\ep$, and the heading column has last
component one, so the full rank is at least two.
For a fixed target put $F(a)=A(a)-Z(a,B(a),\theta)$.
Since $A'=B$, $B'=-A$, at a root
$J=-F'=-L(1+z^2)N_z$.
At a stationary root the kernel is
$V=\partial_a-Z\partial_B$; differentiating the displayed determinant
gives $VJ=-F''=-L(1+z^2)^2N_{zz}$ there.
The corank-one fold criterion is precisely nonvanishing of this kernel
derivative. This argument does not extend the chart to $B=0$.
\end{proof}

The decomposition already bounds where sources can lie.

\begin{proposition}[Exact cell bounds]\label{prop:cellbounds}
Let $R=|X|>2$, and let $n\ge1$ be the period count of a rotating source of
$(X,\theta)$ with $\varepsilon=1$.  Then
\begin{equation}\label{eq:cellband}
 \frac{\sqrt{R^2-4}}{n+1}\le L(m)\le\frac Rn .
\end{equation}
If moreover $X=R\ee(a_X)$ with $a_X\in(-\pi/2,\pi/2)$ and
$\varepsilon_R=\arcsin(2/R)$, then
\begin{equation}\label{eq:anglebound}
 |a-a_X|\le\varepsilon_R,\qquad
 k\ge k_*-\varepsilon_R,\qquad k_*=\max\{|\sin a_X|,|\sin(a_X-\theta)|\},
\end{equation}
so that, when $k_*>\varepsilon_R$, only the cells $n\le R/L\bigl((k_*-\varepsilon_R)^2\bigr)$
can contain sources.  For $\varepsilon=-1$ replace $X$ by $-X$.
\end{proposition}
\begin{proof}
Orthogonality of $\ee(a)$ and $\ep(a)$ in \eqref{eq:decomposition} gives
$R^2=A^2+B^2$, and \eqref{eq:bound} gives $nL\le A<(n+1)L$, $|B|<2$; this
is \eqref{eq:cellband}.  Scalar products with $\ee(a)$ and $\ep(a)$ give
$A=R\cos(a-a_X)>0$ and $B=R\sin(a_X-a)$; both angles lie in
$(-\pi/2,\pi/2)$ and the cosine is positive, so $|a-a_X|<\pi/2$ and
$|B|<2$ gives the first bound in \eqref{eq:anglebound}.  Since
$k\ge|\sin a|$, $k\ge|\sin b|$ and $b=a-\theta$, the Lipschitz bound of
the sine gives the second; the cell bound follows from $nL\le A\le R$ and
$L'>0$.
\end{proof}

\begin{proposition}[Finite positive-period fibers]\label{prop:positive-finite}
Let
\[
 \Omega=\{(x,y,\theta): (x,y)\ne0,\quad 0<|\theta|<\pi,
       \quad x\cos(\theta/2)+y\sin(\theta/2)\ne0\}.
\]
For every target in $\Omega$, the number of rotating sources with actual
period index $n\ge1$, across both rotation signs, is finite. These numbers
are uniformly bounded on every compact subset of $\Omega$. Every such
source satisfies the explicit index restriction
\begin{equation}\label{eq:positive-index-bound}
 1\le n<\frac{|X|}{L(\sin^2(\theta/2))}.
\end{equation}
In particular, no such source exists if $|X|\le L(\sin^2(\theta/2))$.
The uniform root-count bound is not effective. These assertions do not
include period index zero or targets on the excluded seam.
\end{proposition}

\begin{proof}
Work first with $\varepsilon=1$ and $a=\arctan z$ on the admissible open
set of Proposition~\ref{prop:inverse}; put $b=a-\theta$ and $R=|X|$.
The circle equations give
\begin{equation}\label{eq:positive-modulus-bound}
 m\ge\max\{\sin^2a,\sin^2b\}\ge\sin^2(\theta/2)=:m_h>0.
\end{equation}
Indeed, at least one of $|a|,|b|$ is at least $|\theta|/2$, and both
angles are principal. At a point with $N\ge\ell>0$, $0<Q_0<L$ gives
\begin{equation}\label{eq:positive-compact-bounds}
 \ell L(m)\le A-Q_0<A\le R,\qquad
 m<L^{-1}(R/\ell)<1,\qquad A>\ell L(m_h)>0.
\end{equation}
Here $L$ is strictly increasing from zero to infinity. In particular,
\eqref{eq:positive-index-bound} follows at $N=n\ge1$.

The positive superlevel $N\ge\ell$ on any maximal admissible component
is compact. To see this, take a sequence in that superlevel and a limit
of its angles in the closure of the principal-angle interval.
An endpoint with $|a|=\pi/2$ or $|b|=\pi/2$ would force $m\to1$,
contrary to \eqref{eq:positive-compact-bounds}. The same bounds exclude
$A\to0$ and $m\to0,1$. Finally the exact circle identity
\[
 \sin^2a-\sin^2b=-B(f+g)
      =\sin\theta\sin(2a-\theta)
\]
excludes $B\to0$: the boundedness of $f,g$ would force $a=\theta/2$,
since $2a-\theta\in(-\pi,\pi)$ even at the principal interval's
endpoints for this fixed strict heading. The limiting equation $B=0$
would then be the excluded target seam. Thus the limit remains
admissible and, by continuity, in the superlevel. It lies in the same
component, since components are relatively closed in the admissible set.

For completeness, $N$ is jointly real analytic in $(a,x,y,\theta)$
on this open set, including the apparent amplitude cuts and the zeros
of $f$ or $g$. The nonzero unit-circle pairs
$(\sin a,f)/\sqrt m$ and $(\sin b,g)/\sqrt m$ admit analytic local
angle lifts. Choose them with $0<\phi_1-\phi_0<2\pi$; equality modulo
$2\pi$ would imply $a=b$, contrary to the nonzero heading. Then
\[
 Q_0=\int_{\phi_0}^{\phi_1}
          \frac{m\sin^2u}{\sqrt{1-m\sin^2u}}\,du
\]
is jointly analytic, as follows by changing to a fixed integration
interval. Simultaneous $2\pi$ shifts of the lifts leave it unchanged.
This also proves the continuity used above without imposing artificial
phase cuts on the domain.

On a maximal component, $N$ cannot be identically a positive constant:
the corresponding superlevel would be the whole nonempty open interval,
contradicting compactness. The analytic identity theorem therefore
makes each positive level set discrete, and compactness makes it finite.
There are finitely many components by \eqref{eq:factor} and the linear
domain boundaries. The index bound leaves finitely many positive
integers. Proposition~\ref{prop:inverse} gives exactly one source per
integer intersection, with residual time in $(0,4kK)$ and actual index
$n$. A source with $B=0$ would lie on the excluded target seam by the
same circle identity. Replacing $X$ by $-X$ covers the other rotation
sign and preserves all these restrictions.

To prove uniformity, let $\mathcal Q\subset\Omega$ be nonempty and
compact, set $R_{\max}=\max_{\mathcal Q}|X|$ and
$m_*=\min_{\mathcal Q}\sin^2(\theta/2)>0$. At every positive integer
intersection,
\[
 n<R_{\max}/L(m_*),\qquad
 m_*\le m<L^{-1}(R_{\max})<1,\qquad A>L(m_*)>0.
\]
The preceding boundary argument, now allowing the target to converge
inside $\mathcal Q$, proves compactness of the joint root set in
$(a,q,n,\varepsilon)$. At each of its points some finite derivative
$\partial_a^j(N-n)$, $j\ge1$, is nonzero; otherwise analyticity would
make $N$ identically $n$ on that component. Joint analyticity keeps this
derivative nonzero on a product neighborhood in angle and target.
Repeated Rolle's theorem bounds the number of roots there for each
fixed target by $j$. A finite cover of the compact root set supplies
a uniform bound by summing these orders. This argument assumes neither
simple roots nor a sign predicate for critical points.
\end{proof}


\begin{corollary}[Period zero and the full rotating fiber]\label{cor:full-finite}
For every $q\in\Omega$ of Proposition~\ref{prop:positive-finite}, the
entire rotating fiber is finite, including actual period index zero.
Its cardinality is uniformly bounded on each compact subset of $\Omega$.
Every period-zero source satisfies
\begin{equation}\label{eq:zero-time-bound}
 |X|\le t<|X|+2\pi.
\end{equation}
The finiteness bound is not effective and gives no exact count.
\end{corollary}

\begin{proof}
Only period zero remains after Proposition~\ref{prop:positive-finite}.
For such a source $t=k\delta$ and
\[
 \delta=A+\int_{\psi}^{\psi+\delta}\dn^2u\,du,
 \qquad 0<\delta<4K.
\]
The last integral is strictly between zero and $4E(m)\le2\pi$.
Since $A\le|X|$ and $k<1$, this gives the upper bound in
\eqref{eq:zero-time-bound}; the lower bound follows from unit speed.

Use the original normalized initial covector $(\gamma,c)$, with
$\gamma$ taken modulo $4\pi$. Its pendulum energy obeys
\[
 \mathsf E=c^2/2-\cos\gamma=2/m-1,
 \qquad |c|\le2/\sqrt m\le2/|\sin(\theta/2)|
\]
by \eqref{eq:positive-modulus-bound}. Thus all period-zero sources of
a fixed target lie in a compact set of initial covectors and positive
times after allowing the boundary $\mathsf E=1$. The original
Hamiltonian flow, unlike elliptic phase representatives at $m=1$,
is analytic there. Every limiting source of energy one belongs to
$C_3\cup C_5$ and is a regular point of the full exponential map at
positive time \cite[Theorem 2.4]{SC}. The inverse function theorem
therefore excludes an accumulation of distinct sources of the same
target at that boundary.

The scalar function $N$ cannot be identically zero on a maximal
admissible component. Otherwise take angles approaching an endpoint
of that open component and the corresponding distinct period-zero
sources. The bounds just proved supply a convergent subsequence in
the original source coordinates. If its limit has energy greater
than one, Proposition~\ref{prop:inverse} applies at that source:
the target is off the seam, both angles are strict principal angles,
and the limiting angle would be in the same admissible component,
contrary to its being an endpoint. An energy-one limit contradicts
the preceding local injectivity. Hence $N\not\equiv0$ on every
component. Analyticity and the exact inverse now make each period-zero
source isolated. Infinitely many such sources would accumulate in the
compact set above, either at an isolated rotating source or at a
regular energy-one source; both alternatives are impossible.

Finally let the target range over a compact subset of $\Omega$.
The time and covector bounds are uniform, with a positive lower time
bound, so the closure of the joint period-zero source set is compact
in the original coordinates. A rotating limit still has period zero:
the only possible upper period boundary would have $t=4kK$ and
heading zero. At such a rotating limit, $N\not\equiv0$ supplies a
nonzero finite angle derivative. The product-neighborhood and Rolle
argument in Proposition~\ref{prop:positive-finite} bounds the nearby
roots uniformly. At an energy-one limit the inverse function theorem
gives at most one source per nearby target. A finite cover of the
compact closure yields a uniform bound for period zero. Adding the
positive-period bound proves the assertion for the full rotating fiber.
\end{proof}


\begin{remark}
The bounds are necessary conditions, not a count.  On the cells with
$n/R$ in a compact subinterval of $(0,1/L_*)$ the exact system, divided by
$R$, is a regular perturbation of $\eta\,L(m)\ee(a)=\ee(a_X)$, $\eta=n/R$,
$\theta=a-b$, whose four solutions (two initial and two terminal cosine
signs) continue by the implicit function theorem, and a compactness
argument excludes others: exactly four regular sources per bulk cell.
Theorem~\ref{thm:generic} below gives the complete eventual count, of
which this is the local mechanism; the compactness argument is that of
Theorem~\ref{thm:generic} localised to one cell.
\end{remark}

\section{The generic ray: global control of the scalar graph}

Assume $d_0=\sin^2b_0-\sin^2a_0\ne0$.
For large $R$, the sign $\varepsilon=-1$ is impossible:
\eqref{eq:decomposition} would give $y=A\cos a-B\sin a>-2$,
whereas $q_R$ has $y=-R\cos a_0<-2$.
For the other sign, $|B|<2$, $A>0$ uniquely determine
\begin{equation}\label{eq:Bdata}
 \begin{split}
 a&=a_0-\arcsin(B/R),\quad b=a-\theta,\quad A=\sqrt{R^2-B^2},\\
 d&=\sin^2b-\sin^2a,\qquad
 f=\frac{B^2+d}{2B},\quad g=\frac{d-B^2}{2B},\quad
 m_R(B)=\sin^2a+f^2.
 \end{split}
\end{equation}
All angles remain principal uniformly for $|B|<2$ and large $R$.
Moreover $d\to d_0$ uniformly, so the condition $m_R<1$ excludes
a fixed neighborhood of $B=0$. Proposition~\ref{prop:inverse} gives
$N_R(B)=(A-Q_0)/L(m_R(B))$ and the exact equation $N_R=n$.

\begin{lemma}[Domain and limiting profile]\label{lem:domain}
For sufficiently large $R$, the admissible $B$ domain consists of two
intervals. Their four endpoints, indexed by cosine signs
$\sigma_0,\sigma_1\in\{1,-1\}$, are
\begin{equation}\label{eq:endpoints}
 B_{\sigma_0\sigma_1}(R)=\frac{Ru}{\sqrt{(R-v)^2+u^2}},\qquad
 u=\sigma_0\cos a_0-\sigma_1\cos b_0,\quad
 v=\sigma_0\sin a_0-\sigma_1\sin b_0.
\end{equation}
On compact subsets of their limiting interiors,
\begin{equation}\label{eq:profile}
 \frac{N_R(B)}R\ \longrightarrow\ \frac1{L(m_0(B))}
 \quad\hbox{in }C^2,\qquad
 m_0(B)=\frac{B^2+2(\sin^2a_0+\sin^2b_0)+d_0^2/B^2}{4}.
\end{equation}
Each limiting component has one nondegenerate maximum of this profile,
at $B_*^\pm=\pm\sqrt{|d_0|}$, with value $1/L_*$.
\end{lemma}
\begin{proof}
At $m=1$, $f=\sigma_0\cos a$ and $g=\sigma_1\cos b$.
The equation $B=f-g$ becomes
$B(1-v/R)=u\sqrt{1-B^2/R^2}$. For $R>2$ the first factor is positive,
so its unique solution with the prescribed sign is \eqref{eq:endpoints}.
Substitution proves sufficiency. The four limits are
$\pm(\cos a_0+\cos b_0)$ and
$\pm|\cos a_0-\cos b_0|$, distinct by $d_0\ne0$.
They are all roots of $m_R=1$. The limiting sign chart, the excluded
neighborhood of zero, and $|B|<2$ therefore give exactly two components.
The modulus is positive because $m_R=0$ would force $a=b=0$, contrary
to the nonzero heading.

Direct differentiation gives
\[
 m_0'=\frac B2-\frac{d_0^2}{2B^3},\qquad
 m_0''=\frac12+\frac{3d_0^2}{2B^4}.
\]
Its two minima have value $m_*$ and curvature $2$.
The standard derivative identity \cite[Section 19.4]{DLMF} is
\begin{equation}\label{eq:Lm}
 L_m(m)=\frac{2E(m)}{1-m}>0.
\end{equation}
Hence the profile curvature at either maximum is $-2L_m(m_*)/L_*^2$.
On compact subsets the reconstructed phases have analytic local lifts,
even when a cosine vanishes. Nonzero heading keeps the residual time
away from both period boundaries. Thus $Q_0$ and its first two derivatives
remain bounded there, proving \eqref{eq:profile}.
\end{proof}

Compact convergence alone cannot exclude extra critical points close to
$m=1$. The next estimate supplies this missing endpoint control.
Define the signed short integral
\begin{equation}\label{eq:short}
 q(s,m)=\int_0^s\frac{u^2\,du}{\sqrt{1-u^2}\sqrt{m-u^2}},
 \qquad |s|<\sqrt m.
\end{equation}
It is odd in $s$. Near a domain endpoint, both cosine signs are fixed and
\begin{equation}\label{eq:tail}
 \begin{aligned}
 Q_0&=cL+\sigma_1q(\sin b,m)-\sigma_0q(\sin a,m),\\
 P_R&=A+\sigma_0q(\sin a,m)-\sigma_1q(\sin b,m),\qquad
 N_R=\frac{P_R}{L}-c,
 \end{aligned}
\end{equation}
where $c=1/2$ for opposite signs and
$c=\ind_{\{\sigma_0\theta>0\}}$ for equal signs.
To check the period coefficient, a phase of cosine sign $\sigma$ has a
lift $(1-\sigma)K+\sigma F(\arcsin(s/\sqrt m)\mid m)$.
Equal signs need a full wrap exactly when $\sigma_0\theta>0$;
opposite signs contribute one half-period. Subtracting the incomplete
second-kind integral proves \eqref{eq:tail} for either sign of $s$.

\begin{lemma}[Uniform endpoint monotonicity]\label{lem:tails}
There are neighborhoods of the four moving endpoints and $R_1$ such
that for every $R>R_1$ and every interior point of those neighborhoods,
$\operatorname{sgn}N_R'=-\operatorname{sgn}m_R'$.
The endpoint limits of $N_R$ are $-c$ and are approached from above.
\end{lemma}
\begin{proof}
The limiting endpoint roots are simple and the fixed angles are strict.
We can therefore choose the neighborhoods uniformly so that
$|\sin a|,|\sin b|\le s_{\max}<1$,
$m\ge m_c>s_{\max}^2$, and $|m_R'|\ge\gamma>0$.
Differentiation under \eqref{eq:short} bounds its needed derivatives
uniformly up to $m=1$. Formula \eqref{eq:Bdata} consequently gives
$P_R=R+O(1)$ and $|P_R'|\le C$ throughout these tails.
By the usual complete-integral asymptotics \cite[Section 19.12]{DLMF},
$L/L_m$ extends continuously to zero at $m=1$ and is bounded, say by $H$.
For sufficiently large $R$, with $P_R\ge R-C_0$,
\[
 |P_R'|\frac L{L_m}\le CH<(R-C_0)\gamma\le P_R|m_R'|.
\]
The identity $N_R'=(P_R'L-P_RL_m m_R')/L^2$ proves the sign.
Finally $P_R>0$ is finite as $m\to1$ at fixed $R$, while $L\to\infty$.
Thus $N_R\to-c$ from above.
\end{proof}

\begin{theorem}[Exact eventual generic count]\label{thm:generic}
When $d_0\ne0$, for sufficiently large $R$ each admissible component
has exactly one critical point, a nondegenerate maximum. If the heights
are $M_-(R)$ and $M_+(R)$, then
\begin{equation}\label{eq:floor}
 \nu_2(q_R)=2\lfloor M_-\rfloor+2\lfloor M_+\rfloor+3
 -\ind_{\{M_-\in\mathbb Z\}}-\ind_{\{M_+\in\mathbb Z\}}.
\end{equation}
Moreover $M_\pm=R/L_*+O(1)$. Every source is regular except an
integer-maximum source, which is a fold.
\end{theorem}
\begin{proof}
Near either profile maximum in Lemma~\ref{lem:domain}, strict negative
second derivative and the implicit function theorem give exactly one
nearby maximum. On the remaining compact interior, $C^1$ convergence
excludes critical points. Lemma~\ref{lem:tails} excludes them on all four
tails, exhausting the moving domain.
Analytic dependence on $1/R$ near either maximum gives
$M_\pm=R/L_*+O(1)$, so both heights are positive for large $R$.

The two opposite-cosine endpoints are the outer endpoints; each has
level limit $-1/2$. The equal-cosine endpoints are the inner ones,
with respective limits $0$ and $-1$. Thus one component has limits
$0,-1/2$, and the other $-1,-1/2$. On the first, level zero has one
root; on the second it has two. Every positive integer below a maximum
has two roots in that component. At an integer maximum those two roots
coalesce to one. Strict monotonicity on each arm proves \eqref{eq:floor}.
There is no $B=0$ source: that condition in \eqref{eq:decomposition}
would force $b_0=-a_0$. The other sign was already excluded, so the count
is complete. Proposition~\ref{prop:fold} applies because
$dz/dB=-(1+z^2)/\sqrt{R^2-B^2}\ne0$.
Equation \eqref{eq:main} follows in this case.
\end{proof}

Figure~\ref{fig:level} shows the level function $N$ of
\eqref{eq:residual}, in the chart $z=\tan a$ of
Proposition~\ref{prop:inverse}, for one target of this kind.

\begin{figure}[htbp]
\centering
\includegraphics[width=\textwidth]{../figures/counts_level_function.pdf}
\caption{The level function $N(z)$ of \eqref{eq:residual} for the target
$(39/5,-52/5,3/10)$, that is $R=13$ and $\theta=3/10$, on its two
admissible components, against $z=\tan a$. The horizontal lines are the
integer levels and the dots the $37$ integer crossings ($21$ and $16$),
one rotating source each, in agreement with \eqref{eq:floor}. When a
maximum crosses an integer level as the target varies, the corresponding
sources merge in a fold (Proposition~\ref{prop:fold}); a noninteger
critical level alone does not represent a source of this target.}
\label{fig:level}
\end{figure}

\section{The half-heading ray and its separate seam}\label{sec:seam}

The only remaining case is $b_0=-a_0$. First take
$\alpha=a_0\in(0,\pi/2)$ and abbreviate $s=\sin\alpha$,
$c_\alpha=\cos\alpha$, $L_*=L(s^2)$.
For $B\ne0$, the exact formulas \eqref{eq:Bdata} reduce to
\begin{align}
 a&=\alpha-\arcsin(B/R),& b&=-\alpha-\arcsin(B/R),\notag\\
 f&=\frac B2+\frac{\sin(2\alpha)\sqrt{R^2-B^2}}{R^2},&
 g&=-\frac B2+\frac{\sin(2\alpha)\sqrt{R^2-B^2}}{R^2},\notag\\
 m_R(B)&=m_c+D_RB^2,&
 m_c&=s^2+\frac{4s^2c_\alpha^2}{R^2},\label{eq:raymod}\\
 D_R&=\frac{(R^2-4s^2)(R^2+4c_\alpha^2)}{4R^4}.&&\notag
\end{align}
For $R>2s$, this formula has admissible interval
\begin{equation}\label{eq:raydom}
 |B|<B_{\max}=\frac{2Rc_\alpha}{\sqrt{R^2+4c_\alpha^2}}.
\end{equation}
Indeed the largest possible absolute angle is less than
$\alpha+\arctan(2c_\alpha/R)<\pi/2$, so no further angle boundary occurs.

\begin{proposition}[Nonzero transverse sources]\label{prop:offseam}
For sufficiently large $R$, the extended scalar $N_R$ on
\eqref{eq:raydom} is even and has exactly one critical point, a
nondegenerate maximum at zero. If $M_0=N_R(0)$, its nonzero-$B$ sources
number $2\max(0,\lceil M_0\rceil)$, and
\[
 M_0=\frac{R}{L(m_c)}-1+\frac{2q_c}{L(m_c)}=\frac R{L_*}-\frac12+O(R^{-1}),
\]
where $m_c$ is the modulus selected at $B=0$, $h_c=2sc_\alpha/R$ and
$q_c=F(\arctan(s/h_c)\mid m_c)-E(\arctan(s/h_c)\mid m_c)$; the counts use
the exact expression, whose $O(R^{-1})$ term is not small at moderate
$R$.  The extension at zero selects only modulus $m_c$; it does not
represent the full $B=0$ fiber. The nonzero-$B$ count formula extends to every
$R>2s$ by Theorem~\ref{thm:ray-all-radii} below; the unimodality
statement here remains restricted to sufficiently large $R$.
\end{proposition}
\begin{proof}
Replacing $B$ by $-B$ sends $(a,b,f,g)$ to $(-b,-a,g,f)$.
The phases become $(-v,-\psi)$ modulo $4K$, preserving the positive
residual time and the integral of the even function $\sn^2$.
Hence $Q_0$, $m_R$, and $N_R$ are even. Their phase parametrization
extends analytically at zero: its sine entries are nonzero and the
residual time stays strictly between zero and $4K$. Let
$h_c=2sc_\alpha/R$. At fixed finite $R$ the value at zero is the long
arc $4K-2u_c$, where $u_c=F(\arctan(s/h_c)\mid m_c)<K$;
this approaches $2K$ as $R\to\infty$.
The normalized profile converges in $C^2$ on compact interiors to
$1/L(s^2+B^2/4)$, whose unique maximum has curvature
$-L_m(s^2)/(2L_*^2)$. The compact argument of
Theorem~\ref{thm:generic} and Lemma~\ref{lem:tails} apply; the two
endpoint signs are opposite and their limits are $-1/2$.
Evenness places the sole maximum exactly at zero.

Let
$q_c=F(\arctan(s/h_c)\mid m_c)-E(\arctan(s/h_c)\mid m_c)$.
At zero the residual is the long arc, so
$M_0=R/L(m_c)-1+2q_c/L(m_c)$.
Here $m_c=s^2+O(R^{-2})$ and $q_c=L_*/4+O(R^{-1})$,
which proves the height estimate. Every nonnegative integer strictly
below the height contributes a pair. An integer maximum is excluded
because it lies at $B=0$, to be counted separately below.
\end{proof}

\begin{proposition}[Complete zero-transverse fiber]\label{prop:seam}
Set $Q=\atanh(s)-s$. For each integer $n\ge1$ there is a unique
$m_n\in(s^2,1)$ minimizing
$R_n^{\rm short}(m)=nL(m)+2q(s,m)$. Write its value as $S_n$.
Then the number of $B=0$ sources is exactly
\begin{equation}\label{eq:seamcount}
 \ind_{\{R>2Q\}}+2\#\{n\ge1:S_n<R\}+\#\{n\ge1:S_n=R\},
\end{equation}
and it equals $2R/L_*+O(1)$. In particular it is finite for every
fixed positive $R$. The thresholds satisfy
\begin{equation}\label{eq:Sbounds}
 nL_*+2Q<S_n<(n+1/2)L_*,\qquad S_{n+1}>S_n+L_*.
\end{equation}
\end{proposition}
\begin{proof}
At $B=0$, the equations force $a=\alpha$, $b=-\alpha$, and
$f=g$. For $m\in[s^2,1)$, put
$u=F(\arcsin(s/\sqrt m)\mid m)$.
The initial phases are $2K-u$ and $u$, with terminal phase $-\psi$.
Their radii in actual period cell $n\ge0$ are exactly
\begin{equation}\label{eq:seambranches}
 R_n^{\rm short}(m)=nL(m)+2q(s,m),\qquad
 R_n^{\rm long}(m)=(n+1)L(m)-2q(s,m).
\end{equation}
At $m=s^2$ the two phases coincide and both radii are
$H_n=(n+1/2)L_*$. This source must be counted once.

For $m>s^2$, differentiation under the short integral gives
\[
 q_m=-\frac12\int_0^s
 \frac{u^2\,du}{\sqrt{1-u^2}(m-u^2)^{3/2}}<0.
\]
Consequently every long branch is strictly increasing from $H_n$ to
infinity. The short branch for $n=0$ decreases from $H_0$ to $2Q$,
the latter limit being unattained. After deduplication at $H_0$ their
union contributes exactly $\ind_{\{R>2Q\}}$.

For $n\ge1$ define
\[
 \mathcal H(m)=-\frac{2q_m}{L_m}
 =\frac{1-m}{2E(m)}\int_0^s
 \frac{u^2\,du}{\sqrt{1-u^2}(m-u^2)^{3/2}}.
\]
It decreases strictly from infinity to zero. To see the derivative
sign, its integral is positive and strictly decreasing, and its positive
prefactor is strictly decreasing because
\[
 E+(1-m)E_m=\frac{(m+1)E-(1-m)K}{2m}>0.
\]
The last inequality follows directly from
$E-(1-m)K=m\int_0^{\pi/2}\cos^2\phi/
\sqrt{1-m\sin^2\phi}\,d\phi>0$.
The lower-endpoint divergence of the integral and its finite limit at
$m=1$ prove the asserted limits of $\mathcal H$.
Thus $(R_n^{\rm short})'=L_m(n-\mathcal H)$ has exactly one zero,
with second derivative $-L_m\mathcal H'>0$.

The short branch decreases from $H_n$ to $S_n$, then increases to
infinity. Combining it with the long branch gives zero, one, or two
distinct sources according as $R<S_n$, $R=S_n$, or $R>S_n$.
At $R=H_n$, its lower-endpoint source is shared by the two branches;
the other short-branch root remains distinct, giving two in total.
This proves \eqref{eq:seamcount}.
Since $L(m)>L_*$ and $q(s,m)>Q$, the lower bound in
\eqref{eq:Sbounds} follows; the initial decrease gives the upper bound.
Evaluating $R_n^{\rm short}+L$ at the next minimizer gives
$S_{n+1}>S_n+L_*$. These bounds prove finiteness and
$\#\{n\ge1:S_n<R\}=R/L_*+O(1)$.
\end{proof}

\begin{proposition}[Full-map singularities away from the branch join]
\label{prop:seam-singularities}
Fix $0<\alpha<\pi/2$ and $m\in(s^2,1)$ in
\eqref{eq:seambranches}, and put $h=\sqrt{m-s^2}$ and
$c_\alpha=\cos\alpha$.
On a short branch with $n\ge1$, the unique minimum $S_n$ is a
corank-one fold of the full exponential map, and all other points
are regular. The short branch with $n=0$ is regular throughout.
Each long branch, including $n=0$, has exactly one singular point,
characterized by
\begin{equation}\label{eq:seam-long-singular}
 R_n^{\rm long}(m)=\frac{2sc_\alpha}{\sqrt{m-s^2}}.
\end{equation}
It has corank one and is a cusp, as proved in
Corollary~\ref{cor:seam-cusps}. The common branch endpoint $m=s^2$
is regular by Proposition~\ref{prop:seam-local-counts}.
\end{proposition}
\begin{proof}
Use local source coordinates $(a,\theta,m)$ with $b=a-\theta$.
They are valid when the signed cosines $f,g$ are nonzero: in the
original phase and time variables, $a_\psi=f$ and $b_t=g/k$.
Write $\sigma=-1$ on the short branch and $\sigma=1$ on the long
branch. Near the source in question,
\[
 \begin{gathered}
 f=\sigma\sqrt{m-\sin^2a},\quad
 g=\sigma\sqrt{m-\sin^2b},\quad B=f-g,\\
 A=\left(n+\frac{1+\sigma}{2}\right)L
       -\sigma\{q(\sin a,m)-q(\sin b,m)\}.
 \end{gathered}
\]
Here the defining integral for $q$ is extended to negative upper
limits, and all angles remain principal. At fixed $\theta$, the
endpoint $X=A\ee(a)+B\ee^\perp(a)$ has planar Jacobian
\[
 J=(A_a-B)B_m-(A+B_a)A_m.
\]
At the seam $a=\alpha$, $b=-\alpha$ one has
$A_a=B=B_m=0$ and $B_a=-2\sigma sc_\alpha/h$. Thus
\begin{equation}\label{eq:seam-full-jacobian}
 J=-H A_m,\qquad H=A-\frac{2\sigma sc_\alpha}{h},
 \qquad X_a=H\ee^\perp(\alpha),\quad X_m=A_m\ee(\alpha).
\end{equation}
Since the third output is $\theta$, these formulas also determine
the rank of the full map.

For a short branch $H>0$. Proposition~\ref{prop:seam} gives
$A_m=0$ exactly at its nondegenerate minimum when $n\ge1$;
there the kernel is $\partial_m$ and
$\partial_mJ=-H A_{mm}\ne0$. This is the corank-one fold
criterion, equivalently the parameter-dependent Morse lemma after
solving for the transverse output. All remaining short points have
$J\ne0$, including the entire $n=0$ branch.

For a long branch $A_m>0$ and
$H_m=A_m+sc_\alpha/h^3>0$. Moreover $H$ tends to $-\infty$
as $m\downarrow s^2$ and to $+\infty$ as $m\uparrow1$.
This proves existence and uniqueness in
\eqref{eq:seam-long-singular}. At that point the full rank is two,
with kernel $\partial_a$. To evaluate the fold criterion, put
$u=a-\theta/2$ at fixed $\theta=2\alpha$. Reflection $u\mapsto-u$
makes $A$ even and $B$ odd. The components of $X$ in the fixed
basis at $\alpha$ are consequently
$U=A\cos u-B\sin u$ even and $V=A\sin u+B\cos u$ odd.
Their Jacobian $U_uV_m-V_uU_m$ is even, so $\partial_aJ=0$
at $u=0$. The fold criterion therefore fails, although
$J_m=-H_mA_m\ne0$ and the critical set is smooth there.
\end{proof}

\begin{corollary}[Long-branch cusps]\label{cor:seam-cusps}
Every long-branch singularity in
Proposition~\ref{prop:seam-singularities} is an ordinary corank-one
cusp of the full exponential map, including period $n=0$.
\end{corollary}
\begin{proof}
Extend $U,V$ to target coordinates using the basis at $\theta/2$,
and write $u=a-\theta/2$. Since $U_m=A_m>0$ at the singularity, solve
locally for $m=m(u,U,\theta)$ and put
$W(u,U,\theta)=V(u,m(u,U,\theta),\theta)$.
The full map is then $(u,U,\theta)\mapsto(W,U,\theta)$.
At the singularity $W_u=0$, and reflection gives $W_{uu}=0$, while
\[
 W_{uU}=\frac{H_m}{A_m}>0,\qquad
 \mathcal C:=W_{uuu}=V_{uuu}-3\frac{V_{um}U_{uu}}{A_m}.
\]
All derivatives here are evaluated at the singular source.
Differentiating the local endpoint formulas and using
$A=2sc_\alpha/h$ gives
\[
 U_{uu}=-\frac{2sc_\alpha m}{h^3},\qquad
 V_{uuu}=-\frac{6sc_\alpha^3m}{h^5},\qquad
 V_{um}=A_m+\frac{sc_\alpha}{h^3}.
\]
Consequently
\begin{equation}\label{eq:seam-cusp-cubic}
 \mathcal C=\frac{6sc_\alpha m}{h^6A_m}
       \{sc_\alpha-(1-m)hA_m\}.
\end{equation}
To determine its sign, define
\[
 I=\int_0^\alpha\frac{\sin^2v}{(m-\sin^2v)^{3/2}}\,dv,
 \qquad
 \mathcal E_\alpha=\int_0^\alpha
       \frac{\cos^2v}{\sqrt{m-\sin^2v}}\,dv.
\]
The long-branch formula gives
$A_m=2(n+1)E(m)/(1-m)+I$. Direct differentiation shows
\[
 \frac{d}{dv}\frac{\sin v\cos v}{\sqrt{m-\sin^2v}}
 =\frac{(1-m)\sin^2v}{(m-\sin^2v)^{3/2}}
       +\frac{\cos^2v}{\sqrt{m-\sin^2v}}.
\]
Integration therefore yields $sc_\alpha/h=(1-m)I+\mathcal E_\alpha$.
The substitution $\sin v=\sqrt m\sin\phi$ gives
$\mathcal E_\alpha=E(\arcsin(s/\sqrt m)\mid m)$, so
$0<\mathcal E_\alpha<E(m)$. Thus
\[
 \mathcal C=\frac{6sc_\alpha m}{h^5A_m}
       \{\mathcal E_\alpha-2(n+1)E(m)\}<0.
\]
For $\lambda=W_u$ and the kernel extension $\partial_u$, one has
$\lambda=\lambda_u=0$, $\lambda_{uu}=\mathcal C\ne0$.
Moreover $(\lambda,\lambda_u)$ has rank two, since its derivatives
in the $U,u$ directions have determinant $W_{uU}\mathcal C\ne0$.
The cusp criterion \cite[Theorem A.1]{SUY} applies to this
three-dimensional map.
\end{proof}
\begin{proposition}[Regular joins and local seam counts]
\label{prop:seam-local-counts}
At the common endpoint $m=s^2$ of each pair of branches
\eqref{eq:seambranches}, the full exponential map is regular.
In the original coordinates $(\gamma_0,c_0,t)$ for the positive
rotation sign, its Jacobian determinant there is $m^2>0$.

Fix $\alpha$ and vary the target radius $R$. In a sufficiently small
source neighborhood of a short-branch minimum $S_n$, $n\ge1$,
the numbers of sources are $0,1,2$ according as $R<S_n$, $R=S_n$,
or $R>S_n$. At the long-branch cusp of period $n\ge0$, write its
radius as $R_n^{\rm cusp}$. The corresponding local counts are $1,1,3$
according as $R<R_n^{\rm cusp}$, $R=R_n^{\rm cusp}$, or $R>R_n^{\rm cusp}$. The extra two sources
on the last side have $B\ne0$. Near a common branch endpoint,
the local count remains one. These statements concern radii
sufficiently close to the respective threshold.
\end{proposition}
\begin{proof}
Use the initial and terminal pendulum amplitudes
$\phi_0,\phi_1$ and $m$ as source coordinates. At the join,
\[
 \phi_0=\pi/2,\quad \phi_1=3\pi/2+2\pi n,\quad
 k=\sqrt m=s,\quad k'=\sqrt{1-m}=c_\alpha.
\]
In a neighborhood of this source the defining-integral endpoint is
\[
 \begin{gathered}
 a=\arcsin(k\sin\phi_0),\quad b=\arcsin(k\sin\phi_1),
 \quad B=k(\cos\phi_0-\cos\phi_1),\\
 A=\int_{\phi_0}^{\phi_1}
       \frac{m\sin^2v}{\sqrt{1-m\sin^2v}}\,dv,
 \qquad X=A\ee(a)+B\ee^\perp(a),\quad\theta=a-b.
 \end{gathered}
\]
These expressions are analytic, including at the join. There,
$a_{\phi_0}=b_{\phi_1}=0$, and differentiation gives
\[
 \begin{gathered}
 X_{\phi_0}=-\frac m{k'}\ee(\alpha)-k\ee^\perp(\alpha),\quad
 X_{\phi_1}=\frac m{k'}\ee(\alpha)-k\ee^\perp(\alpha),\\
 \theta_{\phi_0}=\theta_{\phi_1}=0,\qquad \theta_m=\frac1{kk'}.
 \end{gathered}
\]
Hence $\det\partial(X,\theta)/\partial(\phi_0,\phi_1,m)=2m/(1-m)$.
The conversion to the original source coordinates is
\[
 \gamma_0=2\phi_0,\quad
 c_0=\frac2k\sqrt{1-m\sin^2\phi_0},\quad
 t=k\int_{\phi_0}^{\phi_1}\frac{dv}{\sqrt{1-m\sin^2v}},
\]
with $\gamma_0$ taken locally modulo $4\pi$. Its determinant at
the join is $2/[m(1-m)]$: indeed $(c_0)_m=-1/(k^3k')$
and $t_{\phi_1}=k/k'$. The ratio is $m^2>0$.
The inverse function theorem proves the local count one there.

At a short minimum, the transverse derivative $V_u=H>0$ in
the coordinates of Proposition~\ref{prop:seam-singularities}.
Reflection and the implicit function theorem force $u=0$ for
nearby targets on the same ray. The nondegenerate radial minimum
then gives the local $0,1,2$ rule.

At a long cusp, use $W(u,R,2\alpha)$ from
Corollary~\ref{cor:seam-cusps}. Analyticity and oddness in $u$
give $W=uG(u^2,R)$ for an analytic $G$. At $(0,R_n^{\rm cusp})$,
\[
 G=0,\qquad G_z=\mathcal C/6<0,\qquad
 G_R=H_m/A_m>0.
\]
The unique local solution of $G(z,R)=0$ therefore satisfies
$z(R_n^{\rm cusp})=0$ and $z'(R_n^{\rm cusp})=-6H_m/(A_m\mathcal C)>0$.
For $R<R_n^{\rm cusp}$ this root is negative and contributes no real $u$;
for $R>R_n^{\rm cusp}$ it gives the two roots $u=\pm\sqrt{z(R)}$,
in addition to $u=0$. At equality all three coincide as one source.
Along the fixed-radius reduction, $B_u=-2sc_\alpha/h\ne0$
at the cusp, so the two new sources have $B\ne0$.
The actual period is locally constant at these sources.
\end{proof}

\begin{lemma}[Regularity of nonzero transverse sources on the ray]
\label{lem:ray-maxwell-regular}
Every rotating source of $(R\ee(\alpha),2\alpha)$ with $B\ne0$
is regular, for every $R>0$ and $0<\alpha<\pi/2$.
\end{lemma}
\begin{proof}
Put $p=t/(2k)$, $\tau=\psi+p$, so $\tau$ is the midpoint
phase, and set
\[
 \begin{gathered}
 \epsilon_p=\int_0^p\dn^2v\,dv,\qquad q_p=p-\epsilon_p>0,\\
 f_1(p)=-q_p\cn p-\dn p\sn p,\qquad
 \Delta=1-m\sn^2\tau\sn^2p>0.
 \end{gathered}
\]
For the positive rotation sign, the Jacobi addition formulas applied
to \eqref{eq:decomposition} give the half-heading components
\begin{equation}\label{eq:ray-maxwell-components}
 \begin{gathered}
 U=X\cdot\ee(\theta/2)=\frac{2q_p\dn\tau}{\sqrt\Delta}>0,
 \qquad
 V=X\cdot\ep(\theta/2)=\frac{-2k\sn\tau\,f_1(p)}{\sqrt\Delta},\\
 B=\frac{2k\sn\tau\sn p\dn\tau\dn p}{\Delta},
 \qquad
 \sin(\theta/2)=\frac{-k\cn\tau\sn p}{\sqrt\Delta}.
 \end{gathered}
\end{equation}
For example, if $\delta=(a+b)/2$, the addition formulas give
$\cos\delta=\dn\tau/\sqrt\Delta$ and
$\sin\delta=k\sn\tau\cn p/\sqrt\Delta$; inserting these into
$U=A\cos\delta-B\sin\delta$ and
$V=A\sin\delta+B\cos\delta$ proves the first two identities.
The negative rotation sign negates $U$ and cannot reach the stated ray.

On the ray $V=0$, while $B\ne0$ forces $\sn\tau\ne0$.
Thus $f_1(p)=0$. In particular $\cn p\ne0$, and
$q_p=-\dn p\sn p/\cn p$.
Substitution into Sachkov's reduced Jacobian
\cite[(2.10)--(2.14)]{SC} gives
\begin{equation}\label{eq:ray-maxwell-jacobian}
 J_2=(1-m)q_p\epsilon_p\sn^2\tau>0.
\end{equation}
Indeed the coefficient of $\sn^2\tau$ simplifies to
$(1-m)q_p\epsilon_p$, and the other coefficient contains $f_1$.
The full Jacobian differs by a nonvanishing factor, proving regularity.
\end{proof}

\begin{theorem}[Exact rotating count on the ray at every radius]
\label{thm:ray-all-radii}
Fix $0<\alpha<\pi/2$. For each $n\ge0$, let
$\mu_n\in(s^2,1)$ be the unique long-branch cusp modulus, and put
\[
 R_n^{\rm cusp}=\frac{2sc_\alpha}{\sqrt{\mu_n-s^2}}.
\]
In actual period $n$, there are exactly two nonzero-$B$ sources of
$(R\ee(\alpha),2\alpha)$ when $R>R_n^{\rm cusp}$, and none
when $R\le R_n^{\rm cusp}$. Both are regular. Consequently the
full rotating-source count is
\begin{equation}\label{eq:ray-global-count}
 \begin{aligned}
 N_{\rm rot}(R,\alpha)
 ={}&\ind_{\{R>2Q\}}
   +2\#\{n\ge1:S_n<R\}+\#\{n\ge1:S_n=R\}\\
   &+2\#\{n\ge0:R_n^{\rm cusp}<R\},\qquad R>0.
 \end{aligned}
\end{equation}
The cusp thresholds increase strictly and satisfy
$R_n^{\rm cusp}>(n+1/2)L_*$, so all sums are finite.
For $R>2s$, the nonzero-$B$ count also equals
$2\max(0,\lceil M_0\rceil)$ with $M_0=N_R(0)$ as in
Proposition~\ref{prop:offseam}. This count identity does not assert
unimodality of the scalar at every radius.
\end{theorem}
\begin{proof}
By \cite[Lemma 5.3]{MS}, for each $m\in(0,1)$ the positive roots
of $f_1$ consist of one
simple root $P_n(m)$ in each interval
$((2n+1)K,2(n+1)K)$, $n\ge0$. To check this directly,
\[
 \left(\frac{f_1(p)}{\cn p}\right)'
       =-\frac{\dn^2p}{\cn^2p}<0.
\]
On the stated interval the quotient decreases from $+\infty$ to
$-q_{2(n+1)K}<0$. On the complementary intervals
$(2nK,(2n+1)K)$ it is $-q_p-\dn p\tan(\operatorname{am}p)<0$.
The endpoints are not positive roots. Simplicity and analyticity give
analytic $P_n(m)$; its actual period is
$\lfloor P_n/(2K)\rfloor=n$.

Fix $n$, abbreviate $p=P_n(m)$ and $w=m\sn^2p$, and retain
$q_p=p-\epsilon_p$. Equation~\eqref{eq:ray-maxwell-components}
for heading $2\alpha$ is equivalent to
\begin{equation}\label{eq:ray-maxwell-phase}
 \sn^2\tau=\frac{w-s^2}{wc_\alpha^2},\qquad
 \operatorname{sgn}(\cn\tau)=-\operatorname{sgn}(\sn p).
\end{equation}
It gives exactly two phases modulo $4K$ when $w>s^2$, one
when $w=s^2$, and none when $w<s^2$. The two phases have
opposite nonzero $\sn\tau$ and hence opposite nonzero $B$.
Their common radius is
\begin{equation}\label{eq:ray-maxwell-radius}
 R_n(m)^2=4(q_p^2+s^2-w)>4s^2.
\end{equation}
This follows from $U=2q_p\dn\tau/\sqrt\Delta$ and
$q_p^2=(1-m\sn^2p)\sn^2p/\cn^2p$.

The equality $w=s^2$ is precisely the long-branch cusp condition:
write $p=2(n+1)K-v$, where
$v=F(\arcsin(s/\sqrt m)\mid m)$, and substitute into $f_1=0$.
It becomes $R_n^{\rm long}(m)=2sc_\alpha/\sqrt{m-s^2}$.
Its unique solution is $\mu_n$ by
Proposition~\ref{prop:seam-singularities}.
For $m<s^2$ one has $w<s^2$. As $m\uparrow1$,
\[
 q_p>(2n+1)K-2(n+1)E\longrightarrow\infty.
\]
The identity $q_p^2=(1-mx)x/(1-x)$, $x=\sn^2p$, implies
$x\to1$, and therefore $w\to1$. Continuity and the unique equality
give $w>s^2$ exactly on $(\mu_n,1)$.

On this interval both phase branches are analytic and regular.
Moreover $\theta_\tau=B\ne0$ and, on $f_1=0$,
$V_p=-2k\sn\tau\,f_1'(p)/\sqrt\Delta\ne0$, while $V_\tau=0$.
Thus the constraints $V=0$, $\theta=2\alpha$ have independent
derivatives in $(p,\tau)$ and leave $m$ as a local parameter.
Lemma~\ref{lem:ray-maxwell-regular} then implies $R_n'(m)\ne0$.
The radius tends to $R_n^{\rm cusp}$ as $m\downarrow\mu_n$
and to infinity as $m\uparrow1$, by
\eqref{eq:ray-maxwell-radius}. Hence it increases strictly and
maps $(\mu_n,1)$ onto $(R_n^{\rm cusp},\infty)$.
This proves the asserted exhaustive count of nonzero-$B$ sources.

The function $R_n^{\rm long}(m)-2sc_\alpha/\sqrt{m-s^2}$
increases strictly in $m$, and increasing $n$ adds $L(m)>0$.
Consequently $\mu_{n+1}<\mu_n$ and the cusp radii increase strictly.
Also $R_n^{\rm cusp}=R_n^{\rm long}(\mu_n)>(n+1/2)L_*$.
Adding the complete $B=0$ count \eqref{eq:seamcount} proves
\eqref{eq:ray-global-count}; the strict inequality at a cusp avoids
counting its single $B=0$ source again.

Finally, for $R>2s$ put $m_c=s^2+(2sc_\alpha/R)^2$.
The long-branch equation at this selected modulus reads
\[
 R_n^{\rm long}(m_c)-R=L(m_c)(n-M_0).
\]
Its sign and strict monotonicity in $m_c$ show that
$R>R_n^{\rm cusp}$ is equivalent to $n<M_0$.
Counting the nonnegative integers with that strict inequality gives
$2\max(0,\lceil M_0\rceil)$.
\end{proof}

\begin{corollary}[Minimizing part of the ray count]\label{cor:ray-cut-status}
Write $C=R_0^{\rm cusp}$ and let $M_{\rm rot}$ count the minimizing
rotating sources of $(R\ee(\alpha),2\alpha)$. Then
\[
 M_{\rm rot}=
 \begin{cases}
 0,&0<R\le2Q,\\
 1,&2Q<R\le C,\\
 2,&R>C.
 \end{cases}
\]
For $2Q<R<C$ the minimizing source is strictly pre-cut; at $R=C$
it is the period-zero cusp at its cut time. For $R>C$ the two
minimizing sources form the period-zero regular Maxwell pair, both
at their cut time. Every other rotating source is strictly post-cut,
so their number is $N_{\rm rot}-M_{\rm rot}$.
In the notation of \cite[(2.13), (2.20)]{SS},
$2Q=R_1^2(2\alpha)$ and $C=R_1^1(2\alpha)$.
\end{corollary}
\begin{proof}
The established rotating cut time is $2kP_0(m)$, with
$K<P_0<2K$, and a geodesic minimizes through that time
\cite[(2.2), Theorem 3.3]{SS}. Every source of actual period
$n\ge1$ is therefore post-cut. The nonzero-$B$ sources have
$p=P_n(m)$, so precisely their period-zero pair is at the cut time.

For $B=0$, period zero has one source exactly when $R>2Q$.
On its short branch $p=u\le K<P_0$, where
$u=F(\arcsin(s/\sqrt m)\mid m)$. On its long branch
$p=2K-u$ and, with $h=\sqrt{m-s^2}$,
\[
 f_1(2K-u)=\frac{h}{2k}
 \left(R_0^{\rm long}(m)-\frac{2sc_\alpha}{h}\right).
\]
The parenthesis increases through zero at $m=\mu_0$.
On $(K,2K)$, $f_1$ is negative before $P_0$ and positive after it.
Since $R_0^{\rm long}$ increases strictly, this source is pre-cut,
at cut, or post-cut according as $R<C$, $R=C$, or $R>C$.
The branch join has $p=K$ and is strictly pre-cut as well.

At $m=\mu_0$, put $v=\operatorname{am}P_0\in(\pi/2,\pi)$.
Then $\sin\alpha=k\sin v$ and
$C=2(P_0-\epsilon_{P_0})$. These are exactly the parametric
equations of $R_1^1(2\alpha)$ in \cite[(2.18)--(2.20)]{SS};
the formula for $R_1^2$ gives $2Q$ directly. In those rectifying
coordinates the present ray has $R_1=-R$, $R_2=0$.
This recovers the known minimizing thresholds without identifying
the additional post-cut sources with minimizers.
\end{proof}

\begin{proposition}[Cusp scale near zero heading]\label{prop:ray-cusp-scale}
Fix an actual period $n\ge0$, put $j=n+1$ and $\theta=2\alpha$,
and define $z=(\theta/(j\pi))^{2/3}$. As $\theta\downarrow0$,
\begin{equation}\label{eq:ray-cusp-small-heading}
 \mu_n=z-\tfrac14z^2+O_j(z^3),\qquad
 R_n^{\rm cusp}=j\pi\left(z+\tfrac18z^2+O_j(z^3)\right).
\end{equation}
In particular $R_n^{\rm cusp}\sim(j\pi)^{1/3}\theta^{2/3}$.
These are fixed-period estimates; no uniform remainder as $n\to\infty$
is asserted.
\end{proposition}
\begin{proof}
Write the Maxwell amplitude as
$\operatorname{am}P_n=j\pi-\delta$ and define
\[
 D=\sqrt{1-m\sin^2\delta},\qquad
 \mathcal Q=F(j\pi-\delta\mid m)-E(j\pi-\delta\mid m).
\]
The root condition is $\mathcal Q=D\tan\delta$.
At $(m,\delta)=(0,0)$ its derivative in $\delta$ is $-1$,
so the implicit function theorem gives an analytic $\delta(m)$
with $\delta(0)=0$. From the defining integral,
\[
 \mathcal Q=m\int_0^{j\pi-\delta}
       \frac{\sin^2 v}{\sqrt{1-m\sin^2 v}}\,dv
 =\frac{j\pi}{2}m+\frac{3j\pi}{16}m^2+O_j(m^3).
\]
Here $\delta=O_j(m)$, and removing the terminal interval contributes
$O_j(m\delta^3)$. Since $D=1+O_j(m^3)$, the root equation gives
$\delta=j\pi m/2+3j\pi m^2/16+O_j(m^3)$.
Its positive leading coefficient places this root in the Maxwell
interval used in Theorem~\ref{thm:ray-all-radii}. On the cusp,
\[
 \begin{aligned}
 \theta&=2\arcsin(\sqrt m\sin\delta)
 =j\pi m^{3/2}\left(1+\tfrac38m+O_j(m^2)\right),\\
 R_n^{\rm cusp}&=2\mathcal Q
 =j\pi m\left(1+\tfrac38m+O_j(m^2)\right).
 \end{aligned}
\]
Consequently $z=m+m^2/4+O_j(m^3)$; inversion and substitution
give \eqref{eq:ray-cusp-small-heading}.
\end{proof}
\begin{remark}[Normalization of the published first threshold]
For $n=0$, this yields
$R_1^1(\theta)=\pi^{1/3}\theta^{2/3}+O(\theta^{4/3})$
in the notation of \cite{SS}. Lemma 2.5(6) of that paper prints
half this leading coefficient. Its exact parametrization (2.19),
$R_1=2(F-E)$, instead gives $R_1=\pi k^2+O(k^4)$,
whereas the displayed expansion on p.~304 gives $\pi k^2/2$.
We use the exact parametrization. The correction concerns this
asymptotic coefficient, not the curve or the optimal synthesis.
\end{remark}

\begin{theorem}[Count blow-up near zero heading]\label{thm:ray-small-heading-count}
For $R>\max(2s,2Q)$, put $m_c=s^2+4s^2c_\alpha^2/R^2$.
Then the full rotating count on the half-heading ray satisfies
\begin{equation}\label{eq:ray-bounded-count-error}
 -4<N_{\rm rot}-\frac{2R}{L_*}-\frac{2R}{L(m_c)}<2.
\end{equation}
Consequently, for each fixed $R>0$, as $\theta=2\alpha\downarrow0$,
\begin{equation}\label{eq:ray-small-heading-count}
 N_{\rm rot}(R,\theta/2)
 =\frac{16R(R^2+2)}{\pi(R^2+4)\theta^2}+O_R(1).
\end{equation}
The remainder is uniform for $R$ in any compact subinterval of
$(0,\infty)$. Exactly two of these sources minimize for sufficiently
small $\theta$; the strictly post-cut count has the same asymptotic law.
\end{theorem}
\begin{proof}
Put $x=R/L_*$. From \eqref{eq:Sbounds},
$nL_*<S_n<(n+1/2)L_*$. If
$A=\#\{n\ge1:S_n<R\}$ and $E_0=\#\{n\ge1:S_n=R\}$, then
\[
 A\ge\max(0,\lfloor x-1/2\rfloor),\qquad
 A+E_0\le\max(0,\lceil x\rceil-1).
\]
Since $R>2Q$, formula \eqref{eq:seamcount} gives
\begin{equation}\label{eq:ray-zero-count-error}
 -2<N_{B=0}-2x<1.
\end{equation}
This includes a fold equality, where $E_0=1$.

For the nonzero-$B$ count, Theorem~\ref{thm:ray-all-radii} gives
$2\max(0,\lceil M_0\rceil)$ and
\[
 M_0=\frac{R+2q_c}{L(m_c)}-1,\qquad
 q_c=F(\arctan(R/(2c_\alpha))\mid m_c)
       -E(\arctan(R/(2c_\alpha))\mid m_c).
\]
The defining integral is positive and its amplitude is below $\pi/2$,
so $0<q_c<L(m_c)/4$. The ceiling bound, including the case
$M_0\le0$, therefore yields
\[
 -2<N_{B\ne0}-2R/L(m_c)<1.
\]
Adding this to \eqref{eq:ray-zero-count-error} proves
\eqref{eq:ray-bounded-count-error} without a large-radius assumption.

As $\theta\downarrow0$, the defining complete integral gives
$L(m)=\pi m+O(m^2)$, while
\[
 s^2=\theta^2/4+O(\theta^4),\qquad
 m_c=\tfrac14(1+4/R^2)\theta^2+O_R(\theta^4).
\]
It follows that
\[
 \frac{2R}{L_*}=\frac{8R}{\pi\theta^2}+O_R(1),\qquad
 \frac{2R}{L(m_c)}=\frac{8R^3}{\pi(R^2+4)\theta^2}+O_R(1).
\]
All expansions and the condition $R>\max(2s,2Q)$ are uniform on the
stated compact radius intervals. This proves
\eqref{eq:ray-small-heading-count}. Finally,
$R_0^{\rm cusp}\to0$ by Proposition~\ref{prop:ray-cusp-scale};
Corollary~\ref{cor:ray-cut-status} then gives exactly two minimizers.
\end{proof}

\begin{remark}[At most one short-minimum comparison]\label{rem:ray-fast-count}
The exact ray count does not require evaluating every threshold.
For $R>2Q$, let $x=R/L_*$, $\rho=2Q/L_*$, and
$N_{\rm sure}=\max(0,\lfloor x-1/2\rfloor)$. The first $N_{\rm sure}$ short minima
are strictly below $R$. Any undecided index lies in
\[
 \{n\ge1:x-1/2<n<x-\rho\}.
\]
Since $0<\rho<1/2$, this set has at most one integer. If it is empty,
$N_{B=0}=1+2N_{\rm sure}$; otherwise add $0$, $1$, or $2$ according as
$R<S_n$, $R=S_n$, or $R>S_n$. Then add the nonzero-$B$ ceiling count
when $R>2s$; for $R\le2s$ that count is zero. For $R\le2Q$
there are no rotating sources. Thus the number of threshold comparisons
does not grow with the period count. Exact numerical decisions near a
threshold still require validated enclosures or an analytic equality.
\end{remark}

\begin{corollary}[Finiteness on the target seam]\label{cor:seam-finite}
For every $X\ne0$ and $0<|\theta|<\pi$, the full rotating fiber is
finite. On the target seam this includes every positive radius, rather
than only the far-field regime. Uniform boundedness on compact target
families, including those crossing the seam, follows in
Proposition~\ref{prop:uniform-finite} below.
\end{corollary}
\begin{proof}
Corollary~\ref{cor:full-finite} treats targets off the seam.
Spatial negation and the reflection
\[
 (x,y,\theta)\longmapsto(-x,y,-\theta)
\]
reduce every remaining target to $(R\ee(\alpha),2\alpha)$ with
$R>0$ and $0<\alpha<\pi/2$.
The exact count in Theorem~\ref{thm:ray-all-radii} is finite by
\eqref{eq:Sbounds} and the cusp-threshold bound in that theorem.
\end{proof}

\begin{proposition}[Uniform finiteness through the seam]\label{prop:uniform-finite}
The cardinality of the full rotating fiber is uniformly bounded on every
compact subset of
\[
 \mathcal D=\{(X,\theta):X\ne0,\quad 0<|\theta|<\pi\}.
\]
In particular, the compact family may cross the target seam. The bound
is not effective and does not give an exact count.
\end{proposition}
\begin{proof}
Every such compact family is contained in a band
\[
 \mathcal T=\{r_0\le|X|\le r_1,\quad
                  \theta_0\le|\theta|\le\theta_1\},
 \qquad 0<r_0\le r_1,\quad 0<\theta_0\le\theta_1<\pi.
\]
Set $m_*=\sin^2(\theta_0/2)>0$ and choose an integer
$M>r_1/L(m_*)$. The circle equations and longitudinal decomposition
give $m\ge m_*$ and $n<r_1/L(m_*)$ for every positive period index;
period zero also satisfies $n<M$. For any of these sources,
\[
 \frac tk=A+\int_\psi^{\psi+t/k}\dn^2u\,du
       <r_1+4(n+1)E(m)\le r_1+2\pi M.
\]
Here the integration interval consists of $n$ full periods and a
strictly shorter residual. Thus
\[
 r_0\le t< T:=r_1+2\pi M,\qquad |c|\le C:=2/\sqrt{m_*},
\]
using the covector bound from Corollary~\ref{cor:full-finite} and unit
speed. These bounds hold also at seam targets and for both rotation signs.

In the original coordinates let $F(\gamma,c,t)=(x,y,\widetilde\theta)$
be the endpoint map with continuous heading. The Hamiltonian vector
field is analytic and its real solutions exist at all finite times:
$|\dot c|\le1$, while the horizontal controls are bounded by one.
Analytic dependence on initial data and time makes $F$ analytic on a
neighborhood of the compact box
$[0,4\pi]\times[-C,C]\times[r_0,T]$.
After affine rescaling, its restriction is a restricted analytic map.
Consider the graph with the additional restrictions
\[
 0\le\gamma<4\pi,\qquad c^2/2-\cos\gamma>1,
 \qquad F(\gamma,c,t)\in\mathcal T.
\]
It is definable in the o-minimal structure $\mathbb R_{\rm an}$ of
restricted analytic functions \cite[Section 2]{VDM}. The half-open
angle interval counts each normalized covector once. For rotating
sources the continuous heading lies in $(-\pi,\pi)$, so no additional
heading representatives enter this graph. By the bounds above, its
fiber over each $q\in\mathcal T$ is exactly the full rotating fiber.

Each such fiber is finite by Corollary~\ref{cor:seam-finite}. The
uniform bound on connected components of fibers in
$\mathbb R_{\rm an}$ \cite[4.4]{VDM} therefore bounds their cardinalities.
This theorem applies to definable sets without requiring closedness or
properness; removing the energy-one boundary causes no difficulty.
The band $\mathcal T$ is definable, and the original compact family
need not be. Restricting the resulting bound to that family proves
the claim.
\end{proof}


Figure~\ref{fig:seam} shows the branches \eqref{eq:seambranches} for
$\alpha=\pi/5$, with the minima $S_n$ for $n\ge1$ and the common endpoints $H_n$.

\begin{figure}[htbp]
\centering
\includegraphics[width=\textwidth]{../figures/counts_seam.pdf}
\caption{The seam branches $R_n^{\rm short}$ (solid) and $R_n^{\rm long}$
(dashed) of \eqref{eq:seambranches}, $n=0,\dots,3$, at $\alpha=\pi/5$,
against $m-\sin^2\alpha$ (logarithmic axis). For $n\ge1$, the short branch dips to its
minimum $S_n$ (triangles); the $n=0$ short branch decreases. Both branches meet at
$H_n=(n+\tfrac12)L_*$ as $m\to\sin^2\alpha$ (dots).}
\label{fig:seam}
\end{figure}

Adding Propositions~\ref{prop:offseam} and \ref{prop:seam} proves
Theorem~\ref{thm:main} for positive $\alpha$.
The identity-fixing isometry $(x,y,\theta)\mapsto(-x,y,-\theta)$,
which sends controls to $(-u_1,-u_2)$, reduces negative $\alpha$ to this
case. Notice that the scalar maximum at $B=0$ does not imply a full-map
fold there: the source chart of Proposition~\ref{prop:fold} is singular
on this seam. Proposition~\ref{prop:seam-singularities} instead uses a
regular source chart: the short-branch minima are full-map folds,
whereas each long branch has one cusp (Corollary~\ref{cor:seam-cusps}). Neither
classification is needed for the count.

\section{Fixed-cell sources near the separatrix}

Return to the generic case $d_0\ne0$. Over the four arms, the coefficient
$c$ in \eqref{eq:tail} has values zero, one half, one, and one half.

\begin{proposition}[Fixed-cell scale]\label{prop:scale}
Fix a nonnegative integer $n$ and an endpoint arm with $n+c>0$.
For all sufficiently large $R$ there is exactly one source of cell $n$
on that arm. Define
\[
 C_{\sigma}=\sigma_0\{\atanh(\sin a_0)-\sin a_0\}
 -\sigma_1\{\atanh(\sin b_0)-\sin b_0\}.
\]
Its modulus satisfies
\begin{equation}\label{eq:scale}
 1-k=8e^{-2}\exp\!\left(-\frac{R+C_{\sigma}}{2(n+c)}\right)
 \bigl(1+O(R^{-1})\bigr).
\end{equation}
There are exactly three cell-zero sources eventually, with exponential
rates $e^{-R},e^{-R/2},e^{-R}$. For each fixed $n\ge1$, there are
eventually four sources. These are fixed-cell statements, not estimates
uniform in $n$ growing with $R$.
\end{proposition}
\begin{proof}
Existence and uniqueness follow from the monotone arm, its limit $-c$,
and its diverging maximum. A fixed-level root must tend to its endpoint:
otherwise the positive compact profile in \eqref{eq:profile} would force
$N_R$ to diverge. The short-integral limit is
$q(s,1)=\atanh(s)-s$, with bounded derivatives on the strict angle
range. Hence along the root
$P_R=R+C_{\sigma}+O(R^{-1})+O(1-m)$.
The standard expansion \cite[Section 19.12]{DLMF} reads
\[
 L(m)=2\log\frac{16}{1-m}-4
       +O\bigl((1-m)|\log(1-m)|\bigr).
\]
The equation $P_R=(n+c)L$ first gives exponential decay of $1-m$.
Substitution back into the remainders yields
$1-m=16e^{-2}\exp(-(R+C_{\sigma})/(2(n+c)))
(1+O(R^{-1}))$. Dividing by $1+k$ proves \eqref{eq:scale}.
The $c=0$ arm approaches level zero from above, so it contributes no
cell-zero root. The other three arms give the stated rates.
\end{proof}

\section{The full normal fiber by heading winding}
\label{sec:winding}

We now retain all pendulum regimes. In the usual vertical coordinates,
$\dot\gamma=c$, $\dot c=-\sin\gamma$ and the pendulum energy is
$\mathsf E=c^2/2-\cos\gamma$. The oscillating stratum $C_1$ has
$-1<\mathsf E<1$; $C_3$ is the separatrix $\mathsf E=1$, $c\ne0$;
$C_4$ and $C_5$ are the equilibria $\mathsf E=-1$ and
$\mathsf E=1$, respectively. These definitions, the endpoint formulas,
and the cut times are the standard ones in \cite{MS,SC,SS}.
For $|\theta|<\pi$, use the smooth target coordinates
\begin{equation}\label{eq:UVW}
 W=\tan(\theta/2),\qquad
 U=-\frac{y-xW}{2},\qquad V=-\frac{x+yW}{2}.
\end{equation}
The letter $U$ here is a target coordinate and no longer the auxiliary
quantity in \eqref{eq:zdata}. Write $\operatorname{am}$ for the continuously
lifted Jacobi amplitude, so $\operatorname{am}(u+2K)=\operatorname{am}(u)+\pi$.

\begin{lemma}[All positive oscillating winding cells]\label{lem:c1cells}
Fix a finite-chart target with $(U,V)\ne(0,0)$ and $W\ne0$.
For every integer $n\ge1$ and sign $s_1\in\{1,-1\}$ there is exactly
one $C_1$ source in
\[
 p=t/2=2nK+r,\qquad -K<r<K.
\]
It is regular and its heading winding is $\ell=-s_1n$.
\end{lemma}
\begin{proof}
A $C_1$ source has phase $\varphi$ modulo $4K$, modulus $0<k<1$,
time $t>0$, and sign $s_1$. Set
$\tau=\varphi+p$, $\eta=\operatorname{am}\tau$ modulo $2\pi$,
$q=p-\mathcal E(p)$, and
$f_2=k^2\cn p\sn p+\dn p\,q$.
The published endpoint equations give
\begin{equation}\label{eq:c1chart}
 U=\frac{q\cos\eta}{k\cn p},\quad
 V=-s_1\frac{f_2\sin\eta}{k\cn p},\quad
 W=-s_1\frac{\sn p\sqrt{1-k^2\sin^2\eta}}{\cn p}.
\end{equation}
There is no finite-chart source with $\cn p=0$: its principal heading is
$\pi$. Direct differentiation yields
\[
 \left(\frac{f_2}{\dn p}\right)'=\frac{k^2\cn^2p}{\dn^2p}.
\]
Its initial value is zero and its derivative is positive on a set of
positive measure in every interval, so $f_2>0$ for all $p>0$.

We first prove the assertion for $UV\ne0$. Put
$Q=\max(2,\sqrt{U^2+4V^2})$ and $z=q/k$.
Since $\dn p\ge|\cn p|$, when $q>k^2$ we have
$f_2\ge|\cn p|(q-k^2)$. Recovering sine and cosine from
\eqref{eq:c1chart} therefore gives
\[
 1\le\frac{U^2}{z^2}+\frac{V^2}{(z-k)^2}.
\]
If $z\ge2k$, this implies $z^2\le U^2+4V^2$; otherwise $z<2$.
The omitted case $q\le k^2$ gives $z\le k<1$. Thus $q/k\le Q$.
Since $q_p=k^2\sn^2p\ge0$, in cell $n$ this implies
\begin{equation}\label{eq:c1bound}
 (2n-1)\frac{K(k^2)-E(k^2)}k\le Q.
\end{equation}
The function $H(k)=(K(k^2)-E(k^2))/k$ increases from zero to infinity:
\[
 H'(k)=\frac{E-(1-k^2)K}{k^2(1-k^2)}>0,
\]
where the numerator equals
$k^2\int_0^{\pi/2}\cos^2u/\sqrt{1-k^2\sin^2u}\,du$.

Fix $n,s_1$ and the signs of the nonzero target coordinates. Use
$\rho=\operatorname{am}r\in(-\pi/2,\pi/2)$, $\eta$, and $k$ as source
coordinates. Equations \eqref{eq:c1chart} imply
\[
 \operatorname{sgn}U=(-1)^n\operatorname{sgn}(\cos\eta),\quad
 \operatorname{sgn}V=-s_1(-1)^n\operatorname{sgn}(\sin\eta),\quad
 \operatorname{sgn}W=-s_1\operatorname{sgn}\rho.
\]
These signs select one half-interval for $\rho$, one quadrant for
$\eta$, and $0<k<1$: a nonempty connected source box. At fixed $k$,
$p=2nK+F(\rho\mid k^2)$ and
$\varphi=F(\eta\mid k^2)-p$ have nonsingular Jacobian in
$(\rho,\eta)$. The absence of $C_1$ conjugate points
\cite[Theorem 2.1]{SC} makes the target map a local diffeomorphism.

This map is proper over the chosen target octant. For targets in a compact
subset, \eqref{eq:c1bound} bounds $k$ away from one. The equation
\[
 |W|=|\tan\rho|\sqrt{1-k^2\sin^2\eta}
\]
then keeps $\rho$ away from both ends of its half-interval.
If $k\to0$, with $n$ fixed and $\rho$ in this compact interval,
uniformly $q,f_2=O(k^2)$, whence $U,V=O(k)$, contradicting their
positive lower absolute bounds. Thus $k$ also stays away from zero.
The first two target equations now keep $\eta$ away from its quadrant
edges. All inverse images of compact target sets are therefore compact.
A proper local diffeomorphism has nonempty open and closed image here,
so it covers the entire connected octant. The octant is simply connected
and the source box connected, hence this covering has exactly one sheet.

For $V=0$, positivity of $f_2$ forces $\sin\eta=0$, and $U\ne0$ gives
$\cos\eta=(-1)^n\operatorname{sgn}U$. The remaining equations reduce to
\begin{equation}\label{eq:c1seam}
 \rho=\arctan(-s_1W),\quad r=F(\rho\mid k^2),\quad
 \frac{q}{k\cos\rho}=|U|.
\end{equation}
Define
$I_\rho(k)=\int_0^\rho\sin^2u/\sqrt{1-k^2\sin^2u}\,du$ and
$I_*=I_{\pi/2}$. Then $q/k=k(2nI_*+I_\rho)$.
For $\rho\ge0$ the bracket is a positive sum of increasing integrals;
for $\rho<0$ it is
\[
 (2n-1)I_*+\int_{|\rho|}^{\pi/2}
 \frac{\sin^2u}{\sqrt{1-k^2\sin^2u}}\,du.
\]
The left side of the last equation in \eqref{eq:c1seam} therefore
increases strictly from zero to infinity. It gives exactly one modulus,
and \eqref{eq:c1seam} recovers exactly one source. Its regularity follows
from the same published nonconjugacy theorem.

It remains to treat $U=0$, $V\ne0$. Put
\[
 w=|W|,\quad h_k=\sqrt{1-k^2},\quad
 L_k=\sqrt{h_k^2+w^2},\quad C_w=\sqrt{1+w^2},\quad
 \epsilon=-s_1\operatorname{sgn}W.
\]
Since $q>0$, \eqref{eq:c1chart} forces $\cos\eta=0$ and
$\sin\eta=-s_1(-1)^n\operatorname{sgn}V$.
It also gives $\rho=\epsilon a(k)$, where $a(k)=\arctan(w/h_k)$.
The remaining scalar equation is
\begin{equation}\label{eq:c1Useam}
 \mathcal F_\epsilon(k)
 =C_w k(2nI_*+\epsilon I_{a(k)})
       +\frac{\epsilon kw}{L_k}=|V|.
\end{equation}
Define
$A_a=\int_0^a\sin^2u/(1-k^2\sin^2u)^{3/2}\,du$ and
$A_*=A_{\pi/2}$. Differentiating, including the moving endpoint, yields
\begin{equation}\label{eq:c1Useam-derivative}
 \mathcal F_\epsilon'
 =C_w(2nA_*+\epsilon A_{a(k)})
       +\frac{\epsilon w}{h_k^2L_k}.
\end{equation}
The sign is positive for $\epsilon=1$. For $\epsilon=-1$, $n\ge1$,
use $2nA_*-A_{a(k)}\ge A_*$ and
\[
 h_k^2A_*=J_k,\qquad
 J_k=\int_0^{\pi/2}\frac{\cos^2u}{\sqrt{1-k^2\sin^2u}}\,du.
\]
The identity $h_k^2A_*=J_k$ follows by integrating the derivative
of $\sin u\cos u/\sqrt{1-k^2\sin^2u}$ between its zero endpoint
values. Cauchy--Schwarz gives
\[
 J_kE\ge\left(\int_0^{\pi/2}\cos u\,du\right)^2=1.
\]
Also $E<1+h_k$, by integrating the strict pointwise inequality
\[
 \sqrt{\cos^2u+h_k^2\sin^2u}<\cos u+h_k\sin u,\qquad 0<u<\pi/2.
\]
Finally
\[
 C_w^2L_k^2-(1+h_k)^2w^2=(w^2-h_k)^2\ge0,
 \qquad J_k>\frac1{1+h_k}\ge\frac w{C_wL_k}.
\]
These inequalities prove strict positivity in
\eqref{eq:c1Useam-derivative} for the negative sign too.
The scalar tends to zero as $k\to0$ and to infinity as $k\to1$,
since $q\ge(2n-1)(K-E)$ while $kw/L_k$ stays bounded.
Thus it supplies exactly one modulus and one source.
Again full-map regularity follows from $C_1$ nonconjugacy.

Finally the continuous heading is
\begin{align*}
 \widetilde\theta
 &=s_1\{\operatorname{am}(\tau-p)-\operatorname{am}(\tau+p)\}\\
 &=-2\pi s_1n+s_1\{\operatorname{am}(\tau-r)-\operatorname{am}(\tau+r)\}.
\end{align*}
The residual is strictly in $(-\pi,\pi)$ because $|r|<K$, hence equals
the principal target heading. This proves $\ell=-s_1n$.
\end{proof}

\begin{proposition}[Nonzero winding at every radius]\label{prop:global-winding}
If $X\ne0$ and $0<|\theta|<\pi$, every nonzero integer winding
class contains exactly one normal source. It belongs to $C_1$ and
is regular.
\end{proposition}
\begin{proof}
The target chart \eqref{eq:UVW} is invertible, so $(U,V)\ne(0,0)$.
Lemma~\ref{lem:c1cells} exhausts all positive $C_1$ cells; their
winding is $-s_1n$. The remaining $C_1$ cell is $0<p<K$ and has
winding zero. The continuous $C_2$ heading is the difference of two
strict principal arcsines in \eqref{eq:ab}, so it has winding zero.
For $C_3$, the continuous heading is
\[
 \widetilde\theta(t)
 =s_1\{\arcsin(\tanh\varphi)-\arcsin(\tanh(\varphi+t))\},
\]
also strictly in $(-\pi,\pi)$. Finally $C_4$ has $X=0$ and $C_5$
has principal heading zero, excluding both. Thus no other stratum
contributes a nonzero winding source.
\end{proof}

\begin{lemma}[Effective separatrix exclusion]\label{lem:c3exclude}
For the angles of Theorem~\ref{thm:main}, set
$h=(a_0-b_0)/2$ and $\beta=(a_0+b_0)/2$.
There is no $C_3$ source of $q_R$ if
\begin{equation}\label{eq:c3radius}
 R>\frac{2\sin^2h}{\cos(|\beta|+|h|)}.
\end{equation}
\end{lemma}
\begin{proof}
Here $|\beta|+|h|<\pi/2$ and \eqref{eq:UVW} gives
$W=\tan h$, $U=Ru$, $V=Rv$, where
\[
 u=\frac{\cos\beta}{2\cos h},\quad
 v=-\frac{\sin\beta}{2\cos h},\quad
 g=u-|W||v|=\frac{\cos(|\beta|+|h|)}{2\cos^2h}>0.
\]
For $C_3$, the exact endpoint equations in midpoint variables $p=t/2>0$
and $\tau$ are
\[
 U=s_2\frac{(p-\tanh p)\cosh p}{\cosh\tau},\quad
 V=-s_1s_2p\tanh\tau,\quad W=-s_1\frac{\sinh p}{\cosh\tau}.
\]
Eliminating $\tau$ and the signs gives the necessary equations
\begin{equation}\label{eq:c3surface}
 \frac{V^2}{p^2}+\frac{W^2}{\sinh^2p}=1,\qquad
 |U|=|W|(p\coth p-1).
\end{equation}
The first left side strictly decreases from infinity to zero, so its
root is unique. It lies in $(|V|,|V|+|W|)$: for $V\ne0$ the lower
endpoint gives a value above one and $\sinh p>p$ makes the upper value
less than one; for $V=0$ the root is $\operatorname{arsinh}|W|<|W|$.
Since $e^{2p}>1+2p$, one has $p\coth p-1<p$. Thus
$Ru-|W|(p\coth p-1)>Rg-W^2>0$ under \eqref{eq:c3radius}, excluding
the second equation. This is a sufficient bound for $C_3$ alone.
\end{proof}

\begin{proof}[Proof of Theorem~\ref{thm:winding}]
The target has $U>0$, $W\ne0$, and $V=0$ precisely when $\beta=0$.
Lemma~\ref{lem:c1cells} supplies exactly one source in every nonzero
winding class. The bounded continuous heading in \eqref{eq:ab} shows
that every $C_2$ source has winding zero. It remains to exclude $C_1$
cell zero and the other regimes.

Suppose first $\beta\ne0$, equivalently
$\sin^2a_0\ne\sin^2b_0$. Proposition~\ref{prop:scale} supplies an
actual cell-zero source on each opposite-cosine arm
$\sigma_1=-\sigma_0$. With the source angles $a,b$ and signed short
integrals
$u_a=F(\arcsin(\sin a/k)\mid m)$ and
$u_b=F(\arcsin(\sin b/k)\mid m)$, phase reconstruction gives exactly
\[
 \delta=2K-\sigma_0(u_a+u_b).
\]
As $R\to\infty$, the sum tends to
$\atanh(\sin a_0)+\atanh(\sin b_0)$, which has the nonzero sign of
$\beta$ by oddness and strict monotonicity of $\atanh(\sin a)$.
The arm $\sigma_0=\operatorname{sgn}\beta$ therefore has $t=k\delta<2kK$
eventually. For
$f_1=\cn p(\mathcal E(p)-p)-\dn p\sn p$, we have $f_1<0$ on $(0,K]$;
thus its first positive root $p_1$ satisfies $p_1>K$.
The published cut time is $2kp_1$ \cite[(2.2)]{SS}.
This source is strictly pre-cut, and the known synthesis makes it the
unique minimizer \cite[Theorems 3.1--3.2]{SS}.
Any $C_1$ cell-zero source would have $t=2p<2K$, its own cut time
\cite[(2.1)]{SS}, and would be another strictly pre-cut
minimizer. This contradiction excludes cell zero.

For $\beta=0$, Proposition~\ref{prop:offseam} gives two distinct
cell-zero sources with $B\ne0$ eventually. They are exchanged by
$B\mapsto-B$, with equal modulus and time. Put $p=t/(2k)$ and
$\tau=\psi+p$. The $C_2$ target formula and Jacobi addition identity give
\[
 V=\frac{kf_1(p)\sn\tau}{\dn p},\qquad
 B=\frac{2k\sn\tau\sn p\dn\tau\dn p}
 {1-m\sn^2\tau\sn^2p}.
\]
The denominator and the $\dn$ factors are positive. Thus $B\ne0$
implies $\sn\tau\ne0$, so $V=0$ forces $f_1(p)=0$.
Cell zero means $0<p<2K$. There is exactly one such root: besides
$f_1<0$ on $(0,K]$, on $(K,2K)$ we have
\[
 (f_1/\cn)'=-\dn^2p/\cn^2p<0,
\]
and the quotient decreases from $+\infty$ to $-2(K-E)<0$.
Hence both sources are at the first cut time $2kp_1$ and minimize.
The published synthesis gives exactly two minimizers at a Maxwell
endpoint \cite[Theorems 3.2 and 3.4]{SS}. A strictly pre-cut $C_1$
cell-zero source would contradict this pair. Negative $h$ follows
either from the same identities or from the isometry used in Section~\ref{sec:seam}.

Lemma~\ref{lem:c3exclude} excludes $C_3$ eventually. The $C_4$ image
has $x=y=0$, and the $C_5$ image has $y=\theta=0$, so neither reaches
the present target. Theorem~\ref{thm:main} now gives the finite
zero-winding count; the remaining classes each contain one regular
$C_1$ source. Spatial negation preserves winding and all counts.
No regularity conclusion for every $C_2$ source is needed.
\end{proof}

\subsection{The exact rotating image and its parity}
\label{sec:rotating-image}

We retain the target coordinates \eqref{eq:UVW}, with
$X\ne0$ and $0<|\theta|<\pi$. The hyperbolic source formulas are
classical \cite[Section 3.3]{MS}; the comparison below applies to
every positive-time rotating source, including all post-cut periods.

\begin{lemma}[The separatrix graph]\label{lem:separatrix-graph}
For $V\in\mathbb R$ and $w>0$, let $P=P(V,w)>0$ be the unique root of
\[
 \frac{V^2}{P^2}+\frac{w^2}{\sinh^2P}=1,
 \qquad \mathcal F(V,w)=w(P\coth P-1).
\]
The complete $C_3$ image in the stated domain is
$\Sigma=\{|U|=\mathcal F(V,|W|)\}$, with exactly one $C_3$ source
at each target. The function $\mathcal F$ is analytic for $w>0$, even
in $V$, increasing for $V>0$, and satisfies
\begin{equation}\label{eq:separatrix-slope}
 0\le\partial_V\mathcal F(V,w)<w\qquad(V\ge0).
\end{equation}
\end{lemma}
\begin{proof}
Existence and uniqueness of $P$ follow from the strictly decreasing
left side, as in \eqref{eq:c3surface}. Its $P$ derivative is nonzero,
so the implicit function theorem gives analyticity. Necessity of the
graph condition follows from the same endpoint equations. For sufficiency,
set $s_1=-\operatorname{sgn}W$, $s_2=\operatorname{sgn}U$ and
\[
 \tau=\atanh(-s_1s_2V/P),\qquad \varphi=\tau-P,\qquad t=2P.
\]
The root equation gives $\cosh\tau=\sinh P/|W|$. Substitution
recovers all three target coordinates. These signs, time and phase are
forced and determine one normalized hyperbolic covector \cite{MS}.
Also $P\coth P>1$, so $U\ne0$ on the graph.

For $V>0$, put $a=w^2/\sinh^2P\in(0,1)$ and $H=P\coth P-1$.
Then $V=P\sqrt{1-a}$ and
\[
 \frac{dV}{dP}=\frac{1+aH}{\sqrt{1-a}}>1,
 \qquad 0<H'=\frac{\sinh P\cosh P-P}{\sinh^2P}<1.
\]
The last inequalities follow from $\sinh(2P)>2P$ and
$1-e^{-2P}<2P$. This proves \eqref{eq:separatrix-slope} for $V>0$;
evenness and analyticity give derivative zero at $V=0$.
\end{proof}

\begin{lemma}[An all-period rotating barrier]\label{lem:rotating-barrier}
Every positive-time $C_2$ source in the stated target domain satisfies
\begin{equation}\label{eq:rotating-barrier}
 |U|>\mathcal F(V,|W|).
\end{equation}
In particular, its full critical-value image is disjoint from $\Sigma$.
\end{lemma}
\begin{proof}
Use $p=t/(2k)>0$, $\tau=\psi+p$ and
$\eta=\operatorname{am}\tau$ modulo $2\pi$. Write
\[
 s=\sn p,\quad c_p=\cn p,\quad d=\dn p,\quad
 q_p=p-\epsilon_p=\int_0^p k^2\sn^2u\,du.
\]
The original endpoint formulas \cite[Section 3.3]{MS} and Jacobi
addition give, in the coordinates \eqref{eq:UVW},
\begin{equation}\label{eq:rotating-midpoint-chart}
 U=s_2\frac{q_p\dn\tau}{d},\qquad
 V=s_2\frac{k(-c_pq_p-ds)\sin\eta}{d},\qquad
 W=-\frac{ks\cos\eta}{d}.
\end{equation}
These $U,V$ are the scaled chart coordinates \eqref{eq:UVW}, not
the unscaled components used temporarily in \eqref{eq:ray-maxwell-components}.
For example, $\cos(\theta/2)=d/\sqrt\Delta$ in that formula,
which supplies the scaling and the sign of $V$.

Put $r=k|s|$, $b=k|c_p|$ and $Q(r)=\atanh r-r$. Since
$d^2-b^2=1-k^2>0$, one has $b<d$ and
\[
 \frac{d}{dp}\{q_p-Q(k|\sn p|)\}
 =k^2s^2\left(1-\frac{k\operatorname{sgn}(s)c_p}{d}\right)>0
 \quad(s\ne0).
\]
The derivative extends as zero at the zeros of $s$. The difference
starts at zero and increases strictly over every nontrivial interval;
thus $q_p>Q(r)$ for every $p>0$, with no period restriction.

Set $z=|\sin\eta|$, $h=|\cos\eta|$,
$D_\tau=\sqrt{1-k^2z^2}$ and $q_*=Q(r)$.
Strict nonzero heading gives $r,h>0$. Equation
\eqref{eq:rotating-midpoint-chart} implies
\[
 |U|=\frac{q_pD_\tau}{d},\qquad w:=|W|=\frac{rh}{d},
 \qquad |V|\le z\left(r+\frac{bq_p}{d}\right)=:\overline V(q_p).
\]
At $P_*:=\atanh r$, one has
$\overline V(q_*)\le z(r+q_*)=zP_*$ and
\[
 \frac{(zP_*)^2}{P_*^2}+\frac{w^2}{\sinh^2P_*}=z^2+h^2=1,
 \qquad \mathcal F(zP_*,w)=\frac{hq_*}{d}.
\]
Monotonicity in $|V|$ and \eqref{eq:separatrix-slope} therefore give
\[
 \mathcal F(V,w)\le \frac{hq_*}{d}
       +\frac{wzb}{d}(q_p-q_*).
\]
Subtracting this from $|U|$, and using $D_\tau\ge h$, $b/d<1$, yields
\begin{equation}\label{eq:rotating-barrier-margin}
 |U|-\mathcal F(V,w)
 \ge\frac{q_*(D_\tau-h)+(q_p-q_*)(D_\tau-wzb)}d
 \ge\frac{(q_p-q_*)h(1-rz)}d>0.
\end{equation}
The non-strict comparison is intentional: equality holds in the last
comparison when $z=0$. Its lower bound is still strictly positive.
All phase and rotation signs are retained in this argument.
\end{proof}

\begin{theorem}[Exact rotating image and regular-fiber parity]
\label{thm:rotating-image-parity}
For $X\ne0$ and $0<|\theta|<\pi$, a rotating source exists if and
only if $|U|>\mathcal F(V,|W|)$. Every target in this region has a
finite nonempty rotating fiber. If every rotating preimage is regular,
the fiber cardinality is odd. At critical values only existence and
finiteness are asserted.
\end{theorem}
\begin{proof}
The preceding lemma proves necessity. For each fixed pair of signs of
$U,W$, the strict region is analytically parametrized by
\[
 (V,w,z)\in\mathbb R\times(0,\infty)^2,
 \qquad W=\pm w,\quad U=\pm(\mathcal F(V,w)+z).
\]
It is therefore diffeomorphic to $\mathbb R^3$. Restrict the rotating
endpoint map to the full inverse of one such target component.
The source is an open three-dimensional manifold, with $\gamma$
modulo $4\pi$; it is not assumed compact or connected.

This map is proper. For a compact target subset, the bounds in the
proof of Proposition~\ref{prop:uniform-finite} bound all initial
covectors and positive times, including every period. A limiting
source has energy at least one. At energy one, $C_5$ has heading zero
and $C_3$ maps to $\Sigma$, both excluded by the compact target subset.
Thus every source limit remains in the restricted domain. The angle's
half-open representative introduces no physical boundary.

There is a regular target with exactly one source in each component.
Indeed any adjacent point of $\Sigma$ has a unique $C_3$ source,
regular by \cite[Theorem 2.4]{SC}, and no $C_2$ source by
Lemma~\ref{lem:rotating-barrier}. The same source bounds and local
injectivity exclude any additional $C_2$ inverse branches arriving
near it. The local energy-greater-than-one side lies in the strict
region by the barrier, giving the required single-source base value.

We use degree modulo two \cite[Section 4]{Milnor}, with the needed
proper-map extension made explicit. For two regular target values,
the component's $\mathbb R^3$ coordinates give an isotopy $h_t$
carrying one to the other. The homotopy $h_t\circ\operatorname{Exp}$
is proper: inverse sets over compact targets lie in the product of
$[0,1]$ with a compact inverse set of $\operatorname{Exp}$.
A nearby regular value of the homotopy has inverse a compact
one-dimensional manifold whose boundary is the two endpoint fibers.
Sard's theorem and local constancy at regular values preserve those
endpoint counts. The even number of boundary points gives equal
parity. Hence every regular value has odd count, since the base count
is one. A missing target would be a regular value with empty fiber,
contradicting this parity. This proves surjectivity on all four
components. Finiteness, also at critical values, follows from
Corollary~\ref{cor:seam-finite}.
\end{proof}

\begin{corollary}[The count next to the separatrix]
\label{cor:separatrix-collar}
Every target of $\Sigma$ has a neighborhood in which the exact rotating
count is $\ind_{\{|U|>\mathcal F(V,|W|)\}}$. The source entering
$C_2$ has actual period zero. A compact subset of $\Sigma$ admits
such a neighborhood, without an effective width bound.
\end{corollary}
\begin{proof}
The local inverse and source-limit exhaustion in the preceding proof
give the count. Time remains bounded while $K\to\infty$ on the
entering rotating branch, so its actual period is zero. A finite
cover gives the compact-subset assertion.
\end{proof}

\begin{remark}
Oddness cannot be extended to critical values: at the first short minimum
$R=S_1$ on the half-heading ray, \eqref{eq:ray-global-count} gives two
zero-transverse sources plus an even nonzero-transverse contribution.
The image criterion and parity do not determine the exact count at an
arbitrary target.
\end{remark}


\begin{corollary}[All-radius zero-winding regimes]
\label{cor:zero-winding-regimes}
Let $X\ne0$, $0<|\theta|<\pi$, and put
$\mathcal G(q)=|U|-\mathcal F(V,|W|)$.
The complete zero-winding normal fiber has the following form:
\[
\begin{array}{c|c|c}
\text{condition}&\text{stratum}&\text{cardinality }N_0(q)\\ \hline
\mathcal G(q)<0&C_1\text{, cell zero}&1\\
\mathcal G(q)=0&C_3&1\\
\mathcal G(q)>0&C_2&\nu_2(q)<\infty.
\end{array}
\]
In the first two rows the single source is regular and is the unique
minimizer. In the third row every minimizer belongs to actual $C_2$
cell zero. Every nonzero winding class contains one regular $C_1$
source, so the full normal fiber is countably infinite at every target
in this domain. No exact arbitrary-target formula for $\nu_2$ is asserted.
\end{corollary}
\begin{proof}
We use the classical cut times $t_{\rm cut}=2K$ on $C_1$,
$t_{\rm cut}=2kp_1$ on $C_2$, and $t_{\rm cut}=\infty$ on $C_3$
\cite[Theorem 3.3]{SC}\cite[(2.1)--(2.3)]{SS}, together with
the published optimal synthesis \cite[Theorem 3.2]{SS}.
In particular, a strictly pre-cut source is the unique minimizer
of its target. Every nontrivial target admits a normal minimizer
\cite[Sections 1--2]{MS}.

For each fixed sign $s_1$, the full $C_1$ cell-zero domain is
parametrized by $0<k<1$, $0<p/K<1$ and
$\eta\in\mathbb R/(2\pi\mathbb Z)$; it is connected.
The continuous heading has absolute value
$\int_{\tau-p}^{\tau+p}\dn u\,du\in(0,\pi)$, because
$2p<2K$ and every full $2K$ integral is $\pi$.
Also $q_p=p-\mathcal E(p)>0$ and
$f_2=k^2\cn p\sn p+\dn p\,q_p>0$ on $0<p<K$.
Equation~\eqref{eq:c1chart} therefore gives $(U,V)\ne(0,0)$.
Thus its entire image lies in the stated target domain.

That image cannot meet $\Sigma$. Such a target would have both a
strictly pre-cut $C_1$ source and the $C_3$ source of
Lemma~\ref{lem:separatrix-graph}, also strictly pre-cut, contradicting
uniqueness in the synthesis. Consequently $\mathcal G$ has a constant
nonzero sign on each connected $C_1$ cell-zero domain. At
$\eta=\pi/2$, \eqref{eq:c1chart} gives $U=0$, $V\ne0$ and
$W\ne0$, hence $\mathcal G<0$. This proves that the entire
$C_1$ cell-zero image lies in $\{\mathcal G<0\}$.

Now choose a normal minimizer of any target in the stated domain.
A positive $C_1$ winding cell has $p>K$, hence lies after its cut
time and cannot supply this minimizer. The equilibria $C_4,C_5$
are excluded by $X\ne0$ and $\theta\ne0$, respectively.
If $\mathcal G<0$, the rotating barrier excludes $C_2$ and the
separatrix graph excludes $C_3$, so the minimizer must be in $C_1$
cell zero; its strict pre-cut uniqueness gives exactly one such source.
If $\mathcal G=0$, the unique $C_3$ source is the only zero-winding
source. If $\mathcal G>0$, neither $C_1$ cell zero nor $C_3$ can
occur, so every minimizer is rotating. Since $p_1<2K$, its time
satisfies $t\le2kp_1<4kK$, giving actual period zero.

The previous winding catalogue and rotating finiteness now give all
three rows and the full-fiber assertion. Regularity in the first two
rows is the known nonconjugacy of $C_1,C_3$ \cite[Theorems 2.1, 2.4]{SC}.
The minimizing existence argument also gives an alternative proof of
surjectivity onto $\{\mathcal G>0\}$ using only the rotating barrier,
without the degree argument used for parity.
\end{proof}


\subsection{Convexity and sharp radial existence intervals}
\label{sec:radial-existence}

The graph in Lemma~\ref{lem:separatrix-graph} comes from the classical
hyperbolic formulas \cite[Section 3.3]{MS}. The following geometry,
combined with the all-period barrier, determines exactly where a fixed
target ray admits rotating sources. It does not determine their number.

\begin{lemma}[Strict convexity at fixed heading]\label{lem:image-convexity}
For each fixed $w>0$, $\mathcal F(V,w)$ is strictly convex in $V$.
The two spatial components of the rotating image at a fixed strict
heading are open convex sets. There is a unique $V_*(w)>0$ such that
\[
 \mathcal F(V_*,w)=V_*\partial_V\mathcal F(V_*,w).
\]
Writing $\sigma(w)=\mathcal F(V_*,w)/V_*$, one has
$0<\sigma(w)<w$. The quotient $\mathcal F(V,w)/V$ decreases strictly
from infinity to $\sigma(w)$ and then increases strictly to $w$.
Both $V_*$ and $\sigma$ are analytic for $w>0$.
\end{lemma}
\begin{proof}
Put $P_0=\operatorname{arsinh}w$, $a=w^2/\sinh^2P$ and
$H=P\coth P-1$. For $P>P_0$, parametrize the graph by
$V=P\sqrt{1-a}$ and $\mathcal F=wH$. Differentiation in $P$ gives
\[
 V'=\frac{1+aH}{\sqrt{1-a}}>0,\qquad
 H'=\coth P-P\operatorname{csch}^2P>0,\qquad
 H''=2H\operatorname{csch}^2P>0,
\]
\[
 V''=\frac{a}{(1-a)^{3/2}}
 \bigl\{(1-a)(H'-2H\coth P)-(1+aH)\coth P\bigr\}<0.
\]
For the last sign, with $x=2P$, observe that
\[
 \sinh^2P\,(H'-2H\coth P)
 =-\tfrac12(x\cosh x+2x-3\sinh x),
\]
\[
 x\cosh x+2x-3\sinh x
 =\sum_{j\ge2}\frac{(2j-2)x^{2j+1}}{(2j+1)!}>0.
\]
Thus $\mathcal F_{VV}=w(H''V'-H'V'')/(V')^3>0$ for $V>0$.
At $V=0$, implicit differentiation gives
\[
 P_V=0,\qquad P_{VV}=\frac{1}{P_0^2\coth P_0},\qquad
 \mathcal F_{VV}(0,w)=\frac{wH'(P_0)}{P_0^2\coth P_0}>0.
\]
Evenness proves the assertion on the whole line. The regions
$U>\mathcal F(V,w)$ and $U<-\mathcal F(V,w)$ are convex, and
the fixed-heading change from $(x,y)$ to $(U,V)$ is invertible and linear.

As $V\to\infty$, the identities
$P-V=Pa/(1+\sqrt{1-a})$, $V<P<V+w$ and
$H=P-1+2P/(e^{2P}-1)$ imply
\begin{equation}\label{eq:image-asymptote}
 \mathcal F=wV-w+O_w(Ve^{-2V}),\qquad
 \mathcal F_V=w+O_w(Ve^{-2V}).
\end{equation}
The derivative estimate follows directly from $\mathcal F_V=wH'/V'$.
Consequently $J(V)=\mathcal F-V\mathcal F_V$ starts at
$J(0)=\mathcal F(0,w)>0$, has derivative $-V\mathcal F_{VV}<0$
for $V>0$, and tends to $-w$. Its unique zero gives the claims about
$\mathcal F/V$, since $(\mathcal F/V)'=-J/V^2$. The slope bound
\eqref{eq:separatrix-slope} gives $0<\sigma<w$, and the implicit
function theorem at $J_V(V_*)<0$ gives analytic dependence.
\end{proof}

\begin{proposition}[Exact radial existence classification]
\label{prop:radial-existence}
Fix $0<|\theta|<\pi$ and a unit spatial direction
$\mathbf e=(e_x,e_y)$. For $q_R=(R\mathbf e,\theta)$, $R>0$, put
\[
 W=\tan(\theta/2),\quad w=|W|,\quad
 u=\tfrac12|e_y-We_x|,\quad v=\tfrac12|e_x+We_y|.
\]
If $v=0$, rotating sources exist precisely for
$R>\mathcal F(0,w)/u$. If $v>0$, put $b=u/v$. With $\sigma=\sigma(w)$
from Lemma~\ref{lem:image-convexity}, the complete classification is
\[
\begin{array}{c|c|c}
\text{direction}&\text{radii on }\Sigma&\text{radii with }C_2\text{ sources}\\ \hline
b\ge w&\text{one transverse }R_-&(R_-,\infty)\\
\sigma<b<w&\text{two transverse }R_-<R_+&(R_-,R_+)\\
b=\sigma&\text{one tangency }R_*=V_*/v&\varnothing\\
0\le b<\sigma&\varnothing&\varnothing.
\end{array}
\]
At every radius on $\Sigma$, the zero-winding source is the unique
regular $C_3$ source. In particular, the target-ray tangency is not
a conjugate point or a singularity of the exponential map.
\end{proposition}
\begin{proof}
The exact image criterion is $g(R)=Ru-\mathcal F(Rv,w)>0$.
Since $u^2+v^2=(1+w^2)/4$, the case $v=0$ has $u>0$ and is immediate.
For $v>0$, $g''=-v^2\mathcal F_{VV}<0$ and
$\lim_{R\downarrow0}g(R)=-\mathcal F(0,w)<0$; the latter is only
a scalar limit, since the spatial origin is excluded.
If $b\ge w$, $g'=v(b-\mathcal F_V)>0$. Equation~\eqref{eq:image-asymptote}
gives $g\to\infty$ for $b>w$ and $g\to w>0$ for $b=w$.
Hence there is one transverse crossing in either case. For $b<w$,
the last three rows follow from the unique strict minimum and the
two monotone branches of $\mathcal F(V,w)/V$.
At $b=\sigma$, $g$ touches zero at its unique maximum; it is negative
elsewhere. Corollary~\ref{cor:zero-winding-regimes} identifies the
sources at each boundary and in the complementary regions.
\end{proof}

The tangent parameter can be found without differentiating an inverse:
$P_*>\operatorname{arsinh}w$ is the unique root of
\[
 \sinh^2P_*=(1+w^2)P_*^2-w^2(P_*\coth P_*-1).
\]
Indeed, $H-PH'=P^2/\sinh^2P-1$ and $H^2+PH'=P^2-H$
reduce $J=0$ to this equation. Uniqueness follows from the lemma;
no monotonicity of the multiplied residual is needed.

In the directional notation of Theorem~\ref{thm:directions},
\[
 u^2-w^2v^2=\tfrac14(1+w^2)^2\Delta,\qquad
 \Delta=e_y(e_y\cos\theta-e_x\sin\theta).
\]
Thus $\Delta\ge0$ gives one onset of rotating existence.
For $\Delta<0$, a finite interval, one $C_3$ tangency, or neither
can occur. Finite intervals occur at every strict heading, since
$\sigma(w)<w$ and every nonnegative ratio $u/v$ is attained by a
direction. These existence thresholds do not identify where the eventual
count formulas become valid. At $v=0$ the threshold reduces to
$2(\atanh s-s)$, $s=w/\sqrt{1+w^2}$, the classical function
in \cite[(2.13), Lemma 2.3]{SS}. No novelty of that threshold or
of the minimizing synthesis is asserted.

\paragraph{A certified finite interval.}
Take $\theta=2\arctan(3/4)$ and $\mathbf e=(-1,-2)/\sqrt5$.
Then $\Delta=-4/25$, $U=V/2>0$ and $R=4V/\sqrt5$.
Strict interval gap signs at $V=1/10,1,4$ are negative, positive,
negative. Proposition~\ref{prop:radial-existence} therefore gives exactly
two boundary crossings. Outward interval bisection at 256 bits gives
\[
 R_-\simeq0.46581453181801560948,\qquad
 R_+\simeq5.22604376092815354863.
\]
The exact rational brackets and strict endpoint signs are recorded in
\nolinkurl{data/e105_radial_window.json}; their $V$ widths are below
$10^{-26}$. A separate 90-digit inversion in $P$ recovers both brackets,
and original Hamiltonian integrations corroborate the boundary targets.
Corollary~\ref{cor:exterior-sector-count} proves that exactly one rotating
source exists for $R_-<R<R_+$, with none outside or on its boundaries.

\begin{figure}[htbp]
\centering
\includegraphics[width=\textwidth]{../figures/counts_radial_window.pdf}
\caption{A finite rotating interval in an exterior direction of $\mathrm{SE}(2)$.
The boundary graph and its gap from the target ray are sampled illustrations;
the two crossing locations have separate interval certificates.
Corollary~\ref{cor:exterior-sector-count} proves count one in the shaded
interval and zero elsewhere. Both boundary targets have a regular $C_3$ source.}
\label{fig:radial-window}
\end{figure}


\subsection{Exact counts in an exterior sector}\label{sec:exterior-sector}

Exteriority in Theorem~\ref{thm:directions} is equivalent to
$|U|<|WV|$. The following restrictions apply at every radius and
retain the actual rotating period index $n=\lfloor t/(4kK)\rfloor$.

\begin{lemma}[Exterior source restrictions]\label{lem:exterior-periods}
Let a rotating source reach an exterior target, and put
$w=|W|$, $b=|U|/|V|$ and $m_h=w^2/(1+w^2)$.
Then $0<b<w$ and, with $m=k^2$, $p=t/(2k)$ and
$\mathcal Q(m)=K(m)-E(m)$,
\begin{equation}\label{eq:exterior-source-bounds}
 p-\mathcal E(p)<\frac{b}{w-b},\qquad
 n<\frac{b}{2(w-b)\mathcal Q(m_h)},\qquad
 2nK<p<(2n+1)K.
\end{equation}
If any rotating source exists at an exterior target, exactly one has
actual period zero. That source is regular and strictly pre-cut;
all additional rotating sources have positive actual period.
\end{lemma}
\begin{proof}
Write $\kappa=\sqrt{1-m}$, $s=\sn p$, $c=\cn p$, $d=\dn p$,
$q_p=p-\mathcal E(p)$, and
$z=|\sin\eta|$, $h=|\cos\eta|$, $D=\sqrt{1-mz^2}$ in the
midpoint chart. Its endpoint equations give
\[
 |U|=\frac{q_pD}{d},\quad
 |V|=\frac{kz|cq_p+ds|}{d},\quad w=\frac{k|s|h}{d}.
\]
Here $s,h,z$ are nonzero. Since $k|c|<d$ and $D\ge h$,
$wk|c|z<D$. Also
$D^2=\kappa^2+w^2d^2/s^2$, so $D>wd/|s|$.
The identity $q_pD=bkz|cq_p+ds|$ and $b<w$ therefore yield
\[
 q_p\le\frac{bkzd|s|}{D-bkz|c|}
 <\frac{b}{1-b/w}\frac{kzd|s|}{D}
 <\frac{b}{w-b}kzs^2<\frac{b}{w-b}.
\]
The heading constraint gives $ms^2\ge m_h$.
Moreover $q_p\ge2n\mathcal Q(m)$ and $\mathcal Q$ is increasing,
which proves the index bound.

We next exclude the second half of each period. The complete-integral
derivatives \cite[Section 19.4]{DLMF}, converted to the parameter $m$, give
\[
 \mathcal Q'(m)=\frac{E(m)}{2(1-m)},\qquad
 \frac{d}{dm}\bigl[\mathcal Q(m)-(1-\kappa)\bigr]
 =\frac{E(m)-\kappa}{2\kappa^2}>0.
\]
Since $\mathcal Q(0)=0$ and $E(m)>\pi\kappa/2>\kappa$,
$\mathcal Q(m)>1-\kappa=: \lambda$.
The square identity
\[
 1-mz^2-(1+\kappa)^2z^2(1-z^2)
 =[1-(1+\kappa)z^2]^2
\]
implies $mhz/D\le\lambda$.
If $(2n+1)K\le p\le(2n+2)K$, then $sc\le0$ and
$q_p\ge\mathcal Q>\lambda$. Put $S=|s|$, $C=|c|$.
When $dS\ge Cq_p$,
\[
 q_pd-\lambda S|dS-Cq_p|
 =d(q_p-\lambda S^2)+\lambda SCq_p>0;
\]
when $dS<Cq_p$, the same difference is
$q_p(d-\lambda SC)+\lambda dS^2>0$, because $d\ge C$ and
$\lambda S<1$. These inequalities imply $|U|>w|V|$, a contradiction.
The endpoints $p=2jK$ have zero heading and are excluded.

For $n=0$ we now have $p<K<p_1$, hence
$t<2kK<2kp_1=t_{\rm cut}$ \cite[(5.17)]{MS}\cite[(2.2)]{SS}.
Strict pre-cut uniqueness and regularity follow from the known synthesis
and conjugate-time bound \cite[Theorem 3.2]{SS}\cite{SC}.
If the target admits a rotating source, Corollary~\ref{cor:zero-winding-regimes}
supplies a rotating cell-zero minimizer. It is the unique cell-zero source.
\end{proof}

\begin{corollary}[An exact uniqueness sector]\label{cor:exterior-sector-count}
For $X\ne0$, $W=\pm3/4$ and $|U|\le|V|/2$, the zero-winding
normal fiber contains exactly one source. It is regular, strictly pre-cut
and uniquely minimizing, and
\[
 \nu_2(q)=\ind_{\{|U|>\mathcal F(V,3/4)\}},\qquad N_0(q)=1.
\]
In particular, this fixed-heading sector contains no rotating conjugate
value. The stated angular boundary is sufficient; optimality is not asserted.
\end{corollary}
\begin{proof}
There is no rotating source when $U=0$. Otherwise $0<b\le1/2<w=3/4$,
and Lemma~\ref{lem:exterior-periods} gives $q_p<2$.
We show that $n\ge1$ is impossible. With the notation of its proof,
put $r=ms^2\ge9/25$. Then
\[
 D\ge\tfrac54\sqrt{1-m},\qquad
 bk|c|z\le\tfrac12\sqrt{m-9/25},\qquad
 (bkzd|s|)^2=b^2\left(\tfrac{25r-9}{16}\right)(1-r)\le\tfrac1{25}.
\]
The last quadratic has maximum $4/25$ at $r=17/25$ before
multiplication by $b^2\le1/4$. Whenever its denominator is positive,
\begin{equation}\label{eq:sector-drift-upper}
 q_p\le
 \frac{1/5}{(5/4)\sqrt{1-m}-(1/2)\sqrt{m-9/25}}.
\end{equation}
For $9/25\le m\le3/5$, that denominator is greater than $1/2$:
use $\sqrt{2/5}>3/5$ and $\sqrt{6/25}<1/2$ at $m=3/5$.
Thus $q_p<2/5$, while $n\ge1$ requires
$q_p\ge2\mathcal Q(m)>3m/2\ge27/50$.
Here $\mathcal Q(m)>\pi m/4>3m/4$ follows from its positive integral.
For $3/5\le m\le4/5$, the denominator is at least
$\sqrt5/4-\sqrt{11}/10>2/9$, using
$\sqrt5>20/9$ and $\sqrt{11}<10/3$.
Hence $q_p<9/10$, whereas $2\mathcal Q(m)>3m/2\ge9/10$.

Finally $\mathcal Q(4/5)>1$. Indeed, retaining five positive terms
in the binomial expansion of its defining integral gives
\[
 \mathcal Q(m)>\frac{\pi m}{4}
 \left(1+\frac{3m}{8}+\frac{15m^2}{64}
             +\frac{175m^3}{1024}+\frac{2205m^4}{16384}\right).
\]
At $m=4/5$, $\pi>157/50$ makes the right side greater than
$2000337/2000000>1$. Thus $m\ge4/5$ would give
$q_p\ge2\mathcal Q(m)>2$, contradicting $q_p<2$.
These three ranges exhaust all admissible moduli.

Every rotating source consequently has $n=0$ and is strictly pre-cut by
Lemma~\ref{lem:exterior-periods}. The exact image theorem supplies it
precisely when $|U|>\mathcal F(V,3/4)$; the other cases and their
regularity follow from Corollary~\ref{cor:zero-winding-regimes}.
\end{proof}

For the ray of Figure~\ref{fig:radial-window}, the corollary gives exactly
one rotating source for $R_-<R<R_+$ and none elsewhere.
The zero-winding source is unique and regular at every positive radius.
In original covector and time coordinates it varies analytically with
$R$, by the inverse function theorem and uniqueness, changing from $C_1$
through $C_3$ to $C_2$ and back through $C_3$ to $C_1$.

The midpoint identity also bounds every positive period away from the
separatrix at a fixed strict-heading target. With $w=|W|$ and the
notation of Lemma~\ref{lem:exterior-periods},
$D^2=1-m+w^2d^2/s^2$ gives
\begin{equation}\label{eq:positive-period-modulus}
 q_p=\frac{|U|d}{D}<\frac{|U|\,|s|}{w}\le\frac{|U|}{w},
 \qquad 2n\mathcal Q(m)\le q_p.
\end{equation}
This inequality does not require exteriority. For $n\ge1$, it implies
$m<\mu_n$, where $\mathcal Q(\mu_n)=|U|/(2nw)$.
The solution $\mu_n\in(0,1)$ exists uniquely because $\mathcal Q$ is
strictly increasing from zero to infinity \cite[Sections 19.4, 19.12]{DLMF}.
Together with $m\ge m_h$, this excludes the period if $\mu_n\le m_h$.
The following pointwise certificate combines this bound with the unique
exterior cell-zero source; it uses interval computation to count the
remaining period-one sources.

\begin{proposition}[A certified exterior three-source fiber]
\label{prop:exterior-three}
Choose the exact source data
\[
 m_0=\frac{9999}{10000},\qquad \alpha_0=\frac{29}{20},\qquad
 p_0=2K(m_0)+F(\alpha_0\mid m_0),
\]
\[
 \cos\eta_0=\frac{3\sqrt{1-m_0\sin^2\alpha_0}}
 {4\sqrt{m_0}\sin\alpha_0},\qquad 0<\eta_0<\frac\pi2,
 \qquad \psi_0=F(\eta_0\mid m_0)-p_0.
\]
Let $q_*$ be the endpoint of $(\psi_0,\sqrt{m_0},2\sqrt{m_0}p_0,+1)$.
Its zero-winding fiber consists of exactly three rotating sources, all
regular: one strictly pre-cut cell-zero minimizer and two post-cut
sources of actual period one. There are no higher-period rotating sources.
In particular, exterior uniqueness does not hold in every direction.
\end{proposition}
\begin{proof}[Computer-assisted proof]
Direct outward evaluation at 256 bits gives $U,V>0$, $W=3/4$ and
\[
 \frac7{10}<\frac UV<\frac{18}{25}<\frac34,\qquad
 4\sqrt{m_0}K(m_0)<2\sqrt{m_0}p_0<6\sqrt{m_0}K(m_0).
\]
Thus $q_*$ is exterior and has a period-one source.
Lemma~\ref{lem:exterior-periods} supplies exactly one regular, strictly
pre-cut cell-zero source. All other sources have positive actual period.
The exact target is defined by the source, not by its approximation
\[
 (-7.6417489246,-23.6582865360,2\arctan(3/4)).
\]

For the scalar inverse of Proposition~\ref{prop:inverse},
$\sin\theta=24/25$ and $\cos\theta=7/25$.
The two roots of $P-2C$, with
$P=C^2+(7z/25-24/25)^2-z^2$, lie below $-7/24$, the
heading lower bound. The roots of $P+2C$ lie inside
\[
 I_{\rm out}=[-0.281164,-0.280274].
\]
The quadratic factorization \eqref{eq:factor} therefore gives one
admissible component, between the latter two roots.
On $I_{\rm out}$, $C<0$, $xz-y>0$ and $\cos\theta+z\sin\theta>0$.
There is no component at infinity, and $m>0$ since $z\ne0$.
The negative rotation sign is excluded by $U>0$ and
$U=\varepsilon q_pD/d$. The $B=0$ seam is excluded by $V\ne0$.
Equivalently, $C<0$ on the entire required heading half-plane,
since $y<0$ and $C(-7/24)<0$.

Thirty-two adjacent rational intervals prove $m>0.9998$ on all
$I_{\rm out}$. Since the target satisfies $U/W<12$ and
\[
 4\mathcal Q(0.9998)>12,\qquad 2\mathcal Q(0.99999)>12,
\]
\eqref{eq:positive-period-modulus} excludes $n\ge2$ everywhere,
and excludes $n\ge1$ wherever $m\ge0.99999$.
Sixteen rational intervals at each end prove that latter inequality on
\[
 [-0.281164,-0.28115]\quad\hbox{and}\quad[-0.28029,-0.280274].
\]
It remains to count $N(z)=1$ on $I=[-0.28115,-0.28029]$, whose
endpoints lie strictly inside the admissible component.

An adjacent 256-bit interval cover of $I$ has $1282$ leaves.
Each either excludes zero from $N-1$, or excludes zero from $N_z$
and proves strict endpoint signs. The latter test counts one root
exactly when the endpoint signs differ. The only two root intervals are
\[
 \left[-\frac{14378669}{51200000},-\frac{7189313}{25600000}\right],
 \qquad
 \left[-\frac{898519}{3200000},-\frac{3594033}{12800000}\right].
\]
The derivative is strictly positive on the first and strictly negative
on the second. Every shared endpoint is excluded as a root by the
same strict tests; no interval remains unresolved.
Proposition~\ref{prop:fold} proves both sources regular.
Their actual period is one, so $t>4kK>2kp_1=t_{\rm cut}$.
Adding the unique cell-zero source gives three.
Corollary~\ref{cor:zero-winding-regimes} excludes other zero-winding strata.

The exact rational cover, target enclosures and signs are in
\nolinkurl{data/e109_exterior_count.json}, produced by
\nolinkurl{research/e109_exterior_count.py}.
A full 384-bit replay using unfactored geometry with mean-value
enclosures reproduces the count in $1328$ nodes and rechecks all
$64$ modulus intervals. It shares the Arb library and elliptic derivative
primitives. A separate exact rational Bernstein calculation checks the
domain and modulus polynomial inequalities without that library.
The replay and rational audit are recorded in
\nolinkurl{data/e109_replay.json} and
\nolinkurl{data/e110_rational_audit.json}.
\end{proof}


\subsection{Sharp angular thresholds for positive periods}
\label{sec:exterior-thresholds}

The previous explicit sector can be sharpened by a variational
characterization. The thresholds below concern existence at some radius
on a fixed ray, rather than the count at a prescribed radius.
Fix $w>0$, put $m_h=w^2/(1+w^2)$, and define
\[
 \mathcal D_w=\{(m,\alpha):m_h<m<1,\quad
       \arcsin\sqrt{m_h/m}<\alpha<\pi/2\}.
\]
For $(m,\alpha)\in\mathcal D_w$, set
\[
 \begin{gathered}
 S=\sin\alpha,\quad C=\cos\alpha,\quad d=\sqrt{1-mS^2},
 \quad k=\sqrt m,\\
 h=\frac{wd}{kS},\qquad z=\sqrt{1-h^2},\qquad D=\sqrt{1-mz^2},\\
 q_n=2n\mathcal Q(m)+F(\alpha\mid m)-E(\alpha\mid m),\qquad
 B_n(m,\alpha;w)=\frac{Dq_n}{kz(Cq_n+dS)},\quad n\ge1.
 \end{gathered}
\]
All factors are positive. The letter $B_n$ here denotes a slope function,
not the seam coefficient $B$ of the scalar inverse.

\begin{theorem}[Sharp exterior period thresholds]
\label{thm:exterior-thresholds}
For each $w>0$ and integer $n\ge1$, the minimum
\begin{equation}\label{eq:angular-threshold}
 \beta_n(w)=\min_{(m,\alpha)\in\mathcal D_w}B_n(m,\alpha;w)
\end{equation}
is attained in $\mathcal D_w$. With $\sigma(w)$ from
Proposition~\ref{prop:radial-existence},
\begin{equation}\label{eq:angular-hierarchy}
 \sigma(w)<\beta_1(w)<\beta_2(w)<\cdots<w,\qquad
 0<w-\beta_n(w)<\frac{w}{1+2n\mathcal Q(m_h)}.
\end{equation}
In particular, $\beta_n(w)\to w$.

Fix a strict heading with $|W|=w$ and an exterior spatial ray with
$b=|U|/|V|<w$. A rotating source of actual period $n$ reaches some
positive radius on that ray if and only if $b\ge\beta_n(w)$.
At $b=\beta_n(w)$ every such period-$n$ source is post-cut and conjugate.
The zero-winding fiber consists of exactly one source, which is regular
and strictly pre-cut, at every positive radius on the ray if and only if
$b<\beta_1(w)$.
Different periods need not occur at the same radius.
\end{theorem}
\begin{proof}
Lemma~\ref{lem:exterior-periods} puts every exterior period-$n$ source
at $p=2nK+F(\alpha\mid m)$ in the stated domain. Its absolute endpoint
coordinates are
\[
 U_n=\frac{q_nD}{d},\qquad V_n=\frac{kz(Cq_n+dS)}d,
 \qquad |W|=w,
\]
so its slope is $B_n$. Conversely every domain point gives a source.
For prescribed signs $s_U,s_V,s_W\in\{-1,1\}$, choose
\[
 \varepsilon=s_U,\quad
 \cos\eta=(-1)^{n+1}s_Wh,\quad
 \sin\eta=(-1)^{n+1}\varepsilon s_Vz.
\]
With a consistent real amplitude lift, $\psi=F(\eta\mid m)-p$ and
$t=2kp$ then give $(U,V,W)=(s_UU_n,s_VV_n,s_Ww)$.
The physical radius is $2\sqrt{U_n^2+V_n^2}/\sqrt{1+w^2}>0$.
For $j>i\ge1$, direct subtraction gives
\begin{equation}\label{eq:angular-period-comparison}
 B_j-B_i=
 \frac{DdS(q_j-q_i)}{kz(Cq_i+dS)(Cq_j+dS)}>0,
 \qquad q_j-q_i=2(j-i)\mathcal Q.
\end{equation}

We first show that every sublevel $B_n\le B_*<w$ is compact inside
$\mathcal D_w$. The exterior drift bound implies
\[
 2n\mathcal Q(m)\le q_n<\frac{B_*}{w-B_*}.
\]
Since $\mathcal Q(m)\to\infty$ as $m\to1$, this bounds $m$ above by
$m_{\max}<1$. The numerator $Dq_n$ is bounded below by
$2n\sqrt{1-m_{\max}}\mathcal Q(m_h)>0$, whereas the denominator
is at most $z(q_n+1)$. Thus the sublevel stays away from $z=0$,
equivalently $mS^2=m_h$, and also from $m=m_h$.
At the remaining boundary $\alpha=\pi/2$,
\[
 B_n(m,\pi/2;w)=\frac{(2n+1)\mathcal Q(m)}{\sqrt{m-m_h}}>w.
\]
Indeed, $\mathcal Q>1-\kappa$, $\kappa=\sqrt{1-m}$, and
\[
 (1-\kappa)^2-w^2(m-m_h)
 =\left(\sqrt{1+w^2}\,\kappa-\frac1{\sqrt{1+w^2}}\right)^2\ge0.
\]
Continuity on the remaining compact modulus interval excludes this
boundary too, including joint approaches to the original corners.

Such an exterior sublevel is nonempty for every fixed $n,w$.
Let $e\downarrow0$ and choose $m=1-e^2$, $C=\sqrt e$,
$S=\sqrt{1-e}$. These points eventually lie in $\mathcal D_w$.
Writing $a=dS/C$ and $G=D/(kzC)$, rational algebra yields
\[
 a^2=1-2e^2+e^3,\qquad
 \frac{G^2}{w^2}=1+(w+w^{-1})^2e+O_w(e^2).
\]
Moreover $2n\mathcal Q<q_n<(2n+1)\mathcal Q$ and
$\mathcal Q(1-e^2)=\log(4/e)-1+o(1)$
\cite[Section 19.12]{DLMF}. Thus $q_n\to\infty$, $eq_n\to0$, and
\[
 q_n\left(\frac{B_n}w-1\right)
 =\frac{q_n(G/w-1)-a}{1+a/q_n}\longrightarrow-1.
\]
Hence $B_n<w$ eventually. Compactness of its sublevel proves that
the positive global minimum is attained in the interior and is less
than $w$. An attained minimum is an actual exterior rotating target,
so Proposition~\ref{prop:radial-existence} gives $\beta_n>\sigma(w)$.
Evaluate \eqref{eq:angular-period-comparison} at a minimizer of
$B_{n+1}$ to get $\beta_n<\beta_{n+1}$. At a minimizer of $B_n$,
$2n\mathcal Q(m_h)\le q_n<\beta_n/(w-\beta_n)$ gives
the final bound in \eqref{eq:angular-hierarchy}.

The domain is connected, and $B_n$ is unbounded above as $z\to0$
at fixed $m$. Its image is therefore $[\beta_n,\infty)$.
Source reconstruction with arbitrary prescribed signs proves the
claimed support on each exterior ray.

At an interior minimum, both derivatives of $B_n$ at fixed $W$ vanish.
The source coordinates $(m,\alpha,\eta)$ are nonsingular since
$k_m>0$, $p_\alpha=1/d>0$ and $\psi_\eta=1/D>0$.
Replacing $\eta$ by $W$ is legitimate because
$W_\eta=k\sn p\sin\eta/d\ne0$.
At fixed $W$, the nonzero row $(1/V,-U/V^2)$ annihilates the
derivative of $(U,V)$, so the exponential map is singular.
Every source on the threshold ray with actual period $n$ is a
minimizer of $B_n$, hence conjugate. It is post-cut because
$p>2nK\ge2K>p_1$.

For $b<\beta_1$ there is no positive-period rotating source.
Lemma~\ref{lem:exterior-periods} and
Corollary~\ref{cor:zero-winding-regimes} give exactly one regular
pre-cut zero-winding source at every radius, including $b=0$.
For $\beta_1\le b<w$, some radius has a positive-period source
as well as the distinct cell-zero minimizer. This proves the sharp
strict inequality for all-radius uniqueness.
\end{proof}

The theorem does not classify the corank or fold/cusp type at a
threshold, nor assert uniqueness of the minimizing parameters in
\eqref{eq:angular-threshold}. It concerns the zero-winding fiber;
the full normal fiber still contains its nonzero-winding sources.
Corollary~\ref{cor:exterior-sector-count}, together with attainment,
already gives $\beta_1(3/4)>1/2$. A complete interval cover gives
the following much narrower enclosure.

\begin{corollary}[A certified angular threshold]
\label{cor:angular-threshold-certificate}
The first exterior threshold at $w=3/4$ satisfies
\begin{equation}\label{eq:angular-threshold-certificate}
 0.70555041348314<\beta_1(3/4)<0.705550413483143.
\end{equation}
\end{corollary}
\begin{proof}[Computer-assisted proof]
An upper witness is given by the exact source parameters
\[
 m_* =\frac{999956771591962}{10^{15}},\qquad
 \alpha_* =\frac{1469648769750615}{10^{15}},\qquad w=\frac34.
\]
Direct outward evaluation at 384 bits proves
$(m_*,\alpha_*)\in\mathcal D_{3/4}$ and
$B_1(m_*,\alpha_*;3/4)<0.705550413483143$.
These rational parameters define a source, not an asserted exact
minimizer. The resulting upper bound is also less than $b_*=0.705551$.

To prove the lower bound, parametrize the full domain by
\[
 \begin{gathered}
 m=1-\tfrac{16}{25}e^{-\xi},\quad
 \Lambda=\sqrt{m-\tfrac9{25}},\quad v=\tfrac\pi2\zeta,\\
 r=\tfrac9{25}+\Lambda^2\sin^2v,\quad
 \alpha=\operatorname{atan2}(\sqrt r,\Lambda\cos v),
 \qquad \xi>0,\quad0<\zeta<1.
 \end{gathered}
\]
Here $r=m\sin^2\alpha$; $\xi,\zeta$ are source coordinates,
unrelated to the target coordinates $x,y$.
Let $m_c=1-(16/25)e^{-16}$. Outward complete-integral evaluation gives
\[
 2\mathcal Q(m_c)>\frac{b_*}{3/4-b_*}.
\]
Lemma~\ref{lem:exterior-periods} therefore excludes $\xi\ge16$
for every source with $B_1\le b_*$, in particular for every minimizer.

On $[0,16]\times[0,1]$, put $q=q_1$ and
\[
 \begin{gathered}
 D_c=\sqrt{1-m+\tfrac9{16}m(1/r-1)},\qquad
 A_c=\frac{5\Lambda^2\sin v\cos v}{4\sqrt r},\\
 C_c=\frac54\Lambda\sin v\sqrt{1-r},\qquad
 \mathcal R=q(D_c-\underline b A_c)-\underline b C_c,
 \quad \underline b=0.70555041348314.
 \end{gathered}
\]
In the open source domain, $B_1=qD_c/(A_cq+C_c)$ with positive
denominator, so $\mathcal R>0$ proves $B_1>\underline b$.
The residual has continuous boundary values on the closed rectangle.
A complete 192-bit cover has $7825$ rational boxes, each with a
strictly positive residual enclosure. Each box uses either direct
interval evaluation or a centered mean-value bound with derivatives
enclosed on the whole box. Nonfinite derivatives disable that tightening;
such a box must pass its direct bound or be subdivided.

The incomplete-integral corner bounds use only strictly positive
amplitudes. Their signs follow from the defining integrals: $q$ increases
with $m$ and $\alpha$, whereas $E(\alpha\mid m)$ decreases with $m$
and increases with $\alpha$. A separate implementation using explicit,
simplified derivatives rechecks all $7825$ boxes at 256 bits in
$7849$ nodes. Exact rational partition reconstruction and an independent
sweep both prove full coverage without gaps or overlapping interiors.
No box remains unresolved.  The cover was assembled in two runs, $7176$
leaves inherited from a first run and $649$ nodes added, and the replay
re-evaluates every leaf.  The two numerical implementations share Arb
special functions; their agreement is not an independent-library review.

The full cover and alternate replay are
\nolinkurl{data/e113_threshold_cover.json} and
\nolinkurl{data/e113_threshold_replay.json}.
The separate sweep, symbolic derivative audit and rational upper witness
are recorded in \nolinkurl{data/e114_certificate_audit.json}.
Together with attainment from Theorem~\ref{thm:exterior-thresholds},
the strict residual bounds prove the strict lower endpoint in
\eqref{eq:angular-threshold-certificate}.
\end{proof}

This encloses the global minimum to an interval of width $3\cdot10^{-15}$.
A separate localization and second-derivative certificate identifies its
unique attaining source and local type in this particular case.

\begin{theorem}[A unique threshold fold and its radial count window]
\label{thm:threshold-fold-window}
Put $w=3/4$ and $\beta=\beta_1(w)$. The function $B_1$ has exactly
one global minimizer in $\mathcal D_w$. In the source coordinates of
Corollary~\ref{cor:angular-threshold-certificate}, it belongs to
\begin{equation}\label{eq:threshold-localization}
 \mathcal C=[9.6027,9.60275]\times[0.91943,0.919431].
\end{equation}
For each prescribed set of signs of $U,V,W$, its reconstructed source
is an ordinary fold of the full exponential map. Its target radius
$R_*$ satisfies
\begin{equation}\label{eq:threshold-radius}
 26.8980<R_*<26.8989.
\end{equation}
The certificate encloses $R_*$ in $(26.89806,26.89887)$; the separate
audit of the same quantity gives only $(26.8979,26.8991)$, so the
displayed digits rest on one implementation.
On that signed ray at slope $b=\beta$, the zero-winding count $N_0$
is two at $R_*$ and one at every other positive radius.

There exists $\delta>0$ such that, for each $\beta<b<\beta+\delta$,
two radii $R_-(b)<R_+(b)$ give the complete count
\begin{equation}\label{eq:threshold-window-count}
 N_0(R)=\begin{cases}
 1,&R\notin[R_-(b),R_+(b)],\\
 2,&R=R_-(b)\ \text{or}\ R=R_+(b),\\
 3,&R_-(b)<R<R_+(b).
 \end{cases}
\end{equation}
The additional sources have actual period one and are post-cut:
one fold at either endpoint, two regular sources inside. The remaining
source is regular and strictly pre-cut. Both radii tend to $R_*$ as
$b\downarrow\beta$. A certified choice of $\delta$ is supplied in
Corollary~\ref{cor:effective-fold-window}.
\end{theorem}
\begin{proof}[Computer-assisted proof]
Let $b_U=0.705550413483143$. At this upper comparison level, repeat
the residual test in Corollary~\ref{cor:angular-threshold-certificate}.
Cut its verified rectangular partition at the four rational boundaries
of $\mathcal C$. Retain pieces inside $\mathcal C$ and prove strict
positivity of $q(D_c-b_U A_c)-b_U C_c$ on every other piece, subdividing
when necessary. The resulting cover has $7465$ excluded boxes and
$1343$ retained boxes, with none unresolved. A 384-bit alternate
implementation rechecks all $7465$ exclusions. Exact rational area
and pairwise overlap checks within each of the $7825$ original boxes
independently verify the partition; missing and duplicated pieces fail
these checks. The analytic tail exclusion is unchanged and its strict
margin is reevaluated. Since $\beta<b_U$, every global minimizer lies
in $\mathcal C$. Retained pieces carry no positive-residual claim.

Write $H=D^2_{\xi,\zeta}B_1$ and
$\mathbf k=(V_\zeta,-V_\xi)$. On the whole convex rectangle
$\mathcal C$, two interval implementations give
\[
 H_{11}>0.0033,\qquad \det H>0.0089,\qquad
 V_\xi>1.0231,\qquad \mathbf k^TH\mathbf k>9.5.
\]
One implementation differentiates the simplified drift gradients in
the transformed chart. The separate audit differentiates the original
$(m,\alpha)$ formulas symbolically and pulls back their Hessian,
including the second derivatives of the coordinate change. It uses
\[
 q_m=\frac{2E+E(\alpha\mid m)-\sin\alpha\cos\alpha/d}{2(1-m)},
 \qquad q_\alpha=\frac{m\sin^2\alpha}{d},
 \qquad d=\sqrt{1-m\sin^2\alpha},
\]
with the standard elliptic derivative identities
\cite[Section 19.4]{DLMF}. The implementations share Arb special
functions; this is not an independent-library review.
The positive principal minors make $B_1$ strictly convex on
$\mathcal C$. Attainment from Theorem~\ref{thm:exterior-thresholds}
therefore gives exactly one global minimizer. It is interior to the
original source domain, so its full $(\xi,\zeta)$ gradient vanishes,
whether or not it lies on the artificial boundary of $\mathcal C$.
Outward evaluation of $R=2\sqrt{U^2+V^2}/\sqrt{1+w^2}$ on
$\mathcal C$ gives \eqref{eq:threshold-radius}.

The source chart $(\xi,\zeta,W)$ is nonsingular by the chart argument
in Theorem~\ref{thm:exterior-thresholds}; explicitly,
$m_\xi=1-m>0$ and
$\alpha_\zeta=\pi\Lambda\sin v/(2\sqrt r)>0$.
Target coordinates $(B=U/V,V,W)$ are nonsingular since $V\ne0$.
At the minimum, $\nabla B=0$ and $V_\xi\ne0$, so the full map
has rank two and kernel $(\mathbf k,0)$. Its transverse quadratic
derivative is $\mathbf k^TH\mathbf k>0$.
To see the full-map fold directly, choose source coordinates $(s,V,W)$.
The map is $(f(s,V,W),V,W)$, with $f_s=0$, $f_{ss}>0$ at the point.
The implicit function theorem gives a critical section $s_c(V,W)$.
Taylor's integral formula writes $f-f(s_c,V,W)$ as $(s-s_c)^2$
times a smooth positive factor. Rescaling the first source coordinate
and subtracting $f(s_c,V,W)$ in the target gives $(z^2,V,W)$.
This proves the ordinary fold type in three dimensions.
The signed reconstruction transfers the result to every sign choice.

The strict inequality $\beta_2>\beta$ excludes all higher periods
on the threshold ray. Every period-one source there minimizes $B_1$,
so uniqueness gives one such source at one radius $R_*$.
Lemma~\ref{lem:exterior-periods} supplies the distinct pre-cut
cell-zero source, proving the count at $b=\beta$.

Choose $\delta$ with $\beta+\delta<\min(\beta_2,b_U)$.
Outside any neighborhood of the unique minimum, the compact
complement inside $\mathcal C$ has minimum strictly greater than
$\beta$. Decreasing $\delta$ thus places every period-one source
in one fold neighborhood. Choose a smaller product neighborhood in
$(z,V)$ so that all solutions considered below remain in its interior.
At fixed $W$, the fold representation retaining the original $B,V$
coordinates is
\[
 B=c(V)+z^2,\qquad c(V_*)=\beta,\qquad c'(V_*)=0,
 \qquad c''(V_*)=\frac{1}{\nabla V^T H^{-1}\nabla V}
 =\frac{\det H}{\mathbf k^TH\mathbf k}>0.
\]
The curvature formula follows by differentiating the constrained
minimum of $B$ at fixed $V$. Hence $c(V)=b$ has exactly two nearby
solutions $V_-(b)<V_+(b)$: between them there are two period-one
sources, at either endpoint one fold, and outside none. Compactness
excludes remote period-one sources, and $b<\beta_2$ excludes higher
periods. Corollary~\ref{cor:zero-winding-regimes} and
Lemma~\ref{lem:exterior-periods} supply exactly one remaining
pre-cut source at every radius, in its appropriate stratum.
Finally $R=2\sqrt{1+b^2}|V|/\sqrt{1+w^2}$ proves
\eqref{eq:threshold-window-count} and convergence of the endpoints.

The localization and local bounds are recorded in
\nolinkurl{data/e115_threshold_localize.json} and
\nolinkurl{data/e115_threshold_local.json}; the separate Hessian
and full exclusion audits are
\nolinkurl{data/e116_threshold_audit.json} and
\nolinkurl{data/e116_localization_replay.json}.
\end{proof}

The same proof gives a square-root opening law. With
$\kappa=c''(V_*)>0$,
\[
 R_+(b)-R_-(b)=2C_R\sqrt{b-\beta}+O(b-\beta),\qquad
 C_R=\frac{2\sqrt{1+\beta^2}}{\sqrt{1+w^2}}\sqrt{\frac2\kappa}.
\]
These statements concern distinct zero-winding sources; a fold counts
once. The full normal fiber remains countably infinite.
Theorem~\ref{thm:window-time-order} below separately orders the post-cut
travel times over the explicit slope range. The following certificate quantifies
a valid slope range for the radial count window.

\begin{corollary}[An effective first fold window]
\label{cor:effective-fold-window}
Put $w=3/4$, $\beta=\beta_1(w)$ and $b_c=0.705551$.
For every prescribed signed exterior ray with
\begin{equation}\label{eq:effective-fold-slopes}
 \beta<b\le b_c,
\end{equation}
the complete count \eqref{eq:threshold-window-count} and the source
types in Theorem~\ref{thm:threshold-fold-window} hold. In particular,
one may take $\delta=0.000000586516857$ in that theorem.
At the exact slope $b=b_c$, the fold radii satisfy
\begin{equation}\label{eq:cap-fold-radii}
 \begin{split}
 26.8301460&<R_-(b_c)<26.8301463,\\
 26.9671811&<R_+(b_c)<26.9671814.
 \end{split}
\end{equation}
This is an interval in the spatial slope $b=|U|/|V|$ at fixed
heading parameter $w$, not an interval of heading angles.
\end{corollary}
\begin{proof}[Computer-assisted proof]
Work first with $U,V,W>0$. In the same source coordinates as above,
set
\[
 \mathcal C_b=[9.55,9.66]\times[0.918,0.921].
\]
The residual test at $b_c$ gives a complete partition with $3111$
excluded boxes and $5167$ retained boxes in $\mathcal C_b$; the
$\xi\ge16$ tail is excluded by the drift inequality. An alternate
384-bit replay checks every excluded box in $3203$ nodes. Exact
parentwise area and overlap checks verify the partition. Separate
edge subdivisions verify that the excluded closed boxes cover all
four sides of $\mathcal C_b$. Thus $B_1>b_c$ on its boundary and
every source with $B_1\le b_c$ lies in its interior. This last
boundary assertion is needed below, not just localization to a
closed rectangle.

Write $B=B_1$, $H=D^2_{\xi,\zeta}B$, and
\[
 \Delta=B_\xi V_\zeta-B_\zeta V_\xi,\qquad
 \mathbf k=(V_\zeta,-V_\xi).
\]
A $2048$-leaf interval cover of $\mathcal C_b$ proves
\[
 H\text{ positive definite},\qquad V_\xi>1,\quad V_\zeta>12,
 \qquad \mathbf k(\Delta)>6,\qquad B_2>b_c+0.02.
\]
An independent symbolic calculation in $(m,\alpha)$, followed by
the full Hessian pullback for both $B$ and $V$, replays every local
leaf at 384 bits in $2080$ nodes. The evaluated kernel derivative is
\[
 \begin{split}
 \mathbf k(\Delta)={}&\mathbf k^TH\mathbf k
 +B_\xi(V_{\xi\zeta}k_\xi+V_{\zeta\zeta}k_\zeta)\\
 &-B_\zeta(V_{\xi\xi}k_\xi+V_{\xi\zeta}k_\zeta).
 \end{split}
\]
No stationary-point assumption is used in these bounds. For $B_2$
the evaluator uses $q_2=q_1+2\mathcal Q(m)$ in the original
endpoint formulas. If any $n\ge2$ had $B_n\le b_c$, the strict
period comparison would give $B_1<B_n\le b_c$, placing its
parameters in $\mathcal C_b$, where $B_n\ge B_2>b_c$ is a
contradiction. Hence no higher period contributes in
\eqref{eq:effective-fold-slopes}.

The positive coordinate derivatives of $V$ make every nontrivial
intersection of $V=v$ with $\mathcal C_b$ a connected graph
$\xi=\xi_v(\zeta)$. Indeed, the allowed $\zeta$ values are determined
by $V(\xi_{\rm left},\zeta)\le v\le
V(\xi_{\rm right},\zeta)$, and both boundary functions increase
strictly. Put $g_v(\zeta)=B(\xi_v(\zeta),\zeta)$. Then
\[
 g_v'=-\frac{\Delta}{V_\xi},\qquad
 g_v''=\frac{\mathbf k(\Delta)}{V_\xi^2}>0
 \quad\text{at every stationary point.}
\]
There is at most one stationary point: two negative-to-positive
crossings of $g_v'$ would require an intervening zero that is not
such a crossing. The graph endpoints lie on the rectangle boundary,
where $B>b_c$. Consequently $g_v=b\le b_c$ has zero, one double,
or two simple roots according as its minimum is above, equal to,
or below $b$.

To exclude several disconnected radial windows, use the positive
Hessian. For $\beta<b\le b_c$, the sublevel
$S_b=\{B\le b\}$ is compact and convex with nonempty interior,
strictly inside $\mathcal C_b$. Its image under $V$ is an interval
$[v_-,v_+]$. Since $V$ is a submersion, $V(\operatorname{int}S_b)$
is open and connected; density of the interior and compactness imply
that it equals $(v_-,v_+)$. For every $v$ in this open interval,
$g_v$ has minimum below $b$, hence two roots. At $v_\pm$ the
minimum equals $b$, giving one fold. Outside the closed interval
there is no root. The positivity of $\mathbf k(\Delta)$ gives
ordinary folds of the full map, while the simple roots are regular.
The signed reconstruction, radius conversion and unique remaining
pre-cut source then prove \eqref{eq:threshold-window-count} for
the whole range \eqref{eq:effective-fold-slopes}.

For \eqref{eq:cap-fold-radii}, interval arithmetic isolates two
solutions of $(B-b_c,\Delta)=0$ in disjoint rational source boxes.
For each box $X$ with center $c$, a nonsingular rational matrix $A$
satisfies
\[
 c-AF(c)+(I-A\,DF(X))(X-c)\subset\operatorname{int}X,
 \qquad F=(B-b_c,\Delta).
\]
The weighted infinity norm of $I-A\,DF(X)$ is less than $0.001$.
Thus $z\mapsto z-AF(z)$ is a contraction of $X$ into itself,
proving a unique root. An original-coordinate derivative implementation
independently repeats both inclusions. Evaluating
$R=2\sqrt{1+b_c^2}\,V/\sqrt{1+w^2}$ on their enclosures gives
\eqref{eq:cap-fold-radii}. The preceding global argument proves
that these roots are the two endpoints, not merely two selected
conjugate sources. Finally,
$b_c-0.705550413483143=0.000000586516857$ and the strict
threshold upper bound give the claimed choice of $\delta$.

The original covers and root certificates are
\nolinkurl{data/e117_localize.json},
\nolinkurl{data/e117_local_window.json} and
\nolinkurl{data/e117_window_geometry.json}.
The replays and boundary audit are
\nolinkurl{data/e118_global_replay.json},
\nolinkurl{data/e118_local_replay.json} and
\nolinkurl{data/e118_geometry_audit.json}.
All interval implementations share Arb special functions.
\end{proof}

The cap $b_c$ is a certified extent, not a maximal one. The wider
target of Proposition~\ref{prop:exterior-three} has $b>b_c$ and
retains its separate pointwise certificate. Ordering the post-cut
travel times requires the additional momentum condition proved next.

\begin{theorem}[Time order in the first fold window]
\label{thm:window-time-order}
In the range of Corollary~\ref{cor:effective-fold-window}, fix a radius
$R_-(b)<R<R_+(b)$. Label the two period-one sources by their coordinates
$\zeta_L<\zeta_R$ in the chart of Corollary~\ref{cor:angular-threshold-certificate}, and let $T_0$ be the
travel time of the remaining pre-cut source. Then
\begin{equation}\label{eq:window-time-order}
 T_0<T_L<T_R.
\end{equation}
The gap $T_R-T_L$ tends to zero at both radial endpoints. At fixed
$|V|$, it is strictly increasing with $b$ wherever the pair exists
and $b\le b_c$. No radial monotonicity of the gap is asserted.
\end{theorem}

\begin{proof}
It suffices first to take $U,V,W>0$, with $W=w=3/4$. Write
$S=\sin\alpha$, $C=\cos\alpha$, $d=\sqrt{1-mS^2}$, $k=\sqrt m$,
$h=wd/(kS)$, $z=\sqrt{1-h^2}$ and $D=\sqrt{1-mz^2}$.
For period one,
\[
 \begin{gathered}
 q=2(K-E)+F(\alpha\mid m)-E(\alpha\mid m),\qquad
 T=2k(2K+F(\alpha\mid m)),\\
 U=qD/d,\qquad V=kz(Cq+dS)/d.
 \end{gathered}
\]
Here $z=\sin\eta$ is distinct from the chart coordinate $\zeta$.

The cotangent Gauss lemma and Hamiltonian homogeneity give
$dT=\lambda_f\cdot d(x,y,\theta)$ for a smooth fixed-start family of
unit-speed normal extremals, including nonminimizing and conjugate
sources \cite[Section 8.6]{ABB}. Indeed, scaling the initial covector
to $T\lambda_0$ in the time-one exponential gives
$T\lambda_f\cdot d(x,y,\theta)=dH(T\lambda_0)=T\,dT$.
In the retained convention the conserved spatial momenta are
$p_x=\sn\psi$ and $p_y=-\dn\psi/k$. Jacobi subtraction at
$\psi=F(\eta\mid m)-2K-F(\alpha\mid m)$ gives
\[
 p_x=\frac{-zCd+ShD}{H_0},\qquad
 p_y=-\frac{Dd+mzShC}{kH_0},\qquad H_0=d^2(1+w^2).
\]
Substitute $x=2(wU-V)/(1+w^2)$ and $y=-2(U+wV)/(1+w^2)$.
At fixed $W$ this yields the smooth source identity
\begin{equation}\label{eq:window-first-variation}
 dT=\ell_U\,dU+\ell_V\,dV,\qquad
 \ell_U=\frac{2D}{kd(1+w^2)},\qquad
 \ell_V=\frac{2zC}{d(1+w^2)}.
\end{equation}
No inverse-map determinant is divided out. Direct differentiation of
the displayed clock and endpoint formulas gives the same identity.

The additional interval condition is
\begin{equation}\label{eq:window-momentum-sign}
 \left.\frac{d\ell_U}{d\zeta}\right|_V>0
 \quad\hbox{on }\{B\le b_c\}\cap\mathcal C_b,
 \qquad B=U/V,
\end{equation}
where $\mathcal C_b$ is the rectangle in the preceding proof. To make
its certificate explicit, put $g=1-m$ and $r=m\sin^2\alpha$. Then
\[
 P=\frac{g}{m(1-r)}+\frac{w^2}{r},\qquad
 \ell_U=\frac{2\sqrt P}{1+w^2},\qquad
 \left.\frac{d\ell_U}{d\zeta}\right|_V
 =\frac{V_\xi P_\zeta-V_\zeta P_\xi}
 {(1+w^2)\sqrt P\,V_\xi}.
\]
A complete mixed cover has $117$ boxes with positive $V_\xi$ and
positive momentum derivative, and $38$ boxes with $B>b_c$.
An exact area and disjointness audit verifies coverage; a separate
sweep rejects missing and duplicated-box controls. The latter boxes
cannot contain a point of $\{B\le b_c\}$.
An alternate 384-bit calculation in the original $(m,\alpha)$ coordinates
replays all $155$ boxes in $287$ nodes. It uses
\[
 V_\xi(\ell_U)_\zeta-V_\zeta(\ell_U)_\xi
 =g\alpha_\zeta\bigl(V_m(\ell_U)_\alpha-V_\alpha(\ell_U)_m\bigr)
\]
and independently simplified derivatives.  Both implementations share
Arb special functions, and the first reuses the residual evaluator of
the threshold-cover replay. Positivity on the entire rectangle is neither
used nor true: at $(\xi,\zeta)=(9.66,.918)$ the derivative is negative,
but $B>b_c$. The certificates are
\nolinkurl{data/e119_time_order.json} and
\nolinkurl{data/e120_time_order_replay.json}; exact algebra and signed
reconstruction checks accompany them.

Fix $V=v>0$ and follow the connected graph $\xi=\xi_v(\zeta)$
from Corollary~\ref{cor:effective-fold-window}. Put
$f(\zeta)=B(\xi_v(\zeta),\zeta)$.
Between its two roots $f(\zeta_L)=f(\zeta_R)=b$, that corollary gives
$f<b$. The intervening interval therefore lies in $\{B\le b_c\}$.
Equation~\eqref{eq:window-first-variation} and integration by parts give
\begin{equation}\label{eq:window-time-gap-integral}
 T_R-T_L
 =v\int_{\zeta_L}^{\zeta_R}(b-f(\zeta))
       \left.\frac{d\ell_U}{d\zeta}\right|_V\,d\zeta>0.
\end{equation}
The unique pre-cut source is minimizing, while both period-one sources
are post-cut, so $T_0<T_L$. Signed reconstruction preserves
$(m,\alpha,T)$ and the $\zeta$ labels, proving the assertion for every
prescribed sign choice. At each radial endpoint the two sources
coalesce, and smoothness of $T$ gives a zero limiting gap.
Finally, on the regular branches at fixed $v$,
\[
 \frac{d(T_R-T_L)}{db}=v(\ell_{U,R}-\ell_{U,L})>0
\]
by \eqref{eq:window-momentum-sign}. Fixing $|V|$ as $b$ varies changes
$R=2|V|\sqrt{1+b^2}/\sqrt{1+w^2}$; this is not radial monotonicity.
\end{proof}

The integral also displays the usual local fold scaling. At a fixed
$v$ with constrained minimum $f(\zeta_0)=c(v)<b_c$, put
$a=f''(\zeta_0)>0$ and $\ell=(d\ell_U/d\zeta)|_V(\zeta_0)>0$.
Taylor expansion in \eqref{eq:window-time-gap-integral} gives
\[
 T_R-T_L=\frac{4\sqrt2}{3}\frac{v\ell}{\sqrt a}
          (b-c(v))^{3/2}+O\bigl((b-c(v))^2\bigr).
\]
This is a local asymptotic at fixed $v$, with no uniform remainder
claimed. First variation and the fold mechanism are classical; the
certificate establishes their time-order consequence throughout the
specified $\mathrm{SE}(2)$ window. No time comparison outside that
window or between arbitrary periods follows.






\subsection{Time comparison along a continuous period label}
\label{sec:period-action}

A different comparison applies when sources in two actual periods can
be joined through the following formal family. Fix $w>0$ and use
$S,C,d,k,h,z,D$ as in the preceding proof, on the open domain
$0<m<1$, $0<\alpha<\pi/2$, $h=wd/(kS)<1$. For a real label $\nu>0$, put
\begin{equation}\label{eq:period-family}
 \begin{gathered}
 q_\nu=2\nu(K-E)+F(\alpha\mid m)-E(\alpha\mid m),\qquad
 p_\nu=2\nu K+F(\alpha\mid m),\\
 U_\nu=q_\nu D/d,\qquad V_\nu=kz(Cq_\nu+dS)/d,\qquad
 T_\nu=2k p_\nu.
 \end{gathered}
\end{equation}
At an integer $\nu=n\ge1$, these are the positive-signed endpoint and
clock formulas in the first quarter of actual period $n$. At noninteger
labels they are formal analytic interpolants, not asserted geodesics.
In particular, physical first variation alone does not determine the
$\nu$ derivative.

\begin{proposition}[Period action on a connected branch]
\label{prop:period-action}
On this formal source domain, with $w$ fixed,
\begin{equation}\label{eq:period-action-identity}
 \begin{gathered}
 dT_\nu=\ell_U\,dU_\nu+\ell_V\,dV_\nu+\mathcal A(m)\,d\nu,\\
 \mathcal A(m)=\frac{4[E(m)-(1-m)K(m)]}{\sqrt m},\qquad
 0<\mathcal A(m)<4,\qquad
 \mathcal A'(m)=\frac{2(K-E)}{m^{3/2}}>0,
 \end{gathered}
\end{equation}
where $\ell_U,\ell_V$ are as in \eqref{eq:window-first-variation}.
Consequently, on any $C^1$ branch
$\nu\mapsto(m(\nu),\alpha(\nu))$ with $U_\nu,V_\nu$ constant,
\begin{equation}\label{eq:period-action-gap}
 0<T_{\nu_2}-T_{\nu_1}
   =\int_{\nu_1}^{\nu_2}\mathcal A(m(\nu))\,d\nu
   <4(\nu_2-\nu_1)\qquad(\nu_1<\nu_2).
\end{equation}
Here both labels must lie on that same branch, parametrized by $\nu$.
\end{proposition}
\begin{proof}
The standard elliptic derivatives \cite[Sections 19.2 and 19.4]{DLMF},
with parameter $m=k^2$, give, on writing $J=2\nu E+E(\alpha\mid m)$,
\[
 (q_\nu)_m=\frac{J-SC/d}{2(1-m)},\quad
 (q_\nu)_\alpha=\frac{mS^2}{d},\quad
 (T_\nu)_m=\frac{J-mSC/d}{k(1-m)},\quad
 (T_\nu)_\alpha=\frac{2k}{d}.
\]
Substitution in $U_\nu,V_\nu$ verifies the $m$ and $\alpha$ components
of \eqref{eq:period-action-identity}. In the remaining direction,
$(q_\nu)_\nu=2(K-E)$ and $(T_\nu)_\nu=4kK$. The identity
$D^2+mz^2C^2=d^2(1+w^2)$ gives
\[
 \ell_U(U_\nu)_\nu+\ell_V(V_\nu)_\nu=4(K-E)/k.
\]
The difference is $\mathcal A(m)$, which proves the full differential
identity without treating a noninteger label as a physical extremal.
The defining integrals give
\[
 \mathcal A(m)=4\int_0^{\pi/2}\cos t\,
       \frac{\sqrt m\cos t}{\sqrt{1-m\sin^2t}}\,dt.
\]
The fraction lies strictly between zero and one in the open integration
interval. Thus $0<\mathcal A<4$. Differentiating the complete integrals
gives $(E-(1-m)K)'=K/2$ and the displayed formula for $\mathcal A'$.
Finally, restrict the differential identity to the assumed branch and
integrate. All source denominators are strictly positive on the domain.
\end{proof}

Dominated convergence gives $\mathcal A(m)\to4$ as $m\to1$. Therefore
a family of such branches uniformly approaching the separatrix has
clock increment asymptotic to four times its label increment. This
requires a connected continuation and uniform approach to $m=1$; it
is not a uniform comparison between arbitrary sources in adjacent cells.

\begin{proposition}[Nine rotating sources with ordered periods]
\label{prop:period-action-example}
The target
\begin{equation}\label{eq:period-action-target}
 (U,V,W)=(6,3,3/4),\qquad
 (x,y,\theta)=(48/25,-264/25,2\arctan(3/4))
\end{equation}
has exactly nine zero-winding normal sources, all rotating and regular.
Their actual-period counts are three in period zero, four in period one
and two in period two; no other rotating periods occur. All nine travel
times are distinct. Every source in period zero arrives before every
source in period one, and every source in period one before every source
in period two, with a gap greater than $2$ between adjacent groups.
Exactly one source is pre-cut; the other eight are post-cut.

The two positive-arc sources in periods one and two having the smaller
auxiliary parameter $a$ in the proof below are also the integer endpoints
of one branch in Proposition~\ref{prop:period-action} over $1\le\nu\le2$.
Their clocks satisfy
\begin{equation}\label{eq:period-action-example-gap}
 3.68729216<T_2-T_1<3.68729218.
\end{equation}
These time comparisons concern this target's rotating sources.
Its nonzero-winding sources are the countably infinite family of
Proposition~\ref{prop:global-winding}.
\end{proposition}
\begin{proof}[Computer-assisted proof]
We first exhaust the possible rotating sources analytically. Put
$p=t/(2k)$, $n=\lfloor p/(2K)\rfloor$ and
$\alpha=\operatorname{am}(p-2nK)\in(0,\pi)$.
A boundary $p=2nK$ would have $W=0$, so it is excluded.
Write $S=\sin\alpha$, $C=\cos\alpha$, $\sigma=S^2$,
$d^2=1-m\sigma$, $\eta=\operatorname{am}(\psi+p)$ and
$D^2=1-m\sin^2\eta$. The sign $s_2$ in
\eqref{eq:rotating-midpoint-chart} equals $+1$, because $U>0$ and
$q=p-\epsilon(p)>0$. Introduce the signed auxiliary scalar
\[
 a=-(-1)^n kS\sin\eta;
\]
it is distinct from the endpoint angle in \eqref{eq:ab}. Then
$0<|a|<1$: $a=0$ would give $V=0$. The midpoint equations become
\[
 m\sigma=\frac{a^2+w^2}{1+w^2},\qquad
 V=a\left(1+\frac{UC}{\sqrt{\sigma-a^2}}\right),\qquad w=3/4.
\]
Since $V=3>|a|$, $\operatorname{sgn}C=\operatorname{sgn}a$.
In particular, $C=0$ cannot reach this target. Solving the second
identity and restoring this cosine sign gives
\begin{equation}\label{eq:nine-source-arcs}
 \begin{gathered}
 \sigma(a)=\frac{a^2[36+(3-a)^2]}{36a^2+(3-a)^2},\qquad
 m(a)=\frac{(16a^2+9)(37a^2-6a+9)}{25a^2(a^2-6a+45)},\\
 q_*(a)=\frac45\sqrt{36+(3-a)^2},\qquad
 a\in(1/3,1)\ \text{or}\ a\in(-1,-3/7).
 \end{gathered}
\end{equation}
The two intervals are exact. Indeed,
\[
 1-m=\frac{(1-a^2)[36a^2-w^2(3-a)^2]}
 {(1+w^2)a^2[36+(3-a)^2]},\qquad
 36a^2-w^2(3-a)^2=\frac{27}{16}(3a-1)(7a+3).
\]
The displayed domains ensure $0<\sigma<1$ and $\sigma>a^2$.
Conversely these formulas reconstruct all endpoint equations with
positive radicals and the forced cosine sign. Thus squaring has
introduced no extra arm or lost sign choice.

Set $\alpha_0=\arcsin\sqrt{\sigma(a)}\in(0,\pi/2)$,
$Q=K-E$ and $B=F(\alpha_0\mid m)-E(\alpha_0\mid m)$.
On the positive arc $\alpha=\alpha_0$; on the negative arc
$\alpha=\pi-\alpha_0$. All sources are therefore exactly the
nonnegative integer intersections of
\begin{equation}\label{eq:nine-source-scalars}
 \begin{aligned}
 f_n^+(a)&=q_*(a)-2nQ-B=0 &&(1/3<a<1),\\
 f_n^-(a)&=q_*(a)-2(n+1)Q+B=0 &&(-1<a<-3/7).
 \end{aligned}
\end{equation}
Each intersection reconstructs one source modulo the $4K$ phase circle,
with actual period $n$. The omitted quarter boundaries have already
been excluded, so these are exhaustive scalar charts for this target.

For the positive arc, exact Bernstein coefficients prove $m>.73$
throughout $[1/3,1]$, and $m>.995$ on $[1/3,.334]$ and $[.999,1]$.
There $q_*\le8\sqrt{97}/15$, and outward complete-integral evaluation gives
\[
 6Q(.73)>8\sqrt{97}/15,\qquad 2Q(.995)>8\sqrt{97}/15.
\]
Since $B>0$ and $Q$ increases with $m$, the first inequality excludes
$n\ge3$, and the second excludes every $n\ge1$ on both tails.
On $[.334,.999]$, the two problems $n=1,2$ have a complete cover of
$1062$ sign-excluded intervals and four retained root intervals.
Each retained interval has opposite endpoint signs and $0\notin(f_n^+)'$,
proving one root. A separate 448-bit implementation, using rational source
derivatives and explicit elliptic partials, replays all exclusions in
$1274$ nodes and recalculates the root derivative and clock enclosures.
A separately included $n=0$ root is the only one on this arc: every such
root has $p=F(\alpha_0\mid m)<K<p_1$, hence is strictly pre-cut, and
classical optimal synthesis implies uniqueness at a fixed target
\cite[Theorems 3.1--3.2]{SS}.

For the negative arc, exact polynomial bounds give $m>.86$ globally
and $m>.99999$ on $[-1,-.99999]$ and $[-.428572,-3/7]$.
Here $q_*\le8\sqrt{13}/5$ and $0<B<Q$, so the inequalities
\[
 5Q(.86)>8\sqrt{13}/5,\qquad Q(.99999)>8\sqrt{13}/5
\]
exclude every $n\ge2$ globally and every $n\ge0$ on the two tails.
The compact remainder has $1953$ excluded intervals and four unique-root
intervals across $n=0,1$. An alternate implementation uses generic
forward derivatives and the original residual $U_n(a)-6$, rather than
$f_n^-$, to replay all exclusions in $11841$ nodes and verify all four
root inclusions. Exact partition audits cover every interval boundary
and reject missing and duplicated intervals. Both implementations share
Arb special functions; these are separate internal replays, not external
review. The certificates are
\nolinkurl{data/e123_first_quarter.json},
\nolinkurl{data/e124_all_quarters.json} and
\nolinkurl{data/e124_remaining_replay.json}.

Simple scalar roots give full exponential-map regularity. For each sign
of $a$, the source change $(a,\sigma)\mapsto(m,\alpha)$ has determinant
\[
 \frac{|a|}{(1+w^2)\sigma^{3/2}\sqrt{1-\sigma}}>0.
\]
Replace the endpoint residuals by
$G_1=q_n-Ud/D$ and
$G_2=a+UaC/\sqrt{\sigma-a^2}-V$; the residual transformation has
determinant $d/D>0$. On the chosen cosine branch,
\[
 (G_2)_\sigma=-\frac{U|a|(1-a^2)}
 {2\sqrt{1-\sigma}(\sigma-a^2)^{3/2}}<0.
\]
Along $G_2=0$, one has $G_1=-f_n^\pm$, so its source determinant is
$-(G_2)_\sigma(f_n^\pm)'\ne0$. The full source change $(m,\alpha,w)\mapsto(m,p,\eta)$ has
determinant $-1/a\ne0$; the standard rotating and target charts are
also nonsingular. These facts give full rank. Reconstructing all nine root-box
centers and integrating the original Hamiltonian system independently
corroborates the target, actual periods and zero lifted heading winding;
these finite integrations do not replace the interval count.

The outward clock enclosures are listed below. Within each arc and period,
``left'' and ``right'' mean smaller and larger $a$.
\begin{center}\small
\begin{tabular}{cccrr}
\toprule
period & arc & branch & lower clock bound & upper clock bound\\
\midrule
0 & positive & sole  & 12.119269 & 12.119271\\
0 & negative & left  & 13.537705 & 13.537707\\
0 & negative & right & 13.667421 & 13.667423\\
1 & positive & left  & 15.805793 & 15.805796\\
1 & positive & right & 16.037329 & 16.037331\\
1 & negative & left  & 17.416405 & 17.416408\\
1 & negative & right & 17.436377 & 17.436379\\
2 & positive & left  & 19.493085 & 19.493088\\
2 & positive & right & 19.551252 & 19.551255\\
\bottomrule
\end{tabular}
\end{center}
These disjoint intervals prove all stated time comparisons. The positive
$n=0$ source is pre-cut, and all $n\ge1$ sources have $p>2K>p_1$.
For a negative-arc $n=0$ source, $K<p<2K$ and $V=3>0$ force
$-\cn p\,q-\dn p\sn p>0$, placing it after the unique Maxwell root
$p_1\in(K,2K)$. Hence the other eight sources are post-cut. Finally,
Corollary~\ref{cor:zero-winding-regimes} identifies this complete rotating
fiber with the zero-winding fiber.

For the connected comparison in \eqref{eq:period-action-example-gap},
the earlier parameterized certificate remains applicable to the two
smaller-$a$ positive-arc sources. In source coordinates
$s=-\log(1-m),\alpha$, its $256$ label intervals cover $[1,2]$.
Uniform Krawczyk contraction proves one root in each source rectangle
for every label. At all $257$ subdivision values, separate root boxes
lie inside the adjacent rectangles; uniqueness identifies the roots,
and the implicit function theorem joins one analytic branch. An independent
448-bit replay checks all $513$ inclusions with maximum contraction below
$.377$, preserving the exact affine label dependence before interval
multiplication. Its data are \nolinkurl{data/e121_period_action.json}
and \nolinkurl{data/e122_period_replay.json}. Endpoint clock evaluation
gives \eqref{eq:period-action-example-gap}.
\end{proof}

This example has $|U|/|V|=2>w$, outside the exterior fold window.
It gives a complete target-specific rotating count and time order, as
well as the selected connected continuation. Comparisons between
arbitrary targets or arbitrary period branches remain open.


\subsection{A complete nine-source interval}

\begin{proposition}[Two consecutive folds on a fixed target line]
\label{prop:nine-source-interval}
Let $(U,V,W)=(u,3,3/4)$ with $11/2\le u\le13/2$.
There are exactly two critical values of the rotating endpoint map on
this segment, denoted by $u_-<u_+$, with
\begin{equation}\label{eq:nine-source-boundaries}
 \begin{split}
 5.5172084304&<u_-<5.5172084325,\\
 6.4218626764&<u_+<6.4218626785.
 \end{split}
\end{equation}
Each has exactly one critical source, an ordinary post-cut fold:
its actual period is one, on the negative scalar arc, at $u_-$ and
three, on the positive arc, at $u_+$: the pair created at $u_-$ is the
family $\mathcal T_1^-$ of Corollary~\ref{thm:line-time-order}, the pair
created at $u_+$ is $\mathcal T_3^+$.
The complete zero-winding counts, including critical sources once, are
\begin{center}\small
\begin{tabular}{lccccc}
\toprule
target coordinate & total & period $0$ & period $1$ & period $2$ & period $3$\\
\midrule
$11/2\le u<u_-$ & $7$ & $3$ & $2$ & $2$ & $0$\\
$u=u_-$         & $8$ & $3$ & $3$ & $2$ & $0$\\
$u_-<u<u_+$     & $9$ & $3$ & $4$ & $2$ & $0$\\
$u=u_+$         & $10$ & $3$ & $4$ & $2$ & $1$\\
$u_+<u\le13/2$ & $11$ & $3$ & $4$ & $2$ & $2$\\
\bottomrule
\end{tabular}
\end{center}
All these sources are rotating. At each target exactly one is pre-cut;
all others are post-cut. Away from $u_\pm$ every source is regular.
\end{proposition}

\begin{proof}[Computer-assisted proof]
The sign argument in Proposition~\ref{prop:period-action-example}
applies throughout $u>3$. With $w=3/4$ it gives the exhaustive scalar charts
\begin{equation}\label{eq:line-source-arcs}
 \begin{gathered}
 a\in\left(\frac{3w}{u+w},1\right)
 \quad\hbox{or}\quad
 a\in\left(-1,-\frac{3w}{u-w}\right),\\
 \sigma=\frac{a^2[u^2+(3-a)^2]}{u^2a^2+(3-a)^2},\qquad
 m=\frac{a^2+w^2}{(1+w^2)\sigma},\qquad
 q_* =\frac{\sqrt{u^2+(3-a)^2}}{\sqrt{1+w^2}}.
 \end{gathered}
\end{equation}
Let $\alpha_0=\arcsin\sqrt\sigma$, $Q=K-E$ and
$B=F(\alpha_0\mid m)-E(\alpha_0\mid m)$. The residuals are
$f_n^+=q_*-2nQ-B$ and $f_n^-=q_*-2(n+1)Q+B$.
As before, $n$ is the actual period in both expressions.

An analytic monotonicity removes the need to search freely in both
$a$ and $u$. Fix $a\ne0$, put $r=(a^2+w^2)/(1+w^2)$, and use
$r<m<1$ as the second coordinate. Then
\[
 u=\frac{3-a}{|a|}\sqrt{\frac{r-ma^2}{m-r}},\qquad
 q_*\sqrt{m-r}=\frac{3-a}{|a|}\sqrt{r(1-r)}=:\Lambda(a),
 \qquad u_m<0.
\]
Changing variable $t=\sqrt m\sin\alpha$ in the defining integral gives
\[
 B_r(m)=\int_0^{\sqrt r}
 \frac{t^2\,dt}{\sqrt{m-t^2}\sqrt{1-t^2}},\qquad B_r'(m)<0.
\]
The derivative here holds $r$ fixed, not $\alpha_0$. Moreover,
\begin{equation}\label{eq:weighted-drift-monotonicity}
 \frac{d}{dm}\bigl[\sqrt{m-r}\,B_r(m)\bigr]
 =\frac{1}{2\sqrt{m-r}}\int_0^{\sqrt r}
 \frac{t^2(r-t^2)\,dt}{(m-t^2)^{3/2}\sqrt{1-t^2}}>0.
\end{equation}
Thus both
$H_n^+=\sqrt{m-r}(2nQ+B_r)$ and
$H_n^-=\sqrt{m-r}(2(n+1)Q-B_r)$ increase strictly in $m$.
For the second expression, use $2(n+1)Q-B_r>0$ and its positive
$m$ derivative. For positive $n\ge1$ and negative $n\ge0$, each $H_n$
increases from zero to infinity. Indeed, $B_r\to Q(r)$ as $m\downarrow r$,
whereas $B_r$ has a finite limit and $Q\to\infty$ as $m\uparrow1$.
There is therefore exactly one solution of $H_n=\Lambda$ for each signed
$a$. In particular $G_n=\sqrt{m-r}\,f_n=\Lambda-H_n$ has
$\partial_uG_n>0$, and $\partial_uf_n>0$ at every root. Each post-cut
family is a single analytic graph $u_n(a)$. This does not assert that it
has only one critical point.

The following finite certificate excludes all additional critical points
on the stated segment. Since $m_u<0$ and $(q_*)_u>0$, bounds at $u=13/2$
apply throughout the band. Four exact polynomial intervals prove $m>.71$
on the positive arc, and
$8Q(.71)>2\sqrt{233}/5\ge q_*$ excludes $n\ge4$.
On the negative arc, $346$ polynomial intervals retain the common
$a$ dependence and prove $5Q(m)>q_*$, excluding $n\ge2$.
All $350$ intervals have a separate exact-coverage and elliptic replay;
their data are \nolinkurl{data/e126_period_bounds.json} and
\nolinkurl{data/e126_period_replay.json}.

Only positive periods $1,2,3$ and negative periods $0,1$ remain post-cut.
Four short endpoint tails are excluded analytically by polynomial bounds
$m>.999$ on the positive arc and $m>.999999$ on the negative arc,
together with $2Q(.999)>q_*$ and $Q(.999999)>q_*$.
The remaining five compact parameter intervals have an exact partition
into $285$ intervals. On each one, signs of $f_n$ at fixed $u$, interpreted
using the monotonicity of $G_n$, either exclude the whole target band or
enclose every possible root in a smaller $u$ interval. On the resulting
rectangle, either $0\notin\partial_af_n$, or the rectangle is contained
in one of two independently certified local uniqueness boxes.
The criterion $f_n=\partial_af_n=0$ is equivalent to full conjugacy by
the nonsingular source and residual changes in the preceding proposition.
The positive $n=0$ source has $p<K$ and is regular.

The two retained boxes contain the points enclosed in
\eqref{eq:nine-source-boundaries}. Both have
$\partial_uf_n>0$ and $\partial_{aa}f_n<0$. To check the full fold
condition, write $b=a\cos\alpha/\sqrt{\sigma-a^2}>0$ and
$\widehat u=(3-a)/b$. The source change $(a,\sigma)\mapsto(a,\widehat u)$
has positive second-coordinate derivative. With
$R=U_{\rm actual}-\widehat u=-(D/d)f_n$, the two endpoint residuals are
$(R+\widehat u-u_*,bR)$. At $R=R_a=0$ their derivative has rank one,
and a cokernel applied to the second kernel derivative gives
$-bR_{aa}=b(D/d)(f_n)_{aa}\ne0$. Restoring $W$ proves an ordinary fold
of the full endpoint map. Each pair is born with increasing $u$.

An alternate implementation replays all $285$ intervals, all four tails
and both larger uniqueness boxes, using the original residual
$ud/D-q_n$ and generic directional derivatives. It uses $1768$ interval
evaluation nodes and rejects ten missing/duplicated cover controls.
The two implementations share Arb special functions; these are internal
independent replays, not external review. The complete data are
\nolinkurl{data/e127_critical_cover.json} and
\nolinkurl{data/e127_critical_replay.json}.

Finally this entire target segment lies strictly inside the rotating image.
If $P$ defines $\mathcal F(3,3/4)$, elementary bounds give
$3<P<4$ and $\coth P<5/4$, hence $\mathcal F(3,3/4)<3$.
Properness over the strict image and Corollary~\ref{cor:zero-winding-regimes}
exclude source arrivals from other regimes. Starting with the complete
$u=6$ count in Proposition~\ref{prop:period-action-example}, the two simple
pair births give the table. The period and arc labels cannot change on a
regular continuation, since the excluded quarter joins, period boundaries
and $a=0$ are the only such label boundaries. The same cut inequalities
as in that proposition give one pre-cut source and all remaining sources
strictly post-cut.
\end{proof}

\subsection{Analytical comparison within a signed period family}

The within-family comparison extends beyond the preceding target segment.
Here $u=U>0$, $V>1$ and $w=W>0$; all three are the target coordinates
already used above, with heading $2\arctan w$. The scalar $a$ remains
an auxiliary source parameter, not a heading angle.

\begin{theorem}[Component time order at general parameters]
\label{thm:general-pair-order}
Fix $V>1$ and $w>0$. For each positive arc with actual period $n\ge1$,
and each negative arc with actual period $n\ge0$, the rotating sources
form a proper real-analytic graph $u=g_n^\pm(a)$ on $(0,1)$ or $(-1,0)$,
respectively. Both ends of each graph tend to $+\infty$.
Write $T(a)$ for its travel time and $\lambda(a,u)$ for the coefficient
of $du$ in first variation at fixed $V,w$.

If $a_L<a_R$ are any two distinct sources in the same signed period
family at the same target $u$, then their moduli and momenta satisfy
\begin{equation}\label{eq:general-source-order}
 m_R<m_L,\qquad \lambda(a_R,u)>\lambda(a_L,u)>0.
\end{equation}
If, in addition, $[a_L,a_R]$ is a nondegenerate connected component of
$\{a:g_n^\pm(a)\le u\}$, then
\begin{equation}\label{eq:general-component-time}
 T(a_R)>T(a_L).
\end{equation}
Where its endpoints vary as regular source branches,
$d(T_R-T_L)/du>0$. In particular, these time conclusions apply whenever
the family has exactly two regular sources at the target. No uniqueness
of a minimum of $g_n^\pm$ is asserted.
\end{theorem}

\begin{proof}
The signed reconstruction in Proposition~\ref{prop:period-action-example}
uses only $u>0$ and $V>|a|$, so it applies for every $V>1,w>0$.
The identity $V=a+uaC/\sqrt{\sigma-a^2}$ forces $\operatorname{sgn}C
=\operatorname{sgn}a$ and excludes $a=0$ and $C=0$. Solving it with
this sign restored gives
\begin{equation}\label{eq:general-signed-source}
 \begin{gathered}
 \sigma=\frac{a^2[u^2+(V-a)^2]}{u^2a^2+(V-a)^2},\qquad
 m=\frac{a^2+w^2}{(1+w^2)\sigma},\qquad
 q_* =\frac{\sqrt{u^2+(V-a)^2}}{\sqrt{1+w^2}},\\
 a\in\left(\frac{wV}{u+w},1\right)
 \quad\hbox{or}\quad
 a\in\left(-1,-\frac{wV}{u-w}\right).
 \end{gathered}
\end{equation}
The positive arc is nonempty for $u>w(V-1)$; the negative one requires
$u>w(V+1)$. Put $\alpha_0=\arcsin\sqrt\sigma\in(0,\pi/2)$,
$Q(m)=K(m)-E(m)$ and $B(\alpha_0,m)=F(\alpha_0\mid m)-E(\alpha_0\mid m)$.
The source equations are
\begin{equation}\label{eq:general-signed-drift}
 q_*=2nQ+B\quad(a>0),\qquad
 q_*=2(n+1)Q-B\quad(a<0).
\end{equation}
The negative arc uses the full amplitude $\pi-\alpha_0$; $n$ remains
the actual period in both equations.

For completeness, the monotone-modulus argument of
\eqref{eq:weighted-drift-monotonicity} applies without fixing $V,w$.
For fixed $a$ put $r=(a^2+w^2)/(1+w^2)$. Then $r<m<1$ and
\[
 u=\frac{V-a}{|a|}\sqrt{\frac{r-ma^2}{m-r}},\qquad u_m<0,
 \qquad q_*\sqrt{m-r}
 =\frac{V-a}{|a|}\sqrt{r(1-r)}.
\]
The rescaled drifts $\sqrt{m-r}(2nQ+B_r)$ and
$\sqrt{m-r}(2(n+1)Q-B_r)$ increase strictly from zero to infinity.
Here $B_r=B(\arcsin\sqrt{r/m},m)$; the positive derivative of its weighted
form is \eqref{eq:weighted-drift-monotonicity}. On the negative arc use
$B_r<Q$ and $B_r'<0$. Thus each equation has exactly one modulus root
for each signed $a$, with nonzero modulus derivative. The implicit
function theorem gives the stated analytic graph and $f_u>0$ at every
root of $f=q_*-(2nQ+B)$ or $q_*-(2(n+1)Q-B)$.
The physical inequality $u>w(V-a)/|a|$ gives $u\to\infty$ as $a\to0$.
As $|a|\to1$, $r\to1$ and the drift is at least $2nQ(r)$ on the positive
arc, or $Q(r)$ on the negative arc. It diverges, forcing $u\to\infty$.
This proves properness, without a critical-point count.

At fixed target, direct differentiation gives
\begin{equation}\label{eq:general-amplitude-monotonicity}
 \sigma_a=\frac{2a(V-a)[u^2(V-a^3)+(V-a)^3]}
 {[u^2a^2+(V-a)^2]^2},\qquad
 (q_*)_a=-\frac{V-a}{(1+w^2)q_*}<0.
\end{equation}
The bracket is positive because $V>1$ and $|a|<1$. Thus $\alpha_0$
increases with $a$ on the positive arc and decreases on the negative arc.
From the defining integrals \cite[Equations 19.2.4--19.2.8]{DLMF},
\[
 B(\alpha,m)=\int_0^\alpha
 \frac{m\sin^2t}{\sqrt{1-m\sin^2t}}\,dt,\qquad
 \partial_m\frac{m\sin^2t}{\sqrt{1-m\sin^2t}}
 =\frac{\sin^2t(2-m\sin^2t)}{2(1-m\sin^2t)^{3/2}}\ge0.
\]
The last derivative is strictly positive for $0<t<\pi/2$. With
$0<\alpha<\pi/2$ fixed, $0<B_m<Q_m$ and $B_\alpha>0$.
Consequently $2nQ+B$ increases in both $m$ and $\alpha$, while
$2(n+1)Q-B$ increases in $m$ and decreases in $\alpha$.
If $m_R\ge m_L$, the signed amplitude change in either case would
make the right drift strictly larger. This contradicts
$q_*(a_R)<q_*(a_L)$ in \eqref{eq:general-signed-drift}.
Hence $m_R<m_L$.

First variation supplies the momentum comparison. Write
$S=\sin\alpha>0$, $C=\cos\alpha$, $z=a/(kS)$, $h=wd/(kS)$,
$d^2=1-mS^2$ and $D^2=1-mz^2$, where $\alpha$ is the full amplitude.
Jacobi subtraction \cite[Equations 22.8.1 and 22.8.3]{DLMF} gives
\[
 p_x=\frac{-zCd+ShD}{1-a^2},\qquad
 p_y=-\frac{Dd+k^2zShC}{k(1-a^2)}.
\]
The actual-period parity factors cancel. Since
$x=2(wu-V)/(1+w^2)$ and $y=-2(u+wV)/(1+w^2)$,
\begin{equation}\label{eq:line-clock-derivative}
 \lambda(a,u)=\frac{2(wp_x-p_y)}{1+w^2}
 =\frac{2D}{kd(1+w^2)}
 =\frac{2u|a|}{\sqrt{(a^2+w^2)[u^2a^2+(V-a)^2]}}>0.
\end{equation}
Here $1-a^2=(1+w^2)d^2$. The endpoint identity $u=q_*D/d$ implies
\begin{equation}\label{eq:general-momentum-product}
 \sqrt m\,q_*\lambda=\frac{2u}{1+w^2}.
\end{equation}
Both $m$ and $q_*$ strictly decrease from $a_L$ to $a_R$ at roots.
This proves \eqref{eq:general-source-order}.

It remains to fix the sign of the clock difference, rather than just
its derivative. Set $g=g_n^\pm$ and let $[A,B]$ be a nondegenerate
connected component of $\{g\le u\}$. Properness makes it compact inside
one signed arc, with $g(A)=g(B)=u$ and $v_0:=\min_{[A,B]}g<u$.
The clock is analytic along the graph, and first variation gives
$T'(a)=\lambda(a,g(a))g'(a)$, including at its critical points.
Define on $[A,B]$
\[
 H(a)=T(a)+\int_{g(a)}^u\lambda(a,v)\,dv.
\]
Then $H'(a)=\int_{g(a)}^u\partial_a\lambda(a,v)\,dv$ and
$H(A)=T(A)$, $H(B)=T(B)$. For almost every $v\in(v_0,u)$,
the section $\{a\in[A,B]:g(a)\le v\}$ is a finite union of intervals
$[l_j(v),r_j(v)]$. Indeed, the nonconstant analytic graph has only
finitely many critical points on this compact interval. Fubini's theorem
therefore yields
\begin{equation}\label{eq:component-horizontal-slices}
 T(B)-T(A)=\int_{v_0}^u\sum_j
 \bigl[\lambda(r_j(v),v)-\lambda(l_j(v),v)\bigr]\,dv>0.
\end{equation}
Every bracket is positive by \eqref{eq:general-source-order}; nonempty
regular sections occur on a set of positive measure. The integrand
$\partial_a\lambda$ is continuous on the compact integration region;
no pointwise sign of this derivative away from the source graph is needed.
This proves \eqref{eq:general-component-time}. For regular endpoints,
first variation directly gives $d(T_R-T_L)/du=\lambda_R-\lambda_L>0$.
Finally, if exactly two roots exist and both are regular, the graph's
infinite endpoint limits force them to bound its unique sublevel component.
\end{proof}

This theorem compares momenta within one signed period family and times
across one connected sublevel component. It does not order times between
separate components, different signs or different periods, or determine
the number of sources at an arbitrary target.

\subsection{Complete time order on the certified target segment}

\begin{corollary}[Time ordering throughout the fixed target segment]
\label{thm:line-time-order}
On the segment of Proposition~\ref{prop:nine-source-interval}, let
$\mathcal T_n^+(u)$ and $\mathcal T_n^-(u)$ be the sets of travel times
in the positive and negative scalar arcs, with actual period $n$.
For nonempty families, with $A<B$ meaning $\max A<\min B$,
\begin{equation}\label{eq:line-family-order}
 \mathcal T_0^+<\mathcal T_0^-<\mathcal T_1^+
 <\mathcal T_1^-<\mathcal T_2^+<\mathcal T_3^+,
\end{equation}
where absent families are omitted. Within each two-source family,
label the sources by $a_L<a_R$. Then $T_L<T_R$ and
\[
 \frac{d}{du}(T_R-T_L)>0
\]
on its regular existence interval. At either fold the merging source is
counted once. Thus all distinct zero-winding sources of any target on
the segment have different travel times.
\end{corollary}

\begin{proof}[Computer-assisted proof]
Proposition~\ref{prop:nine-source-interval} gives exactly two regular
sources in each present post-cut signed family away from its fold.
Theorem~\ref{thm:general-pair-order} makes them the endpoints of one
sublevel component and proves their strict time order and increasing gap.
At a fold the two sources coincide and are counted once.

For comparisons between families, \eqref{eq:line-clock-derivative}
also holds on the positive period-zero branch. With $w=3/4$,
\begin{equation}\label{eq:line-clock-bound}
 0<\lambda(a,u)<\frac2w=\frac83,
\end{equation}
because the exact reciprocal-square identity is
\[
 \frac1{\lambda^2}-\frac{w^2}{4}
 =\frac{a^2}{4}
   +\frac{(a^2+w^2)(V-a)^2}{4u^2a^2}>0.
\]
Continuous clock and momentum limits extend the integrated bounds to
fold endpoints. The earlier $14$-interval paired-momentum certificate,
its $38$ root enclosures, two fold corridors and three initial clock-gap
anchors remain separately replayed corroboration in
\nolinkurl{data/e129_pair_cover.json} and
\nolinkurl{data/e129_pair_replay.json}; the general analytical theorem
supplies the within-family proof here.

Finally, any inter-family clock gap has derivative of absolute value less
than $8/3$, since both individual derivatives belong to $(0,8/3)$.
The certified $u=6$ clocks in Proposition~\ref{prop:period-action-example}
give $32$ comparisons between its distinct signed period families.
After subtracting $4/3$, every gap still exceeds $.0457$; since
$|u-6|\le1/2$, their order persists wherever both branches exist.
This includes the direct comparison of $\mathcal T_1^+$ with
$\mathcal T_2^+$ when $\mathcal T_1^-$ is absent. Both
$\mathcal T_3^+$ clocks increase from their common upper-fold clock.
Subtracting the largest earlier clock at $u=6$ and $4/3$ from its
certified lower bound leaves a margin greater than $2.524$.
Thus $\mathcal T_3^+$ follows every earlier family. The exact rational
comparison audit is \nolinkurl{data/e129_analytic_audit.json}.
Equation~\eqref{eq:line-family-order} follows, including the fold targets
by continuity and counting their single critical sources once.
\end{proof}

This is a time-order result for the complete zero-winding fiber on the
specified segment. It makes no comparison with nonzero winding classes
or with branches at arbitrary targets.


\subsection{Simultaneous scaling of the endpoint and period}
\label{sec:dilation}

Fix $U_0>0$, $V_0>1$, $w>0$ and $\nu_0\ge1$. We now follow selected
period families while both the endpoint and their effective label grow:
\begin{equation}\label{eq:dilation-family}
 (U,V,W)=(sU_0,sV_0,w),\qquad
 P_{s\nu_0},\quad N_{s\nu_0-1},\qquad s\ge1.
\end{equation}
Here $P_n$ and $N_n$ denote the positive and negative scalar arcs with
\emph{actual} period $n$. At noninteger labels these are the formal
analytic families of Section~\ref{sec:period-action}; they represent
geodesics only when the actual labels are integers. The results below
count these two families, not all periods or the full endpoint fiber.

Put $c=\sqrt{1+w^2}$, $Q=K-E$, $a=\epsilon t$ with
$\epsilon\in\{1,-1\}$ and $0<t<1$, and use the short amplitude
$\alpha\in(0,\pi/2)$. The reconstruction
\eqref{eq:general-signed-source}--\eqref{eq:general-signed-drift} reads
\begin{equation}\label{eq:dilation-source-clock}
 \begin{gathered}
 R=\sqrt{U^2+(V-\epsilon t)^2},\qquad q=R/c,
 \qquad B=F(\alpha\mid m)-E(\alpha\mid m),\\
 q=2s\nu_0Q+\epsilon B,\qquad
 T=4s\nu_0\sqrt m K+\epsilon\,2\sqrt m F(\alpha\mid m).
 \end{gathered}
\end{equation}
For $\epsilon=-1$, the full amplitude is $\pi-\alpha$. Its incomplete
integrals are $2K-F$ and $2E-E(\alpha\mid m)$; this explains the
minus-one shift in \eqref{eq:dilation-family}.

\begin{lemma}[Signed action and a positive intercept]
\label{lem:dilation-action}
On either formal source family, with $w$ fixed and effective label $\nu$,
\begin{equation}\label{eq:dilation-signed-action}
 dT=\ell_U\,dU+\ell_V\,dV+\mathcal A(m)\,d\nu,
 \quad (\ell_U,\ell_V)=\frac{2(U,V-\epsilon t)}{c\sqrt m R},
\end{equation}
where $\mathcal A$ is given in \eqref{eq:period-action-identity}.
Define, with $\sin^2\alpha=(t^2+w^2)/(c^2m)$,
\[
 j=c[E(\alpha\mid m)-(1-m)F(\alpha\mid m)],\qquad
 p=\frac{t\cos\alpha}{c\sin\alpha\sqrt{1-m\sin^2\alpha}},
 \qquad I=j-tp.
\]
Then $j_t=p$ at fixed $m$, $p=(V-\epsilon t)/R$ at a source, and
\begin{equation}\label{eq:dilation-positive-intercept}
 \begin{split}
 I={}&cm(1-m)\int_0^\alpha
       \frac{\sin^2z}{(1-m\sin^2z)^{3/2}}\,dz\\
 &+\frac{w^2\cos\alpha}
 {c\sin\alpha\sqrt{1-m\sin^2\alpha}}>0.
 \end{split}
\end{equation}
Consequently every regular $C^1$ source branch parametrized by $s$ in
\eqref{eq:dilation-family} satisfies
\begin{equation}\label{eq:dilation-euler}
 \frac{d(T/s)}{ds}=-\frac{\epsilon\,2I}{c\sqrt m\,s^2}.
\end{equation}
\end{lemma}
\begin{proof}
The defining elliptic derivatives \cite[Sections 19.2 and 19.4]{DLMF}
give $j_t=p$ and $j_m=cF(\alpha\mid m)/2$, with the other variable
fixed. The amplitude derivative is included in the latter identity.
The geometric reconstruction is
\[
 U=cq\sqrt{1-p^2},\qquad V=\epsilon t+cqp.
\]
It gives
$\ell_U dU+\ell_V dV=2\,dq/\sqrt m+\epsilon\,2p\,dt/(c\sqrt m)$.
On using $q=2\nu Q+\epsilon B$, the clock decomposes as
\[
 T=2q/\sqrt m+\nu\mathcal A(m)+\epsilon\,2j/(c\sqrt m).
\]
Differentiate this expression. The $dt$ terms use $j_t=p$, and its
remaining $dm$ coefficient vanishes because
$\mathcal A'=2Q/m^{3/2}$ and $2mj_m-j=cB$. This proves
\eqref{eq:dilation-signed-action} for both signs, without assigning
physical first variation to a noninteger period.

Write $S=\sin\alpha$, $C=\cos\alpha$, $d=\sqrt{1-mS^2}$.
Since $t^2=c^2mS^2-w^2$,
\[
 I=c[E(\alpha\mid m)-(1-m)F(\alpha\mid m)-mSC/d]
   +w^2C/(cSd).
\]
The expression in brackets vanishes at $\alpha=0$ and its derivative
is $m(1-m)\sin^2\alpha/(1-m\sin^2\alpha)^{3/2}$.
This proves \eqref{eq:dilation-positive-intercept}. Finally, apply
\eqref{eq:dilation-signed-action} along the scaling and subtract the
clock decomposition from $sT'$. The complete-period action cancels and
$sT'-T=-\epsilon\,2(j-tp)/(c\sqrt m)$, proving
\eqref{eq:dilation-euler}.
\end{proof}

\begin{lemma}[Uniform normalized limit]
\label{lem:dilation-limit}
Let $R_0=\sqrt{U_0^2+V_0^2}$ and $q_0=R_0/c$. There is a unique
$m_0\in(0,1)$ satisfying $2\nu_0Q(m_0)=q_0$. Every source in
\eqref{eq:dilation-family} that exists at scale $s$ satisfies, uniformly
in the sign and source choice as $s\to\infty$,
\begin{equation}\label{eq:dilation-common-limit}
 m=m_0+O(s^{-1}),\qquad T/s=C_0+O(s^{-1}),\qquad
 C_0=4\nu_0\sqrt{m_0}K(m_0).
\end{equation}
This assertion does not assume or imply existence at every scale.
\end{lemma}
\begin{proof}
The reverse triangle inequality and $0<B<Q$ give
\[
 |q/s-q_0|<\frac1{cs},\qquad
 |2\nu_0Q-q_0|<\frac{Q+1/c}{s}.
\]
Since $q_0>1/c$ and $\nu_0\ge1$, all such moduli belong to the fixed
compact subinterval determined by
\[
 0<\frac{q_0-1/c}{2\nu_0+1}<Q(m)
 <\frac{q_0+1/c}{2\nu_0-1}=:Q_{\rm hi}<\infty.
\]
Moreover $|Q(m)-q_0/(2\nu_0)|\le(Q_{\rm hi}+1/c)/(2\nu_0s)$.
Use $Q'=E/[2(1-m)]>0$ to invert this bound on the compact interval.
The clock assertion follows from \eqref{eq:dilation-source-clock},
since $0<F(\alpha\mid m)<K(m)$ there.
\end{proof}

If a positive source continues regularly for every $s\ge1$ along this
scaling, Lemmas~\ref{lem:dilation-action}--\ref{lem:dilation-limit}
give $T_P(1)>C_0$. Such a negative continuation gives $T_N(1)<C_0$.
These are conditional comparisons: continuation is an essential hypothesis.
The obstruction below shows that even joint initial family existence
does not guarantee persistence to large scales.

\begin{theorem}[Eventual counts and clocks away from the critical surface]
\label{thm:dilation-strict}
Suppose $U_0>wV_0$ and put
\[
 r=U_0/R_0,\quad \beta=V_0/R_0,\quad t_b=wV_0/U_0,
 \qquad m_\star=\frac{(\beta+wr)^2}{c^2}.
\]
If $m_0>m_\star$, then for all sufficiently large $s$ each selected
sign family has exactly two regular sources. If $m_0<m_\star$, both
families are eventually empty. They are also eventually empty when
$U_0\le wV_0$.

In the four-source case, their limiting absolute source coordinates
$t_L<t_R$ satisfy $t_Lt_R=t_b$. Define $I_L,I_R$ by
\eqref{eq:dilation-positive-intercept} at $(t_i,m_0)$ and put
$D_i=2I_i/(c\sqrt{m_0})$. Then $0<D_L<D_R$ and
\begin{equation}\label{eq:dilation-clock-offsets}
 T_{\epsilon,i}(s)=sC_0+\epsilon D_i+O(s^{-1}),\qquad
 T_{+,i}(s)-T_{-,i}(s)=2D_i+O(s^{-2}).
\end{equation}
Thus eventually
$T_{-,R}<T_{-,L}<T_{+,L}<T_{+,R}$.
The scale threshold depends on the fixed base parameters and is not
quantified here. Equality $m_0=m_\star$ is treated separately below.
\end{theorem}
\begin{proof}
Put $\tau=1/s$ and, for $0<x<1$, define
\begin{equation}\label{eq:dilation-residual}
 \begin{gathered}
 z_\epsilon=V_0-\epsilon\tau x,\qquad
 \sigma_\epsilon=\frac{x^2(U_0^2+z_\epsilon^2)}{U_0^2x^2+z_\epsilon^2},
 \quad \mu_\epsilon=\frac{x^2+w^2}{c^2\sigma_\epsilon},\\
 H_\epsilon(x,\tau)=\frac{\sqrt{U_0^2+z_\epsilon^2}}c
       -2\nu_0Q(\mu_\epsilon)
       -\epsilon\tau B(\arcsin\sqrt{\sigma_\epsilon},\mu_\epsilon).
 \end{gathered}
\end{equation}
These are analytic near every interior limiting source, with the exact
reflection $H_-(x,\tau)=H_+(x,-\tau)$. At $\tau=0$ the geometric
modulus obeys
\begin{equation}\label{eq:dilation-square}
 \mu_\infty(x)=\frac{r^2x^2+r^2w^2+\beta^2+\beta^2w^2/x^2}{c^2},
 \qquad
 \mu_\infty(x)-m_\star=\frac{(rx-\beta w/x)^2}{c^2}.
\end{equation}
The limiting physical arc is $(t_b,1)$, with modulus one at either
endpoint and a unique minimum at $x=\sqrt{t_b}$. For $m_0>m_\star$
the equation $\mu_\infty(x)=m_0$ has two simple roots in that arc.
They solve
\[
 r^2 y^2+(r^2w^2+\beta^2-c^2m_0)y+\beta^2w^2=0,
 \qquad y=x^2,
\]
so $t_Lt_R=t_b$. The analytic implicit-function theorem applied to
$H_\epsilon$ gives one source near each root for sufficiently small
$|\tau|$. All other sources are excluded by compactness: a sequence
of extra sources at unbounded scales would have a limiting absolute
coordinate in $[t_b,1]$, while Lemma~\ref{lem:dilation-limit} forces
its modulus to $m_0$. The endpoints, with modulus one, cannot occur;
an interior limit must be $t_L$ or $t_R$, contradicting local uniqueness.
The same argument proves eventual emptiness when $m_0<m_\star$.
If $U_0<wV_0$, the positive physical arc eventually disappears and the
negative arc is empty. If $U_0=wV_0$, any possible positive source has
$x\to1$ and hence $m\to1$, contradicting the uniform compact modulus
bound. Simplicity transfers to full-map regularity in the source chart.

For either limiting root, $p(t_i,m_0)=\beta$. Expanding the source
condition gives, with $B_i=B(\alpha_i,m_0)$,
\[
 m_{\epsilon,i}=m_0-\frac{\epsilon}{s}
       \frac{B_i+\beta t_i/c}{2\nu_0Q'(m_0)}+O(s^{-2}).
\]
Using $(\sqrt m K)'=E/[2\sqrt m(1-m)]$ in the clock formula yields
the first expansion in \eqref{eq:dilation-clock-offsets}.
The reflection of $H_\epsilon$ extends to the locally unique analytic
roots and their normalized clocks. Their even powers of $\tau$ cancel
in the difference, giving its stated $O(s^{-2})$ remainder.
At fixed $m_0$, $p(t)>\beta$ for $t_L<t<t_R$: the sign of
$p(t)^2-\beta^2$ is the opposite of the quadratic above, with a
positive denominator. Since $j_t=p$,
\[
 I_R-I_L=\int_{t_L}^{t_R}[p(t)-\beta]\,dt>0.
\]
Together with Lemma~\ref{lem:dilation-action}, this proves the order.
\end{proof}

\begin{theorem}[Critical scaling classification]
\label{thm:dilation-critical}
Assume $U_0>wV_0$ and $m_0=m_\star$, and put $t=\sqrt{wV_0/U_0}$.
Write $H=H_+$ in \eqref{eq:dilation-residual}, and evaluate the following
partial derivatives at $(x,\tau)=(t,0)$, keeping the base target fixed:
\begin{equation}\label{eq:dilation-critical-coefficients}
 A=-H_{xx}/2>0,\qquad L=H_\tau,\qquad b=H_{x\tau},\qquad
 c_2=H_{\tau\tau}/2,\qquad \kappa=c_2+b^2/(4A).
\end{equation}
For all sufficiently large $s$, the complete selected-family counts are
\[
\begin{array}{c|cc}
 &P_{s\nu_0}&N_{s\nu_0-1}\\ \hline
 L>0&2&0\\
 L<0&0&2\\
 L=0,\ \kappa>0&2&2\\
 L=0,\ \kappa<0&0&0
\end{array}
\]
All counted sources are regular. The case $L=\kappa=0$ is excluded.
If $\epsilon L>0$, their absolute coordinates are
$t\pm\sqrt{\epsilon L/A}\,s^{-1/2}+O(s^{-1})$.
If $L=0,\kappa>0$, they are
$t+[\epsilon b/(2A)\pm\sqrt{\kappa/A}]s^{-1}+O(s^{-2})$.
No explicit starting scale or critical clock-gap exponent is asserted.
\end{theorem}
\begin{proof}
The square identity \eqref{eq:dilation-square} gives
$H(t,0)=H_x(t,0)=0$ and
$A=8\nu_0Q'(m_0)r^2/c^2>0$.
The analytic implicit-function theorem gives a unique nearby critical
point $x_c(\tau)$, a strict maximum on a fixed small interval about $t$.
Its endpoints retain negative residual values. Taylor expansion gives
\[
 x_c(\tau)=t+\frac b{2A}\tau+O(\tau^2),\qquad
 g(\tau):=H(x_c(\tau),\tau)=L\tau+\kappa\tau^2+O(\tau^3).
\]
Strict concavity gives exactly two simple roots when $g>0$ and none
when $g<0$. The negative family has maximum value $g(-\tau)$.
The compact exclusion in the preceding proof applies also at equality:
every possible source at unbounded scales must converge to the unique
limiting root $t$. Hence these local counts are complete for the selected
families. Finally, when $L\ne0$, substitute $\tau=h^2$, $x=t+hy$ and
divide the residual by $h^2$. The limiting polynomial is
$-Ay^2+\epsilon L$, with two simple roots for the winning sign.
When $L=0,\kappa>0$, use $x=t+\tau y$ and divide by $\tau^2$;
the limiting polynomial is $-Ay^2+\epsilon by+c_2$.
Analytic implicit inversion of these two polynomials proves the expansions.
\end{proof}

The coefficients are explicit in the critical target parameters.
For $0<t<1$, $w>0$, set
\begin{equation}\label{eq:dilation-critical-target}
 \begin{gathered}
 D=t^4+w^2,\quad r=w/\sqrt D,\quad \beta=t^2/\sqrt D,
 \quad m=\frac{(t^2+w^2)^2}{c^2D},\\
 R_0=2c\nu_0Q(m),\quad (U_0,V_0)=R_0(r,\beta),
 \quad \alpha=\arcsin\sqrt{D/(t^2+w^2)}.
 \end{gathered}
\end{equation}
Require $\nu_0\ge1$ and $V_0>1$. These parameters describe exactly the
critical targets under consideration. In particular,
\begin{equation}\label{eq:dilation-explicit-L}
 L=\frac{t}{c\sqrt D}
   \left[\frac{E(m)(t^2+w^2)}{Q(m)(1-t^2)}-t^2\right]
   -B(\alpha,m),
\end{equation}
which is independent of $\nu_0$. There is no single-transition assumption.

\begin{proposition}[Certified critical examples]
\label{prop:dilation-critical-examples}
At $w=3/10$, the parameter choices $t=3/5$ and $t=4/5$, both with
$\nu_0=10$, have $L>0$ and $L<0$, respectively. There are at least three
distinct zeros of $L$ at this same $w$. Each interval below contains
exactly one zero; with its listed label it satisfies $V_0>1$ and the
listed sign of $\kappa$:
\[
\begin{array}{c|r|c}
 \text{interval for }t&\nu_0&\operatorname{sgn}\kappa\\ \hline
 (97563304426967701622,97563304426967701623)/10^{22}\ (\approx0.00976)&50000&-\\
 (74222693418376774839,74222693418376774840)/10^{20}&10&+\\
 (87334017356145181518,87334017356145181519)/10^{20}&10&-
\end{array}
\]
Thus all four outcomes of Theorem~\ref{thm:dilation-critical} occur.
The three intervals do not establish an exhaustive zero count.
\end{proposition}
\begin{proof}[Computer-assisted proof]
The 384-bit certificate and its producer are
\begin{center}\small
\nolinkurl{data/e136_critical_certificate.json}\\
\nolinkurl{research/e136_critical_certificate.py}.
\end{center}
They check the real domain and $V_0>1$ throughout each rational input interval.
It bounds $L$ with opposite signs at each pair of endpoints and bounds
$L'$ away from zero throughout the interval, proving existence and
uniqueness of its zero. The derivative $L'$ here varies the base critical
target in \eqref{eq:dilation-critical-target}; it is not a partial
derivative at a fixed endpoint. It also encloses $\kappa$ with the stated
strict sign throughout each interval. The two rational point inputs
have the asserted strict $L$ signs. Full coefficient formulas and their
chain-rule reductions are recorded in the accompanying proof audit
\nolinkurl{references/e136/proof_audit.md}. An independent implementation
at 140 decimal digits differentiates the full-amplitude residual and
places all 45 coefficient and target values inside the raw outward
intervals. That replay corroborates the certificate; the interval signs
and the analytical classification establish the stated outcomes.
\end{proof}

\begin{proposition}[Joint existence without persistence]
\label{prop:dilation-obstruction}
There is an exact target with at least two regular $P_2$ sources and two
regular $N_1$ sources, but no source in either scaled family $P_{2s}$,
$N_{2s-1}$ for any real $s\ge1024$.
\end{proposition}
\begin{proof}[Computer-assisted proof with a uniform tail exclusion]
Use a positive source of actual period two, with
\[
 m=41/50,\qquad \alpha=3/10,\qquad w=3/100.
\]
Set $c=\sqrt{1+w^2}$, $S=\sin\alpha$, $C=\cos\alpha$,
$d=\sqrt{1-mS^2}$, and put
\[
 a=\sqrt{c^2mS^2-w^2},\quad p=\frac{aC}{cSd},\quad
 q=4Q(m)+B(\alpha,m),\quad
 U_0=cq\sqrt{1-p^2},\quad V_0=a+cqp.
\]
These exact expressions give
\[
 (U_0,V_0,w)\approx(2.06467730131716,4.34439453433560,.03).
\]
The 320-bit certificate \nolinkurl{data/e135_joint_obstruction.json}
encloses four distinct roots in the signed $a$ intervals
\[
\begin{array}{c|r}
 N_1&(-2729390928,-2729390924)/10^{10}\\
 N_1&(-2548570550,-2548570545)/10^{10}\\
 P_2&(2250499601,2250499604)/10^{10}\\
 P_2&(2660391869,2660391872)/10^{10}
\end{array}
\]
Opposite endpoint signs and a nonzero derivative enclosure prove one
regular root in each box. No exclusion of further initial roots is used.
The same certificate proves
$27/10^5<m_\star-m_0<29/10^5$.

For a hypothetical scaled source write $z=V_0-\epsilon t/s$.
The square identity gives
\[
 m\ge M(z):=\frac{(z+wU_0)^2}{c^2(U_0^2+z^2)}.
\]
For $S_0=1024$, the certificate checks $U_0>w(V_0+1/S_0)$, so
$M$ increases on $[V_0-1/S_0,V_0+1/S_0]$. Conversely, the source
equation and $B<Q$ imply, for every $s\ge S_0$,
\[
 Q(m)\le\frac{q_0+1/(cs)}{2\nu_0-1/s}
 \le\frac{q_0+1/(cS_0)}{2\nu_0-1/S_0},\qquad \nu_0=2.
\]
Outward evaluation gives
\[
 Q(M(V_0-1/S_0))-
 \frac{q_0+1/(cS_0)}{2\nu_0-1/S_0}>0.000226.
\]
This contradicts both bounds for every real $s\ge S_0$, proving the
entire tail empty. The separate full-amplitude replay
\nolinkurl{data/e135_joint_replay.json} corroborates the four initial
roots and interval quantities at 110 digits. It is not a second outward
certificate. No exact disappearance scale, initial exhaustive count or
clock-order reversal is asserted.
\end{proof}

The last example prevents an argument for arbitrary finite-target time
order based solely on continuation to infinity. Bounded source components
require a separate comparison. Likewise, the critical case $L=\kappa=0$
and a complete global description of the zeros in
\eqref{eq:dilation-explicit-L} remain open.

\section{Boundary directions: a logarithmic drift scale}\label{sec:boundary}

Fix $\alpha\in(-\pi/2,\pi/2)$ and first take
$b_0=-\pi/2$, $\theta=\alpha+\pi/2$.
Write $s=\sin\alpha$, $c=\cos\alpha>0$ and $q_R=(R\ee(\alpha),\theta)$.
The identity-fixing isometry $(x,y,\theta)\mapsto(-x,y,-\theta)$
gives the other boundary with $\alpha$ replaced by $-\alpha$.

\begin{lemma}[Complete boundary scalar chart]\label{lem:boundary-chart}
For sufficiently large $R$, the entire rotating inverse has one
component, parametrized by $z\in\mathbb R$. Put
\begin{equation}\label{eq:boundary-aux}
 w=\tanh z,\quad J=\sqrt{R^2+2Rw+1},\quad
 D=\sqrt{c^2+(s+J)^2},\quad d=\frac{c^2}{D^2}.
\end{equation}
The exact endpoint quantities are
\begin{equation}\label{eq:boundary-chart}
 \begin{gathered}
 B=-\frac{Rc}{D},\quad A=\frac{R(s+J)}D,\quad
 \sin a=\frac{1+sJ}D,\quad\cos a=\frac{cJ}D,\\
 \sin b=-\frac{s+J}D,\quad\cos b=\frac cD,\quad
 f=-\frac{c(R+w)}D,\quad g=-\frac{cw}D,\quad
 m=1-d\,\operatorname{sech}^2z.
 \end{gathered}
\end{equation}
Here $z$ is a new coordinate, not the tangent coordinate of
Proposition~\ref{prop:inverse}. With $q_a=q(\sin a,m)$ from
\eqref{eq:short}, define
\begin{equation}\label{eq:boundary-G}
 G(z)=\int_0^z
 \sqrt{1-d(z)\frac{\cosh^2v}{\cosh^2z}}\,dv,\qquad
 P_R=A-q_a+G,\qquad N_R=\frac{P_R}{L(m)}-\frac14.
\end{equation}
The intersections $N_R(z)=n\in\mathbb Z_{\ge0}$ are in bijection with
the rotating sources, with actual period index $n$.
\end{lemma}
\begin{proof}
The negative spatial sign is excluded when $Rc>2$.
For the positive sign, $a=\alpha-\arcsin(B/R)$ and
$b=-\pi/2-\arcsin(B/R)$. Principal $b$ forces $B<0$ and gives
$\cos b=-B/R$. Since $m=\sin^2b+g^2$, the condition $m<1$ is
equivalent to $|Rg/B|<1$. Set $w=Rg/B$.
Solving the algebraic endpoint equations with $A/(-B)>0$ gives
\eqref{eq:boundary-chart}. Conversely every $-1<w<1$ gives strict
principal $a,b$, $A>0$, $0<m<1$, and
$\sin(a-b)=c$, $\cos(a-b)=-s$. It is therefore admissible.
There is no $B=0$ source: that would require
$\alpha=\theta/2$, contrary to this boundary condition.
Moreover
\[
 \frac{dB}{dw}=\frac{R^2c(s+J)}{JD^3}>0.
\]
Both ends have $m=1$ and are excluded. The point $g=0$ at $w=0$
is included.

For completeness, set
$u_a=F(\arcsin(\sin a/\sqrt m)\mid m)$ and
$j=F(\arcsin(w/\sqrt m)\mid m)$.
The latter is well defined because
$m-w^2=(1-w^2)(1-d)>0$.
The full phase lifts are $\psi=2K-u_a$ and $v=3K-j$.
Their residual difference is $\delta=K+u_a-j$.
For $w<0$, $3K<v<4K$ and $K<\psi<3K$, so $\delta>0$.
For $w\ge0$, set $u_b=j-K$; then $\delta=u_a-u_b>0$ because
$a>b$. Both phase lifts are in $(0,4K)$, so $\delta<4K$.
Jacobi addition for the elliptic epsilon function gives
\[
 Q_0=(K-E)+q_a-G,\qquad
 G=j-E(\arcsin(w/\sqrt m)\mid m)
           +\frac{w\sqrt{m-w^2}}{\sqrt{1-w^2}}.
\]
Differentiating the last expression at fixed $m$ and checking its
value at zero gives
$G=\int_0^w\sqrt{m-u^2}/(1-u^2)^{3/2}\,du$.
Substitute $u=\tanh v$ to obtain \eqref{eq:boundary-G};
$d(z)$ is held fixed during this integral.
The longitudinal equation is exactly $N_R=n$, with reconstructed
time $t=\sqrt m(\delta+4nK)$. Proposition~\ref{prop:inverse}
then proves the claimed bijection.
\end{proof}

\begin{lemma}[Uniform estimates on the whole scalar line]\label{lem:boundary-uniform}
Put $Q_\alpha=\atanh(s)-s$, $C_R=\log(4R/c)-1$, and
$t_z=\operatorname{sech}^2z$. There are constants $\kappa_1,\kappa_2$
depending only on the fixed angle such that, for all real $z$
and sufficiently large $R$,
\begin{align}
 P_R&=R-Q_\alpha+z+p_R,&
 |p_R|&\le \kappa_1/R,\notag\\
 |p_R'|+|p_R''|&\le \kappa_1(1+|z|)t_z/R^2,\label{eq:boundary-P}\\
 L_R:=L(m)&=4(C_R+\log\cosh z)+e_R,&
 |e_R|&\le \kappa_2/R,\notag\\
 |e_R'|+|e_R''|&\le
 \kappa_2t_z\{R^{-1}+R^{-2}(\log R+|z|)\}.\label{eq:boundary-L}
\end{align}
\end{lemma}
\begin{proof}
Uniformly for $-1\le w\le1$,
$J=R+w+O(R^{-1})$ and $D=R+s+w+O(R^{-1})$, also after
two $w$ derivatives of the remainders. Hence $d=O(R^{-2})$ and
$d',d''=O(t_zR^{-3})$.
The angle $a=\alpha+\arcsin(c/D)$ stays in a fixed strict principal
interval; $m-\sin^2a$ stays bounded away from zero.
The required partial derivatives of the short integral are therefore
bounded. It follows that $A-R,q_a-Q_\alpha=O(R^{-1})$, with their
first two $z$ derivatives $O(t_zR^{-2})$.

Expand $\sqrt{1-x}=1-\sum_{j\ge1}a_jx^j$, where $0<a_j\le1$.
For $d\le1/2$, \eqref{eq:boundary-G} becomes
\[
 G-z=-\sum_{j\ge1}a_jd(z)^jT_j(z),\qquad
 T_j(z)=\operatorname{sech}^{2j}z\int_0^z\cosh^{2j}v\,dv.
\]
Integration by parts gives
\[
 T_0=z,\qquad T_j=\frac{\tanh z}{2j}
             +\frac{2j-1}{2j}t_zT_{j-1}.
\]
In particular $|T_j|\le|T_1|\le1$.
Iteration expresses $T_j$ as at most $j+1$ terms of the forms
$w t_z^l$ and $z t_z^j$, with coefficients bounded by one.
Differentiating these expressions yields
$|T_j'|+|T_j''|\le Cj^3(1+|z|)t_z$.
The convergent sum $\sum j^3d^j$, together with the bounds on
$d',d''$, proves $G-z=O(R^{-2})$ and the derivative bounds in
\eqref{eq:boundary-P}. This treats both infinite tails.

Set $\epsilon=1-m=d t_z$.
The convergent complementary-parameter series in
\cite[Section 19.12]{DLMF}, differentiated termwise for
$0<\epsilon\le\epsilon_0<1$, give
\[
 L=4\log(4/\sqrt\epsilon)-4+\Lambda(\epsilon),
\]
\[
 |\Lambda|\le C\epsilon(1+|\log\epsilon|),\quad
 |\Lambda'|\le C(1+|\log\epsilon|),\quad
 |\Lambda''|\le C/\epsilon.
\]
Here $\epsilon',\epsilon''=O(\epsilon)$ and
$|\log\epsilon|=O(\log R+|z|)$. Composition gives the stated
remainder bounds. The other error is $4\log(D/R)=O(R^{-1})$;
its first two derivatives are $O(t_z/R)$.
Since every fixed power of $|z|$ times $t_z$ is bounded,
these function and derivative estimates are uniform on $\mathbb R$.
\end{proof}

\begin{theorem}[Boundary count and folds]\label{thm:boundary-count}
For sufficiently large $R$, $N_R$ has exactly one critical point,
a nondegenerate maximum $M_R=N_R(z_R)$. Its endpoint limits are
$-1/2$ and $0$, both approached from above, and
\begin{equation}\label{eq:boundary-count}
 \nu_2(q_R)=2\lfloor M_R\rfloor+1-\ind_{\{M_R\in\mathbb Z\}}
 =\frac{R}{2(\log(4R/c)-1)}+O(1).
\end{equation}
All these sources are regular except an integer-maximum source,
which is a full-map fold.
\end{theorem}
\begin{proof}
Set $H_R=P_R'L_R-P_RL_R'$, so $N_R'=H_R/L_R^2$.
Lemma~\ref{lem:boundary-uniform} gives, uniformly on the whole line,
\begin{equation}\label{eq:boundary-H}
 \frac{H_R}4=C_R+I(z)-(R-Q_\alpha)\tanh z+O(1),\qquad
 I(z)=\log\cosh z-z\tanh z\in[-\log2,0].
\end{equation}
All error products contain a decaying $t_z$ when multiplying an
unbounded factor, so they are uniformly $O(1)$.
For $z\le0$ this is positive eventually; for
$z\ge\atanh(1/2)$ it is negative eventually.
On the remaining fixed interval,
$P_R\asymp R$, $L_R''=4t_z+O(R^{-1})>0$,
$P_R''=O(R^{-2})$, and $L_R=O(\log R)$. Thus
$H_R'=P_R''L_R-P_RL_R''<0$ there.
This proves one nondegenerate maximum and excludes critical points
escaping toward either end.

Equations \eqref{eq:boundary-P}--\eqref{eq:boundary-L} give
$N_R(-\infty)=-1/2$ and $N_R(+\infty)=0$.
The derivative signs show approach from above.
Level zero has one root on the left; each positive integer below
$M_R$ has two, and an integer maximum has one. This proves the
floor formula.
At the maximum \eqref{eq:boundary-H} gives
$R\tanh z_R=C_R+O(1)$ and $z_R=O(\log R/R)$.
Substitution yields the useful refinement
\[
 M_R=\frac R{4C_R}-\frac14-\frac{Q_\alpha}{4C_R}
                  +O(C_R^{-2}),
\]
which proves the count asymptotic with bounded error.
The changes $z\mapsto w\mapsto B\mapsto a$ have nonzero
derivative at every finite point, and $B$ never vanishes.
Proposition~\ref{prop:fold} therefore gives the regularity and fold
statements for the full map.
\end{proof}

\begin{proposition}[Full fiber on the boundary]\label{prop:boundary-full}
Eventually the zero-winding fiber on either boundary consists exactly
of the sources in Theorem~\ref{thm:boundary-count}. Its unique
cell-zero source is strictly pre-cut and is the unique minimizer.
There are no $C_3,C_4,C_5$ sources.
\end{proposition}
\begin{proof}
On the lower boundary, \eqref{eq:UVW} becomes
$W=\tan((\alpha+\pi/2)/2)>0$, $U=RW/2$, $V=R/2$.
The $C_3$ scalar equation \eqref{eq:c3surface} has $p>V$ and gives
\[
 p-V=\frac{W^2p^2}{(p+V)\sinh^2p}
       \le\frac{W^2V}{\sinh^2V}.
\]
Here $p/\sinh^2p$ decreases on $p>0$.
Likewise $2p/(e^{2p}-1)$ decreases, so
\[
 p\coth p-1-V\le T_R-1,\qquad
 T_R=\frac{W^2V}{\sinh^2V}+\frac{2V}{e^{2V}-1}\longrightarrow0.
\]
For $T_R<1$, the required $C_3$ value $W(p\coth p-1)$ is
strictly less than the actual $U=WV$. This excludes $C_3$.
The equilibrium images again exclude $C_4,C_5$.

Since $N_R(0)\to\infty$, the unique cell-zero rotating root has
$z<0$ eventually. Thus $f<0$, $g>0$ and
$\delta=2K+u_a+u_b$, where
$u_b=F(\arcsin(\sin b/\sqrt m)\mid m)$.
The initial short integral $u_a$ stays bounded, whereas $u_b\to-\infty$:
the terminal angle tends to $-\pi/2$ and $m\to1$.
The joint divergence follows by lower bounding the positive integral
on each fixed shorter interval and then letting that interval approach
$\pi/2$. Since $m\to1$, also $K\to\infty$, and $|u_b|<K$ throughout,
so the divergence of $u_b$ is compatible with the phase bounds; as $u_a$
stays bounded, $u_a+u_b<0$ eventually. Hence $\delta<2K$ eventually and
$t=k\delta<2kK<2kp_1=t_{\rm cut}$.
The published synthesis makes this source the unique minimizer.
Any $C_1$ cell-zero source would also be strictly pre-cut and
minimizing, a contradiction. Proposition~\ref{prop:global-winding}
supplies the remaining classes. Reflection gives the upper boundary.
\end{proof}

\section{Exterior directions and the exact oscillating face}\label{sec:exterior}

We first complete the $U=0$ face at all radii.
\begin{proposition}[The $U=0$ normal fiber]\label{prop:Useam-full}
For $U=0$, $V\ne0$, $W\ne0$, every normal source belongs to $C_1$.
There is exactly one source in every integer winding class.
The cell-zero source is strictly pre-cut and uniquely minimizing;
all sources are regular.
\end{proposition}
\begin{proof}
The positive cells were proved in Lemma~\ref{lem:c1cells}.
For cell zero, positive $p=r$ forces $\epsilon=1$ in
\eqref{eq:c1Useam}, hence $s_1=-\operatorname{sgn}W$.
The derivative \eqref{eq:c1Useam-derivative} remains positive when
$n=0$. The scalar tends to zero at $k=0$.
As $k\to1$, $a(k)\to\pi/2$ and $r=F(a(k)\mid k^2)\to\infty$,
whereas $E(a(k)\mid k^2)$ remains bounded. Thus
$q=r-E(a(k)\mid k^2)\to\infty$, proving that the scalar tends
to infinity. It yields one cell-zero source for each $|V|>0$.
Its time satisfies $t<2K=t_{\rm cut}$.

For $C_2$, midpoint coordinates give
\[
 U=s_2\frac{(p-\mathcal E(p))
            \sqrt{1-k^2\sin^2\eta}}{\dn p},\qquad p=t/(2k)>0.
\]
Every factor except the sign is positive, excluding $U=0$.
The $C_3$ formula preceding \eqref{eq:c3surface} also has nonzero
$U$ for $p>0$. The equilibrium images exclude $C_4,C_5$.
The cut-time and nonconjugacy inputs prove the remaining assertions.
\end{proof}

\begin{theorem}[An explicit exterior threshold]\label{thm:exterior}
Fix $\alpha\in(-\pi/2,\pi/2)$ and $0<|\theta|<\pi$ with
$|b_0|=|\alpha-\theta|>\pi/2$, and let
$q_R=(R\ee(\alpha),\theta)$. Put
\[
 h=\theta/2,\quad\beta=\alpha-h,\quad W=\tan h,\quad
 u=\frac{\cos\beta}{2\cos h},\quad v=-\frac{\sin\beta}{2\cos h},
 \quad g=|Wv|-|u|.
\]
Then $g>0$. With
$\eta_0=\min((|b_0|-\pi/2)/2,\pi/4)$, define
\begin{equation}\label{eq:exterior-radius}
 R_{\rm ext}=
 \max\left\{\frac2{\cos\alpha},\,
             \frac2{\sin\eta_0},\,\frac{|W|}{g}\right\}.
\end{equation}
For $R>R_{\rm ext}$ all normal sources of $q_R$ belong to $C_1$,
with exactly one in every integer winding class.
All are regular; the unique zero-winding source is the unique
minimizer. Its sign is $s_1=-\operatorname{sgn}\theta$.
\end{theorem}
\begin{proof}
The negative $C_2$ spatial sign is excluded by $R\cos\alpha>2$.
The positive sign would require
$b=b_0-\arcsin(B/R)$ with $|B|<2$.
Since $|\arcsin(B/R)|<\eta_0<|b_0|-\pi/2$,
this $b$ cannot be principal. Thus there are no $C_2$ sources,
including $B=0$.
The fixed-angle assumptions imply $|b_0|<3\pi/2$ and
$\cos b_0<0$. Direct calculation gives
\[
 u^2-W^2v^2=\frac{\cos\alpha\cos b_0}{4\cos^4h}<0,
\]
so $v\ne0$ and $g>0$.
A $C_3$ source would satisfy \eqref{eq:c3surface}.
Its first equation forces $p>|V|=R|v|$, and
$p\coth p-1>p-1$ then gives
\[
 |U|-|W|(p\coth p-1)<|W|-Rg<0.
\]
The second equation is impossible. The equilibrium images exclude
$C_4,C_5$.

Proposition~\ref{prop:global-winding} already supplies all nonzero
classes. A minimizing trajectory exists by the known synthesis and
is normal in this contact problem \cite{ABB}. It must now belong to $C_1$.
No positive cell can minimize: $p>(2n-1)K\ge K$ there, so
$t>2K=t_{\rm cut}$. Hence a cell-zero source exists.
It has $0<p<K$, is strictly pre-cut, and is the unique minimizer.
Any further cell-zero source would also be a strictly pre-cut
minimizer, contradicting the synthesis. Its sign follows from
$\operatorname{sgn}W=-s_1\operatorname{sgn}r$ and $r=p>0$.
All $C_1$ sources are regular.
\end{proof}

The nonhorizontal part of Theorem~\ref{thm:directions} follows by combining
Theorems~\ref{thm:winding}, \ref{thm:boundary-count},
\ref{thm:exterior}, and Propositions~\ref{prop:boundary-full},
\ref{prop:global-winding}.
The boundary law does not follow by putting $m_*=1$ in the
interior formula: its modulus complement is of order
$R^{-2}\operatorname{sech}^2z$, producing the logarithmic drift
in \eqref{eq:boundary-L}.

\section{Source reversal and horizontal directions}\label{sec:horizontal}

\begin{proposition}[Reversal on the full source space]\label{prop:reversal}
There is a smooth involution $\mathcal S$ of the positive-time unit-energy
source space such that
\[
 \operatorname{Exp}\circ\mathcal S=\iota\circ\operatorname{Exp},
 \qquad
 \iota(X,\theta)=(-\operatorname{Rot}(-\theta)X,-\theta).
\]
It preserves the pendulum stratum, modulus, travel time, regularity,
full-map local singularity type and minimizing status. For a strict
principal heading, it sends winding $\ell$ to $-\ell$.
\end{proposition}
\begin{proof}
This is the established reflection $\varepsilon^5$ of \cite{MS}, in the
following full-source convention. If $g(s)$ is the source trajectory on
$[0,T]$, put $g_{\mathrm r}(s)=g(T)^{-1}g(T-s)$. Its controls are
$-u(T-s)$. With $u_1=\sin(\gamma/2)$, $u_2=-\cos(\gamma/2)$, the vertical
coordinates are
\[
 \gamma_{\mathrm r}(s)=\gamma(T-s)+2\pi\pmod{4\pi},
 \qquad c_{\mathrm r}(s)=-c(T-s).
\]
They satisfy $\dot\gamma_{\mathrm r}=c_{\mathrm r}$ and
$\dot c_{\mathrm r}=-\sin\gamma_{\mathrm r}$, with unchanged pendulum
energy. If $\Phi_T$ is the vertical flow and
$r(\gamma,c)=(\gamma+2\pi,-c)$, then $r\Phi_T=\Phi_{-T}r$ and
$\mathcal S(\lambda,T)=(r\Phi_T(\lambda),T)$ is a smooth involution.
The endpoint identity follows from the path construction. Both outer
maps are diffeomorphisms, including their time dependence, so they
preserve the full-map local type. Reversing and translating competitors
gives a length-preserving bijection of the two endpoint problems, proving
the minimizing assertion. The lifted heading ends at
$-\widetilde\theta(T)$, which gives the strict-heading winding rule.
\end{proof}

\begin{proposition}[Horizontal boundary transport]\label{prop:horizontal}
For $q_R=(\sigma R,0,\theta)$, $\sigma=\pm1$ and $0<|\theta|<\pi$,
the eventual zero-winding fiber has the boundary count
\[
 N_0(q_R)=\frac{R}{2(\log(4R/|\sin\theta|)-1)}+O(1).
\]
Its exact eventual floor formula, regularity and fold statements are
those of Theorem~\ref{thm:boundary-count}, transported by $\mathcal S$.
The unique minimizer is a strictly pre-cut $C_2$ source. The threshold
remains non-effective.
\end{proposition}
\begin{proof}
The inverse target is
$(-\sigma R\cos\theta,\sigma R\sin\theta,-\theta)$. Its spatial
direction is nonhorizontal, with absolute vertical component
$|\sin\theta|$, and its terminal-body transverse component is zero.
It is therefore a boundary target of Section~\ref{sec:boundary}.
Proposition~\ref{prop:reversal} gives the required source bijection,
including winding zero, full-map type and minimizing status.

To complete Theorem~\ref{thm:directions}, put
$\delta_0=e_y$, $\delta_1=e_y\cos\theta-e_x\sin\theta$. After spatial sign
normalization in the nonhorizontal case,
$\mathbf e=\pm\ee(\alpha)$ and
$\delta_0\delta_1=\cos\alpha\cos(\alpha-\theta)$. Thus the sign of this product
is exactly the interior/boundary/exterior condition already proved.
In the interior,
$m_*=1-\min(\delta_0^2,\delta_1^2)$; on either boundary the nonzero transverse
component has magnitude $|\sin\theta|$. The horizontal case is the
transport just established.
\end{proof}

\section{Pure orientation and heading \texorpdfstring{$\pi$}{pi}}\label{sec:pi}

\begin{lemma}[The fiber over pure orientation]\label{lem:orientation}
For $(X,\theta)=(0,\theta)$, with $\theta\in(-\pi,\pi]$, the only
positive-time sources are in $C_4$. There is exactly one for every
$\ell$ with $\Theta_\ell=\theta+2\pi\ell\ne0$, at time
$T=|\Theta_\ell|$ and rotation sign $s_1=-\operatorname{sgn}\Theta_\ell$.
All are regular. A nonzero heading has one minimizer, except at $\pi$,
where the two shortest lifts minimize. At the identity there is no
positive-time winding-zero source and no positive-time minimizing source.
\end{lemma}
\begin{proof}
Let $a=p_x$ and $b=p_y$ be the conserved canonical spatial momenta.
The standard Hamiltonian \cite[Section 2.2]{WKB} and its equations are
\[
 H=\tfrac12[p_\theta^2+(a\cos\theta+b\sin\theta)^2]=\tfrac12,
 \quad \dot X=u(\cos\theta,\sin\theta),\quad
 u=a\cos\theta+b\sin\theta,\quad \dot\theta=p_\theta .
\]
Hence $a\dot x+b\dot y=u^2$. If $X(T)=X(0)$ with $T>0$,
integration forces $u\equiv0$. The energy condition gives
$p_\theta^2=1$, and differentiating $u=0$ gives
$-a\sin\theta+b\cos\theta=0$. These two equations imply $a=b=0$.
Thus the source is a pure $C_4$ rotation, with
$X=0$ and $\widetilde\theta=-s_1T$, proving the source list.
This arbitrary-time argument uses the same conserved-momentum
positivity as the period-drift argument in \cite[Section 2.3]{WKB};
here the identity applies on an arbitrary time interval.
Regularity follows from \cite[Theorem 2.4]{SC}; equivalently its full
determinant in $(\gamma,c,T)$ is $(T^2-\sin^2T)/16>0$.
The cut time is $\pi$ \cite[(2.4)]{SS}. The shortest lifted rotations
give the stated minimizers. At the identity the distance is attained
at zero time, which is outside the source domain being counted.
\end{proof}

\begin{theorem}[Exact heading-$\pi$ catalogue]\label{thm:pi}
At every target $(X,\pi)$ there is exactly one source in each winding
class. For $R=|X|>0$ all sources belong to $C_1$; for $R=0$ they are
the $C_4$ sources of Lemma~\ref{lem:orientation}. All are regular.
For $R>0$ set
\[
 H(k)=\frac{K(k^2)-E(k^2)}k.
\]
For $n\ge0$, let $a=2n+1$ and solve $R=2aH(k)$ for $k\in(0,1)$.
The pair with windings $n,-n-1$ has common time
$T_n(R)=2aK(k^2)$ and phases determined by
\[
 \sn\phi=s_1x/R,\quad \cn\phi=-y/R,\qquad
 \ell=\begin{cases}n,&s_1=-1,\\-n-1,&s_1=1.\end{cases}
\]
The phase is unique modulo $4K$. Exactly the first pair minimizes.
Different pairs have different travel times.
\end{theorem}
\begin{proof}
The $C_1$ lifted heading is
$s_1(\operatorname{am}\phi-\operatorname{am}(\phi+T))$.
Strict increase of the real amplitude forces
$T=2(2n+1)K$ at group heading $\pi$. The half-period identities in
the original spatial equations then give
\[
 X=2aH(k)(s_1\sn\phi,-\cn\phi).
\]
Every $C_2$ or $C_3$ lifted heading lies strictly in $(-\pi,\pi)$;
$C_5$ has heading zero. Thus these are all sources when $R>0$.
Moreover
\[
 H(k)=k\int_0^{\pi/2}\frac{\sin^2u}{\sqrt{1-k^2\sin^2u}}\,du,
 \qquad
 H'(k)=\int_0^{\pi/2}
       \frac{\sin^2u}{(1-k^2\sin^2u)^{3/2}}\,du>0.
\]
The endpoint limits are zero and infinity, so the modulus and the full
phase-circle point are unique. Opposite $s_1$ give distinct covectors,
since $\cos(\gamma/2)=s_1\dn\phi$ never vanishes.
The lifted heading is $-s_1a\pi$, giving the displayed labels.
Nonconjugacy follows from \cite[Theorem 2.1]{SC}, including these
half-period boundaries. The cut time $2K$ \cite[(2.1)]{SS} selects
exactly $a=1$. Time ordering is proved next.
\end{proof}

\begin{proposition}[Heading-pi time spectrum]\label{prop:pi-times}
For $R>0$ the pair times in Theorem~\ref{thm:pi} satisfy
\[
 R+2a<T_n(R)<R+a\pi,\quad a=2n+1,\qquad
 4<T_{n+1}(R)-T_n(R)<2\pi.
\]
At fixed $R$ the gaps increase strictly with $n$ and tend to $2\pi$.
At fixed $n$ they decrease strictly with $R$ and tend to $4$.
For fixed $R$, as $a=2n+1\to\infty$,
\[
 T_n(R)=a\pi+\frac{R^2}{a\pi}
           -\frac{3R^4}{4a^3\pi^3}+O_R(a^{-5}).
\]
\end{proposition}
\begin{proof}
Allow $a>0$ to be real in $R=2aH(k)$, $T=2aK$. The complete-integral
derivatives give $K_k=kH'$ and $E_k=-H$, hence
\[
 T_R=k,\qquad T_a=2E,\qquad
 T_{aa}=\frac{2H^2}{aH'}>0.
\]
Integrating $2<2E<\pi$ over an $a$ interval of length two proves the
gap bounds and convexity proves their strict increase. Since $k$ decreases
with $a$, the radial derivative of a gap is negative. The endpoint limits
follow by dominated convergence in that same integral for $T_a$.
The quantity $K-H$ decreases strictly from $\pi/2$ to $1$, since
$(K-H)'=(k-1)H'<0$; this proves the time-offset bounds.
Finally, with $z=2R/(a\pi)$, the convergent small-modulus series
\cite[Section 19.5]{DLMF} give
$k=z-3z^3/8+3z^5/16+O(z^7)$ and
$T/(a\pi)=1+z^2/4-3z^4/64+O(z^6)$.
\end{proof}

\section{Zero heading: exact fibers and exceptional axes}\label{sec:zero}

In this section write $\mathcal D(k)=K(k^2)-E(k^2)$; unadorned $K,E$
at a source mean $K(k^2),E(k^2)$. Source phases still have period $4K$.

\begin{theorem}[Zero-heading winding and axes]\label{thm:zero-axes}
For every $X\ne0$ at heading zero, each nonzero winding class has exactly
one regular $C_1$ source. For $n=|\ell|\ge1$ it is determined by
\[
 R=4nH(k),\quad T=4nK,\quad s_1=-\operatorname{sgn}\ell,\quad
 \sn\phi=s_1x/R,\quad \cn\phi=-y/R .
\]
All these sources are post-cut. There are no $C_3$ sources.
For $(x,0,0)$, $x\ne0$, winding zero consists of the unique straight
$C_5$ source, which is regular and uniquely minimizing.

For $(0,y,0)$, $y\ne0$, winding zero consists of four regular $C_2$
sources for each positive root label $n$: two with $p=2nK$ and two with
$p=p_n^1\in((2n-1)K,2nK)$, where $f_1(p_n^1)=0$.
The central pair has actual period index $n$, whereas this $f_1$ pair
has index $n-1$. Thus actual cell zero contains two sources; each
positive actual cell $j$ contains the central pair labelled $j$ and
the $f_1$ pair labelled $j+1$.
Exactly the first $f_1$ pair minimizes. For this fixed vertical target,
the total number of sources with travel time at most $t$ satisfies
\begin{equation}\label{eq:zero-vertical-time}
 N_{\le t}(0,y,0)=\frac{t^2}{\pi|y|}+\frac t\pi+O_y(1).
\end{equation}
\end{theorem}
\begin{proof}
The continuous $C_1$ amplitude now forces $T=4nK$, and substitution in
the original map gives the displayed radial equation. The proof of
Theorem~\ref{thm:pi} applies with $a=2n$. The $C_3$ lifted heading is
$s_1[\arcsin(\tanh\phi)-\arcsin(\tanh(\phi+T))]$; it is nonzero and
strictly between $-\pi$ and $\pi$. Thus it cannot have group heading zero.
Every $C_2$ source here has winding zero. In the chart
\eqref{eq:UVW}, at this heading $U=-y/2$, $V=-x/2$, $W=0$.
Writing $\eta=\operatorname{am}(\psi+p)$ and $q=p-\mathcal E(p)>0$,
the original map gives
\begin{equation}\label{eq:zero-full-chart}
 U=\frac{s_2q\sqrt{1-k^2\sin^2\eta}}{\dn p},\quad
 V=\frac{s_2k f_1(p)\sin\eta}{\dn p},\quad
 W=-\frac{k\sn p\cos\eta}{\dn p}.
\end{equation}
In particular $U\ne0$, excluding $C_2$ on the horizontal axis.
The remaining straight $C_5$ source minimizes since $T\ge|x|$, with
equality forcing that straight control. Its regularity follows from
\cite[Theorem 2.4]{SC}.

For the vertical axis put $u=|y|/2$ and $s_2=-\operatorname{sgn}y$.
The two zero factors in \eqref{eq:zero-full-chart} give precisely
$p=2nK$, $\eta=0,\pi$, or $f_1(p)=0$, $\eta=\pi/2,3\pi/2$.
The function $f_1/\cn p=-q-\dn p\,\sn p/\cn p$ has derivative
$-\dn^2p/\cn^2p<0$ between poles. Thus it has exactly one zero in
each interval $((2n-1)K,2nK)$ and none in the preceding half-cell.
These are the root intervals of \cite[Lemma 5.3]{MS}. We take this
occasion to correct equation (5.15) there: the continuous $k\to0$ limit
of the $n$-th root is $n\pi$, not the displayed $2\pi n$, as
$f_1(p,0)=-\sin p$ and the adjacent first-root limit show.
The central modulus solves $2n\mathcal D(k)=u$ uniquely.
For the $f_1$ pair put $h=\sqrt{1-k^2}$, $\chi=\arctan(u/h)$.
Its equations are
\[
 p=2nK-F(\chi\mid k^2),\quad
 2n\mathcal D(k)=I(u,h),\quad
 I(u,h)=\int_0^u\sqrt{\frac{1+z^2}{h^2+z^2}}\,dz .
\]
The last integral equals
$F(\chi)-E(\chi)+\sqrt{1-k^2\sin^2\chi}\tan\chi$, by differentiation
in $u$. At fixed $u$, substituting $z=h\tan v$ in its modulus derivative
gives
\[
 \frac{d}{dk}[2n\mathcal D(k)-I(u,h)]
  =\frac{k}{h^2}[2nE-E(\chi\mid k^2)]>0.
\]
The endpoint values are $-u$ at $k=0$ and infinity at $k=1$:
$I(u,h)=\log(1/h)+O_u(1)$ and
$\mathcal D(k)=\log(4/h)-1+o(1)$. This proves unique reconstruction
of both phase choices for every $n$. These four sources are exhaustive.
Sachkov's reduced Jacobian \cite[(2.10)--(2.14)]{SC} (see also
\cite{CT}), whose full-map prefactor is nonzero here, is
$J_2=\alpha\sn^2\tau+\beta\cn^2\tau$ with $\tau=\psi+p$,
\begin{align*}
 \alpha&=(1-k^2)\sn p\,\{\cn p\dn p\,(p-2\mathcal E(p))
        +\sn p\,(\dn^2p+\mathcal E(p)q)\},\\
 \beta&=f_1(p)(\cn p\,\mathcal E(p)-\dn p\sn p),
\end{align*}
the brace being the $\alpha_1$ of Lemma~\ref{lem:zero-order} below.
At the central pair $\sn p=\sn\tau=0$, $\cn p=(-1)^n$ and
$f_1(p)=(-1)^{n+1}q$, so $J_2=\beta=-2nEq<0$. At the $f_1$ pair
$\beta=0$ and $\cn\tau=0$; substituting $\sn p=-q\cn p/\dn p$ from
$f_1(p)=0$ turns the brace into
$-\cn p\,\mathcal E(p)(\dn^2p+q^2)/\dn p$, and
$\cn^2p\,(\dn^2p+q^2)=\dn^2p$, so $J_2=\alpha=(1-k^2)\mathcal E(p)q>0$.
Thus all are regular. The cut time $2kp_1^1$ \cite[(2.2)]{SS}
selects exactly the first $f_1$ pair; all remaining sources are post-cut.

For fixed $u$, both sequences have $k_n\to0$, since their increasing
scalar left sides diverge with $n$ at every fixed positive modulus.
Using $\mathcal D(k)=\pi k^2/4+O(k^4)$ and
$I(u,\sqrt{1-k^2})=u+O_u(k^2)$ yields
\[
 k_n^2=\frac{2u}{\pi n}+O_u(n^{-2}),\qquad
 T_n^2=8\pi u n+O_u(1)
\]
in both families; subtracting $2kF(\chi)$ in the $f_1$ time changes
its square by $O_u(1)$. The two families, with two phases each, therefore contribute
$t^2/(2\pi u)+O_u(1)$ sources in total. The $C_1$ pair times are
$2\pi n+O_R(n^{-1})$, contributing $t/\pi+O_R(1)$. Their sum proves
\eqref{eq:zero-vertical-time}, without any ordering assumption between
the two $C_2$ time sequences.
\end{proof}

\begin{theorem}[Generic zero-heading count]\label{thm:zero-count}
Fix $xy\ne0$, put $u=|y|/2$, $v=|x|/2$, $\rho=\sqrt{u^2+v^2}$,
$\nu=v/\rho$, and $S=\rho/[2\mathcal D(\nu)]$. Define
\begin{align}
 k_{\min}^2&=\frac{2v^2}{z+\sqrt{z^2-4v^2}},
       \qquad z=1+u^2+v^2,\label{eq:zero-kmin}\\
 \phi&=\arctan(u/h),\qquad h=\sqrt{1-k^2},\nonumber\\
 a&=\arcsin\frac{vh}{k\sqrt{u^2+h^2}},\qquad
 A_L=a-\phi,\quad A_R=\pi-a-\phi,\nonumber\\
 N_j(k)&=\frac{u d_j/h-F(A_j\mid k^2)+E(A_j\mid k^2)}
                    {2\mathcal D(k)},
 \quad d_j=\sqrt{1-k^2\sin^2 A_j}.\label{eq:zero-index}
\end{align}
Let $M=\sup_{k_{\min}\le k<1}N_L(k)$. The winding-zero fiber consists
of finitely many $C_2$ sources, with exact count
\begin{equation}\label{eq:zero-count}
 N_0(x,y,0)=1+2\lfloor S\rfloor+2\lfloor M\rfloor
                   -2\ind_{\{S\in\mathbb Z\}}-\ind_{\{M\in\mathbb Z\}}.
\end{equation}
The unique complementary cell-zero source minimizes; every other
zero-winding source is post-cut.
\end{theorem}

We prove the count with the distinction between an inverse-formula
branch and a connected fixed-label source branch kept explicit.
In this section, a \emph{complementary cell} means the latter branch;
its label $n$ is not generally the actual period index
$\lfloor p/(2K)\rfloor$. In particular, the complementary cell-zero
interval is $0<p<p_1^1$, a proper subset of actual period cell zero.

\begin{lemma}[Complementary scalar and physical cells]\label{lem:zero-scalar}
The sources on $\cos\eta=0$ are in bijection with
$N_L(k)=n$ or $N_R(k)=n$, $n\in\mathbb Z_{\ge0}$,
subject to $p=2nK+F(A_j\mid k^2)>0$ and one gluing identification
at $k=k_{\min}$. Both index branches are bounded, with limits
$N_L\to1/2$, $N_R\to-1/2$ as $k\to1$.
\end{lemma}
\begin{proof}
In a sign-normalized cell write $s=\sin A$, $c=\cos A$,
$d=\sqrt{1-k^2s^2}$ and
\begin{equation}\label{eq:zero-physical}
 \begin{gathered}
 q=2n\mathcal D(k)+F(A\mid k^2)-E(A\mid k^2),\\
 e=2nE+E(A\mid k^2),\qquad
 u=hq/d,\quad v=k(s+qc/d).
 \end{gathered}
\end{equation}
Eliminating $q$ gives $\sin A+(u/h)\cos A=v/k$, whose two
representatives are exactly \eqref{eq:zero-index}. The existence
inequality is $k^4-(1+u^2+v^2)k^2+v^2\le0$, giving
\eqref{eq:zero-kmin}. Conversely these equations recover both target
components. The actual phase is fixed by
$s_2=\operatorname{sgn}U$ and
$\sin\eta=\operatorname{sgn}(UV)(-1)^{n+1}$, with
$\psi=F(\eta\mid k^2)-p$. A simultaneous half-period shift changes the
sign-normalized representative and its integer label, not the source.
For $n\ge1$ the physical positive-$v$ interval runs from an $f_1$ root
in $(-\pi/2,0)$ to one in $(\pi/2,\pi)$; for $n=0$ it begins at $A=0$.

At $k\to1$, $A_L=-\pi/2+O(h)$ and $A_R=\pi/2+O(h)$.
The quantities $u d_j/h$ remain bounded, whereas the incomplete
drifts are respectively $-\mathcal D(k)+O(1)$ and
$\mathcal D(k)+O(1)$. The error is bounded by rescaling the $O(h)$
endpoint layer by $h$ in its integral. This proves the limits and,
together with continuity at $k_{\min}$, boundedness.
\end{proof}

The real-label interpolation below is an auxiliary analytic family.
For noninteger $n$, $q,e$ are defined by \eqref{eq:zero-physical};
they need not equal $p-\mathcal E(p),\mathcal E(p)$ at that value of $p$.
Write $\widehat f=-cq-ds$ for its boundary factor. At integer labels
its zeros are the physical $f_1$ zeros.

\begin{lemma}[Strict ordering for real cell labels]\label{lem:zero-order}
Allow $n\ge1$ to be real in \eqref{eq:zero-physical}. In its positive-$v$
source interval there is exactly one root of
\[
 \alpha_1=cd(q-e)+s(d^2+eq).
\]
At this root $A=A_n(k)$ the physical target $v_n(k)$ increases strictly
from zero to infinity as $k$ runs from zero to one. For fixed $v>0$,
the corresponding value $u_n^*(v)$ increases strictly with $n$.
On the connected physical curve $v=\text{constant}$ this is the unique
nondegenerate minimum of $u$.
\end{lemma}
\begin{proof}
All derivatives below keep the real label fixed unless indicated.
Direct differentiation gives
\begin{equation}\label{eq:zero-restricted}
 \det\frac{\partial(u,v)}{\partial(A,k)}
       =\frac{k^2\alpha_1}{h d^2},\qquad
 v_A=\frac{k(cd-h^2qs)}{d^3}.
\end{equation}
The latter is positive for negative amplitude, negative after $\pi/2$,
and has exactly one zero in $(0,\pi/2)$, where $q\tan A/d=1/h^2$.
At that fixed-modulus maximum,
$d v_{\max}/dk=(d^2/s+k^2ce/d)/h^2>0$.
For fixed $n$, its limits are zero and infinity; the latter follows
already by taking $A=0$. Hence a fixed positive $v$ defines one connected
source curve, with two amplitudes above a unique coalescence modulus.

Put $D_\alpha=cd+es$ and $\beta_1=ce-ds$. At an $\alpha_1$ root,
\begin{equation}\label{eq:zero-root}
 qD_\alpha=d\beta_1,\quad
 \widehat f=-de/D_\alpha,\quad
 \alpha_{1,A}=\frac{e^2}{dD_\alpha},\quad
 cd-qs=\frac{d^2}{D_\alpha}.
\end{equation}
In the negative quarter $\beta_1>0$, so any root has $D_\alpha>0$
and crosses upward. In the positive quarter $\beta_1$ decreases
strictly, since $\beta_{1,A}=-s(e-k^2sc/d)<0$; a root can only
occur before its zero and again crosses upward. After $\pi/2$ any
root would cross downward, but $\alpha_1$ is positive at both ends
of that last physical quarter. The $\widehat f$ endpoint signs follow from
$s\alpha_1=eq$ there. These signs prove uniqueness.
Since $\alpha_1(0)=2n(K-2E)$, the root has positive amplitude below
the unique modulus $k_0$ solving $K(k_0^2)=2E(k_0^2)$, negative
amplitude above it, and zero amplitude at it. The modulus $k_0$ is
unique because $(K-2E)'=K_k+2H>0$ with opposite endpoint signs.
The root is smooth, with $D_\alpha>0$ throughout.

Set $B_\alpha=d^3+k^2ces$. This is positive: for $s<0$,
$D_\alpha>0$ implies $e|s|<cd$, whence
$k^2ce|s|<k^2c^2d<d^3$. At the root,
\[
 v_A=\frac{kB_\alpha}{d^2D_\alpha}>0.
\]
For the half-time $\tau_n=k(e+q)$, the exact derivatives
$q_k=k(e-sc/d)/h^2$, $e_k=-q/k$, $q_A=k^2s^2/d$, $e_A=d$
give
\begin{equation}\label{eq:zero-time-bridge}
 \begin{gathered}
 \tau_n'=\frac{dI}{h^2e^2D_\alpha},\qquad
 v_n'=\frac{B_\alpha}{dD_\alpha}\tau_n',\\
 I=ds\{(e^2-d^2)^2+2e^2(1-2k^2c^2)\}
       +ce(1-4k^2s^2)(e^2-d^2).
 \end{gathered}
\end{equation}
We give the sign argument, including noninteger $n$. The closed forms
below, namely $F_R$ and its two boundary values, $P$, the Bernstein
coefficients, and the factorization of $I$ through $G$, are polynomial
identities in the substituted variables; they were verified by computer
algebra in the script \path{research/e33_zero_heading_count.py} of the
repository, and the substitutions are collected in
Appendix~\ref{app:subs}.

For positive amplitude put $z=d/e$ and $\tan A=rz$. Then
$e>2$, $d^2>\max(4z^2,c^2)$, $rz^2<1$, and
\[
 I=\frac{d^3F_R}{rz^3\sqrt{1+r^2z^2}},\qquad
 F_R=d^2\{2+r(1-z^2)\}^2+2r^2z^2+3rz^2-3r-4.
\]
This expression increases with $d^2$. At the active lower boundary
$d^2=c^2$, its value is
\[
 \frac{r(r+1)(rz^2-1)\{(2r+1)z^2-1\}}{1+r^2z^2}>0,
\]
because $4z^2\le c^2$ makes both linear factors negative.
At the other active boundary put $w=z^2$, $\xi=rw\in(0,1)$.
Multiplication by $w$ gives
\[
 P=4w^2(\xi-2)^2+w(-8\xi^2+19\xi-4)+3\xi(2\xi-1).
\]
It is positive termwise for $\xi\ge1/2$. For $\xi<1/2$, the active
condition is $w\ge1/4-\xi^2$, and
\[
 P(1/4-\xi^2,\xi)=\frac{\xi(\xi-1)(2\xi-3)(2\xi-1)(4\xi^2-4\xi-1)}4>0.
\]
Moreover $P_{ww}>0$. At the lower endpoint, the Bernstein coefficients
of $P_w$ on $0\le2\xi\le1$ are
\[
 (4,\ 43/8,\ 31/6,\ 35/8,\ 7/2).
\]
All are positive, completing the subcritical sign proof.

For negative amplitude put $\chi=-A$, $z=d/e$ and $\tan\chi=rz$.
The root relation gives $0<r<1$ and $q/e=(1+rz^2)/(1-r)>1$.
We first establish $e>2n-1/2$. The complementary expression
$D_0=e\sin\chi-d\cos\chi$ is negative at the root and has derivative
$D_0'=\cos\chi(e+k^2\sin\chi\cos\chi/d)>0$.
At $\chi_0=\pi/6$ it is positive already for $n=1$: if
$w(\vartheta)=\sqrt{1-k^2\sin^2\vartheta}$ and $g=w-\cos\vartheta$, then
$g$ increases, since
$w^2-k^4\cos^2\vartheta=(1-k^2)(1+k^2\cos^2\vartheta)>0$.
Subtracting the equality at $k=1$ gives
\[
 2E-E(\chi_0)-\sqrt3w(\chi_0)
 =\int_0^{\chi_0}g+2\int_{\chi_0}^{\pi/2}g-\sqrt3g(\chi_0)>0,
\]
using $2\pi/3>\sqrt3$. Adding $2(n-1)E\sin\chi$ proves the
same comparison for every real $n\ge1$. Thus $\chi<\pi/6$ and
\[
 e>2(n-1)E+2E-E(\pi/6)
   >2(n-1)+\sqrt3\sqrt{1-k^2/4}>2n-\tfrac12.
\]

The paired integrands satisfy
$w(\vartheta)w(\pi/2-\vartheta)>h$, hence
$w(\vartheta)+w(\pi/2-\vartheta)>
h\{1/w(\vartheta)+1/w(\pi/2-\vartheta)\}$. Integrating gives $E>hK$,
so $K-2E<0$ at $k^2=3/4$ and $k_0^2>3/4$.
Above $k_0$, $h^2<1/4$ and $2h^2\mathcal D(k)<3\pi/16<3/4$;
the latter follows from $\mathcal D(k)\le\pi k^2/(4h)$ and
$h(1-h^2)<3/8$. Consequently
\begin{equation}\label{eq:zero-margin}
 e(e^2-1)-h^2\{2n\mathcal D(k)+3e\}
 >\frac{(n-1)(32n^2+8n-3)}4\ge0 .
\end{equation}
The strict lower bound is obtained by substituting $e>2n-1/2$
in the increasing cubic $e(e^2-7/4)-3n/4$.

To connect this margin to $I$, set
\[
 \begin{split}
 G={}&(1-r)(1+rz^2)\{1+(2r-1)z^2\}\\
     &-h^2rz^2\{2-r(1-z^2)\}^2,\\
 B_0={}&eq(e^2-d^2)-h^2d^2(e+q)^2.
 \end{split}
\]
Exact substitution gives $I=d^3G/[z^3(1+r^2z^2)^{3/2}]$ and
\[
 G-\frac{(1-r)^2B_0}{e^4}
 =(1-r)z^2\{2r(1+rz^2)+h^2[2-r(1-z^2)]^2\}>0 .
\]
Here $q>e$, $q<2n\mathcal D(k)$ and
$(e+q)^2/q<2n\mathcal D(k)+3e$, so \eqref{eq:zero-margin}
implies $B_0>0$. This proves $I>0$ supercritically.
At $k=k_0$, directly $\tau_n'=(e^2-1)/(h^2e)>0$.
Thus \eqref{eq:zero-time-bridge} proves $v_n'>0$ everywhere.
Its limits are zero and infinity: at the upper end
$q>(2n-1)\mathcal D(k)$ and \eqref{eq:zero-root} imply
$v_n\ge kq/2\to\infty$.

For fixed $v$, the unique alpha intersection is therefore smooth in $n$.
Along its increasing-amplitude branch,
\[
 \left.\frac{du}{dk}\right|_v=-\frac{k^2\alpha_1}{hd^2v_A},
 \qquad
 \left.\frac{d\alpha_1}{dk}\right|_v
       =-\frac{\alpha_{1,A}v_n'}{v_A}<0
 \quad\hbox{at the root}.
\]
It is a nondegenerate minimum, with
\begin{equation}\label{eq:zero-curvature}
 u''=\frac{k^2\alpha_{1,A}v_n'}{hd^2v_A^2}>0 .
\end{equation}
Both ends of the connected fixed-$v$ curve have $u\to\infty$:
$q\ge(2n-1)\mathcal D(k)\to\infty$ and bounded $v$ force
$c/d\to0$, hence $h/d\to1$. There is no other critical point.
Finally the envelope derivative with respect to the real label is
\begin{equation}\label{eq:zero-label}
 \frac{\partial u_n^*(v)}{\partial n}
      =\frac{2\mathcal D(k)h d^2}{B_\alpha}>0 .
\end{equation}
This proves strict ordering without interpolating integer samples.
\end{proof}

\begin{proof}[Proof of Theorem~\ref{thm:zero-count}]
Each positive complementary cell has two, one or zero sources according
as $u>u_n^*(v)$, $u=u_n^*(v)$ or $u<u_n^*(v)$, by the unique
minimum in Lemma~\ref{lem:zero-order}. These thresholds diverge with $n$:
if $u$ stayed bounded with $v$ fixed, phase existence would bound $k$
away from zero, while $q\ge(2n-1)\mathcal D(k)$ would diverge,
contradicting $|q-\rho|\le1$ from \eqref{eq:zero-physical}.

Cramer's rule for the inverse branches gives
\[
 N_j'=\frac{k\alpha_1}{2\mathcal D(k)h^2(cd-qs)}.
\]
The denominator is positive on $L$ and negative on $R$. Every alpha
root for a real label at least one has $cd-qs=d^2/D_\alpha>0$;
therefore any stationary value at height at least one lies on $L$.
By \eqref{eq:zero-label} there is at most one such real-label fold for
fixed $(u,v)$, and \eqref{eq:zero-curvature} makes it a strict maximum.
At the inverse gluing endpoint $cd-qs=0$ and
$\alpha_1=d^2/s>0$, so $N_L$ initially increases. Its limit $1/2$
then shows that any maximum of height at least one is interior.
It follows that the positive-cell count is two/one/zero according
as $n<M$, $n=M$, $n>M$.

The physical cell-zero interval is $0<p<p_1^1$; all its sources are
strictly pre-cut. Existence of a minimizer, the other-regime exclusions
in Theorem~\ref{thm:zero-axes}, and the known cut locus \cite{SS}
give exactly one such source: at heading zero with $x\ne0$ the target
is off the cut locus. Thus the complementary count is
$1+2\lfloor M\rfloor-\ind_{\{M\in\mathbb Z\}}$.

At a central source $p=2nK$, \eqref{eq:zero-full-chart} gives
\[
 q=2n\mathcal D(k),\quad
 U=s_2q\sqrt{1-k^2\sin^2\eta},\quad
 V=s_2(-1)^{n+1}kq\sin\eta .
\]
Therefore $q=\rho$, with a unique modulus and two/one/zero phases
according as $k>\nu$, $k=\nu$, $k<\nu$.
This is the two/one/zero test $n<S$, $n=S$, $n>S$, giving central
count $2\lfloor S\rfloor-\ind_{\{S\in\mathbb Z\}}$.
The source strata meet only at $\sn p=\cos\eta=0$, where $k=\nu$ and
$N_L=S$. If $S$ is an integer, exactly one source is counted in both
lists. Subtracting it proves \eqref{eq:zero-count}.
\end{proof}

\begin{proposition}[Rank at the critical overlap]\label{prop:zero-rank}
The complementary equality in Theorem~\ref{thm:zero-count} is a fold
of its restricted two-dimensional map for every modulus. For the full
map it is a fold except when $k=k_0$ and $p=2nK$, where the corank
is two. A central coalescence is an ordinary cusp off $k_0$ and has
corank two at $k_0$. At the critical target in cell $n$, $S=M=n$ and the exact
$C_2$ count is $4n-2$.
\end{proposition}
\begin{proof}
Equation \eqref{eq:zero-curvature} gives the restricted fold, including
$k_0$. Away from $\sn p=0$, the derivative $W_\eta$ in
\eqref{eq:zero-full-chart} is nonzero, so solving for $\eta$ transports
this fold to the full map. At a central source use
$p=2nK+F(A\mid k^2)$, $\eta=\pi/2$, $s_2=1$ and $A=0$.
With $q=2n\mathcal D(k)$, $e=2nE$, $q_k=2nkE/h^2$, the full
derivative in $(A,k,\eta)$ is
\[
 \begin{pmatrix}
 0&k(e-q)/h&0\\
 (-1)^{n+1}k&(-1)^{n+1}(q+kq_k)&0\\
 0&0&0
 \end{pmatrix}.
\]
The source and target coordinates are nonsingular here. Its rank is two
unless $q=e$, when it is one. That equality is precisely $K=2E$.
Other signs give the same ranks by reflection.

For $q\ne e$, put $b=\eta-\pi/2$ and fix the first two target
coordinates $U,V$ at their central values. The implicit function theorem
gives even analytic functions $A(b),k(b)$ with $A(0)=0$.
Differentiating $U,V$ twice at $b=0$ gives
\[
 k''(0)=-\frac{qk}{e-q},\qquad
 A''(0)=\frac{qe}{h^2(e-q)}.
\]
The third target coordinate therefore has the reduced expansion
\[
 W(b)=\frac{(-1)^n kqe}{2h^2(e-q)}b^3+O(b^5).
\]
Its cubic coefficient is nonzero. Allowing $V$ to vary with $U$ fixed,
the same reduction gives $\partial_V\partial_bW=-1$ at the source:
$A_V=1/[(-1)^{n+1}k]$ and $W_b=(-1)^nkA+O(A^2)$.
These are the nonzero cubic and unfolding derivative of the ordinary
cusp criterion. At $k_0$ this reduction is unavailable; only the
corank-two assertion is made here. At the critical target, the central
equality gives $S=n$, and the restricted fold gives $M=n$.
Substitution in \eqref{eq:zero-count} proves the final count.
\end{proof}

\begin{proposition}[Zero-heading far-field law]\label{prop:zero-farfield}
For a fixed non-axis unit spatial direction $\mathbf e$, put
$\nu=|e_x|\in(0,1)$. Then
\[
 N_0(R\mathbf e,0)
   =\frac{R}{K(\nu^2)-E(\nu^2)}+O_{\mathbf e}(1)
   =\frac{4R}{L(\nu^2)}+O_{\mathbf e}(1).
\]
The statement is not uniform toward the coordinate axes.
\end{proposition}
\begin{proof}
In \eqref{eq:zero-physical}, $(u,v)$ is $q$ times a unit vector
plus $(0,k\sin A)$, giving $|q-\rho|\le1$. Since the incomplete
drift has absolute value at most $2\mathcal D(k)$ on the chosen
representatives, and $N_L(\nu)=S$,
\[
 S\le M\le\frac{\rho+1}{2\mathcal D(k_{\min})}+1.
\]
For $r=k_{\min}^2$, its quadratic equation is equivalently
$\rho^2(\nu^2-r)=r(1-r)$, so
$0<\nu^2-r\le1/(4\rho^2)$.
At fixed $\nu\in(0,1)$,
$\mathcal D(k_{\min})=\mathcal D(\nu)+O_\nu(\rho^{-2})$,
bounded away from zero for large $\rho$. Thus $M=S+O_\nu(1)$.
Equation \eqref{eq:zero-count} with $\rho=R/2$ proves the result.
The exact vertical and horizontal exceptions are in
Theorem~\ref{thm:zero-axes}.
\end{proof}

\section{Numerical corroboration and open questions}
\label{sec:corroboration}

\paragraph{Numerical corroboration.}
The repository records symbolic checks of selected identities and
numerical or interval checks of particular sources and targets. These
include the exact inverse of Proposition~\ref{prop:inverse} against the
endpoint map at 180 to 360 digits on several hundred sources; twelve integrations of
the Hamiltonian system without elliptic functions; and, for two targets,
the complete rotating fiber certified by interval arithmetic with
analytic tail exclusions, 37 sources at $(39/5,-52/5,3/10)$
(Figures~\ref{fig:level} and \ref{fig:geodesics}) and 301 sources at
$(60,-80,3/10)$. The general analytic count and image results do not
depend on these two target computations. The
scripts are in the repository
\url{https://github.com/moiseevigor/se2-conjugate-locus}.
Further finite checks accompany the all-radius ray count and its
small-heading asymptotic, and the zero-heading central cusp reduction.
They use independent defining-integral inversions, original-map
evaluations and selected Hamiltonian integrations. Their recorded
sample sizes and precision do not constitute a full-domain interval
certificate. The finite radial interval in
Figure~\ref{fig:radial-window} has separate outward hyperbolic-root
brackets and independent original-flow checks. Its existence boundaries
are certified; its interior count follows from the analytic proof of
Corollary~\ref{cor:exterior-sector-count}. The wider exterior witness in
Proposition~\ref{prop:exterior-three} has a separate complete 256-bit
count certificate with a full 384-bit alternate replay. Its pointwise
count depends on those interval bounds and the analytic period exclusions.
Exact rational arithmetic separately checks its domain and modulus
polynomials. Eight original Hamiltonian integrations and 110-digit
original-map checks corroborate the sources; these finite checks alone
do not establish completeness. The angular-threshold theorem is analytic;
its recorded finite checks include a separate parameter substitution,
defining-integral evaluations and 32 original Hamiltonian integrations
covering sign choices and several actual periods. The narrow numerical
enclosure in Corollary~\ref{cor:angular-threshold-certificate}, however,
depends on its complete interval cover and alternate replay. An independent
exact rational sweep rechecks coverage and rejects missing/duplicated-box
controls; symbolic identities check the explicit replay derivatives.
The separate localization certificate in
Theorem~\ref{thm:threshold-fold-window} excludes $7465$ boxes at the upper
comparison level and retains $1343$ inside a small rectangle. All exclusions
pass an alternate 384-bit replay. A parentwise rational area and overlap
audit checks coverage independently of the sweep. Uniform Hessian and rank
bounds, independently recalculated by symbolic differentiation in the
original elliptic coordinates, prove a unique ordinary fold at this first
threshold and its local radial count window. Both implementations share Arb;
Corollary~\ref{cor:effective-fold-window} extends the count window to
an explicit slope cap. Its larger-domain certificate has $3111$ global
exclusions and $2048$ local boxes; alternate implementations replay them
in $3203$ and $2080$ nodes respectively. Independent edge checks prove
strict boundary exclusion, and fresh root inclusions verify both cap
endpoints. These computations supply the quantitative extent and the
endpoint enclosures; the connected-graph and convex-image arguments
establish completeness throughout that range. The time-order theorem
uses a separate cover of $117$ positive-momentum boxes and $38$ slope
exclusions. A full alternate calculation in the original elliptic
coordinates checks all $155$ boxes in $287$ nodes at 384 bits. Exact
rational identities verify the simplified derivatives; sixteen original
Hamiltonian integrations corroborate the signed reconstruction. These
finite integrations do not replace the interval signs or the positive
integral argument for the universal time comparison. The separate
cross-period example has $256$ parameter intervals and $257$ joint and
endpoint boxes. A fresh 448-bit forward-derivative implementation replays
all $513$ root inclusions, and exact checks verify the complete label chain
and joins. Five independent polynomial identities check the differential
and physical source-chart formulas. The target is now fully enumerated:
analytic endpoint elimination gives two exhaustive arcs. The first has
$1062$ excluded intervals, replayed independently in $1274$ nodes; the
second has $1953$, replayed with the original endpoint residual in
$11841$ nodes. Exact polynomial tail bounds exclude all other periods
and boundary arrivals. Nine unique simple roots give the full rotating
count and disjoint clock intervals give the time order. Nine original
Hamiltonian integrations corroborate the signed reconstruction and actual
period labels; they do not replace the exhaustive interval proof.
The fixed-line count in Proposition~\ref{prop:nine-source-interval}
uses the additional monotone-modulus argument and a complete critical-source
cover: $350$ higher-period polynomial intervals, four endpoint tails and
$285$ retained-family intervals. All have separate alternate replays; the
last uses $1768$ interval nodes. The two larger fold boxes are independently
included and their intrinsic quadratic terms are nonzero. Properness and
the exact count at $u=6$ then give the full $7/8/9/10/11$ pattern.
The within-family comparison in Theorem~\ref{thm:general-pair-order}
is analytical: defining-integral monotonicity, first variation and horizontal
slicing prove it for general $V>1,W>0$. Separate exact identities, $16$
roots refined to $90$ digits and four direct-integral checks guard the
reductions. A finite scan of $543$ detected pairs guided the proof but
provides neither root completeness nor a universal sign certificate.
For Corollary~\ref{thm:line-time-order}, $41$ rational inter-family
comparisons retain the fixed-line separation. The earlier $14$-interval
sibling certificate, $38$ root enclosures, two fold corridors and three
clock-gap anchors remain independently replayed corroboration. The new
checks and proof audit are recorded under \nolinkurl{references/e130/}.
The dilation results in Section~\ref{sec:dilation} have analytical
proofs based on the scalar source equation and compact exhaustion.
Their computer-assisted examples use separate 320-bit and 384-bit
outward certificates. Independent full-amplitude calculations at 110
and 140 digits corroborate the continuation obstruction and critical
examples, respectively; they do not replace outward bounds. The critical
replay checks 45 values inside the raw intervals, 20 finite-scale maxima,
16 roots and eight wrong-negative-label controls. Exact algebra checks
the derivative reductions. Source definitions, domains, producer hashes
and proof-status matrices are under \nolinkurl{references/e135/},
\nolinkurl{references/e136/} and \nolinkurl{references/e137/}.

\clearpage
\begin{figure}[htbp]
\centering
\includegraphics[width=0.62\textwidth]{../figures/counts_geodesics.pdf}
\caption{The $37$ rotating geodesics from the identity (cross) to
$(39/5,-52/5,3/10)$ (dot), coloured by the period index $n$ of
\eqref{eq:bound}, the number of pendulum periods.}
\label{fig:geodesics}
\end{figure}

\paragraph{Open questions.}
Effective values of the interior and boundary thresholds $R_0$ in
Theorems~\ref{thm:main}--\ref{thm:directions}; finite-target counts at
strict heading in arbitrary spatial directions, beyond the all-radius
half-heading ray, exterior-sector results, the fixed target line in
Proposition~\ref{prop:nine-source-interval}, and the exact catalogues
at headings zero and $\pi$; certified evaluation of further exterior
angular thresholds, uniqueness of their minimizing parameters and the
local singularity types beyond the first threshold at $w=3/4$; continuation of
the first fold window beyond the certified cap in
Corollary~\ref{cor:effective-fold-window}; time comparisons between
arbitrary branches in different positive periods at strict heading, beyond
the complete fixed-line ordering of Corollary~\ref{thm:line-time-order};
time comparisons between separate sublevel components within one signed
period family, beyond Theorem~\ref{thm:general-pair-order};
finite-target comparisons on bounded scaling components, which
Proposition~\ref{prop:dilation-obstruction} shows cannot be covered by
continuation to infinity; explicit starting scales in
Theorems~\ref{thm:dilation-strict}--\ref{thm:dilation-critical}, the
simultaneous critical degeneracy $L=\kappa=0$ (one certified instance,
experiment E138 in the repository, is not treated here), and a complete critical
transition diagram;
effective spatial neighborhoods preserving these counts and orders;
and the incidences of the higher conjugate
sheets \cite{CT} with the fibers counted here.

\paragraph{AI assistance.}
This paper was developed with OpenAI Codex and Claude (Anthropic); the
author is responsible for the content.

\appendix
\section{Substitutions in the proof of Lemma~\ref{lem:zero-order}}
\label{app:subs}

With $s=\sin A$, $c=\cos A$, $d=\sqrt{1-k^2s^2}$ and $h=\sqrt{1-k^2}$,
the closed forms in the proof of Lemma~\ref{lem:zero-order} are
polynomial identities after the substitutions below. They were verified
by computer algebra in the script
\path{research/e33_zero_heading_count.py} of the repository, which also
checks the derivative identities \eqref{eq:zero-restricted},
\eqref{eq:zero-root}, \eqref{eq:zero-time-bridge} and
\eqref{eq:zero-label} modulo the relations $s^2+c^2=1$, $d^2+k^2s^2=1$
and $h^2+k^2=1$.

\emph{Positive amplitude.} With $z=d/e$ and $\tan A=rz$,
\[
 s=\frac{rz}{\sqrt{1+r^2z^2}},\qquad c=\frac1{\sqrt{1+r^2z^2}},\qquad
 e=\frac dz,\qquad k^2=\frac{(1-d^2)(1+r^2z^2)}{r^2z^2},
\]
and $I=d^3F_R/[rz^3\sqrt{1+r^2z^2}]$. The boundary $d^2=c^2$ is
$d^2=1/(1+r^2z^2)$, on which $F_R$ equals the displayed rational
function of $(r,z)$. On the boundary $d^2=4z^2$, with $w=z^2$ and
$\xi=rw$, one has $wF_R=P(w,\xi)$. The Bernstein coefficients are those
of $P_w(1/4-\xi^2,\xi)$, a polynomial of degree four in $\xi$, on
$0\le2\xi\le1$.

\emph{Negative amplitude.} With $\chi=-A$, $z=d/e$ and $\tan\chi=rz$,
\[
 s=-\frac{rz}{\sqrt{1+r^2z^2}},\qquad c=\frac1{\sqrt{1+r^2z^2}},\qquad
 e=\frac dz,\qquad h^2=\frac{d^2(1+r^2z^2)-1}{r^2z^2},
\]
and $I=d^3G/[z^3(1+r^2z^2)^{3/2}]$, with $G$ evaluated at this $h^2$.
In the comparison of $G$ with $B_0$, $q/e=(1+rz^2)/(1-r)$ and
\[
 \frac{B_0}{e^4}=\frac qe\,(1-z^2)-h^2z^2\Bigl(1+\frac qe\Bigr)^2 .
\]
Finally, the lower bound in \eqref{eq:zero-margin} is the identity
$(2n-\tfrac12)\{(2n-\tfrac12)^2-\tfrac74\}-\tfrac34n
=(n-1)(32n^2+8n-3)/4$.

\begin{thebibliography}{99}
\raggedright
\bibitem{MS} I. Moiseev and Y. L. Sachkov.
Maxwell strata in sub-Riemannian problem on the group of motions of a plane.
\emph{ESAIM Control Optim. Calc. Var.} 16 (2010), 380--399.
\href{https://doi.org/10.1051/cocv/2009004}{doi:10.1051/cocv/2009004}.
\bibitem{SC} Y. L. Sachkov.
Conjugate and cut time in the sub-Riemannian problem on the group of motions of a plane.
\emph{ESAIM Control Optim. Calc. Var.} 16 (2010), 1018--1039.
\href{https://doi.org/10.1051/cocv/2009031}{doi:10.1051/cocv/2009031}.
\bibitem{SS} Y. L. Sachkov.
Cut locus and optimal synthesis in the sub-Riemannian problem on the group of motions of a plane.
\emph{ESAIM Control Optim. Calc. Var.} 17 (2011), 293--321.
\href{https://doi.org/10.1051/cocv/2010005}{doi:10.1051/cocv/2010005}.
\bibitem{DBRS} R. Duits, U. Boscain, F. Rossi, and Y. Sachkov.
Association fields via cuspless sub-Riemannian geodesics in SE(2).
\emph{J. Math. Imaging Vision} 49 (2014), 384--417.
\href{https://doi.org/10.1007/s10851-013-0475-y}{doi:10.1007/s10851-013-0475-y}.
\bibitem{LR} A. Lerario and L. Rizzi.
How many geodesics join two points on a contact sub-Riemannian manifold?
\emph{J. Symplectic Geom.} 15 (2017), 247--305.
\href{https://doi.org/10.4310/JSG.2017.v15.n1.a7}{doi:10.4310/JSG.2017.v15.n1.a7}.
\bibitem{WKB} Y. Wang, S. Ku, and A. Bravo-Doddoli.
Metric lines in the special Euclidean group on the plane.
\emph{Involve} 19 (2026), 471--490.
\href{https://doi.org/10.2140/involve.2026.19.471}{doi:10.2140/involve.2026.19.471}.
\bibitem{DLMF} NIST Digital Library of Mathematical Functions,
\href{https://dlmf.nist.gov/}{dlmf.nist.gov}, Sections 19.2, 19.4, 19.5,
19.12, 22.8 and 22.16.
\bibitem{ABB} A. Agrachev, D. Barilari, and U. Boscain.
\emph{A Comprehensive Introduction to Sub-Riemannian Geometry}.
Cambridge Studies in Advanced Mathematics 181, Cambridge University
Press, 2020.
\href{https://doi.org/10.1017/9781108677325}{doi:10.1017/9781108677325}.
\bibitem{Milnor} J. W. Milnor.
\emph{Topology from the Differentiable Viewpoint}.
University Press of Virginia, Charlottesville, 1965, Section 4.
\href{https://people.dm.unipi.it/benedett/MILNOR-TDVPOINT.pdf}{Original edition}.
\bibitem{VDM} L. van den Dries and C. Miller.
Geometric categories and o-minimal structures.
\emph{Duke Math. J.} 84 (1996), 497--540.
\href{https://doi.org/10.1215/S0012-7094-96-08416-1}{doi:10.1215/S0012-7094-96-08416-1}.
\href{https://people.math.osu.edu/miller.1987/newg.pdf}{Authors' revision,
20 February 2001}, Section 2 and statement 4.4.
\bibitem{SUY} K. Saji, M. Umehara, and K. Yamada.
$A_k$ singularities of wave fronts.
\href{https://arxiv.org/abs/0804.0701v4}{arXiv:0804.0701v4},
25 October 2008, Appendix A, Theorem A.1.
\bibitem{CT} I. Moiseev.
All conjugate times of the rotating geodesics in the sub-Riemannian problem on SE(2).
Preprint, 2026.
\bibitem{TF} I. Moiseev.
Moving terminal angles and exponentially small transitions in the sub-Riemannian problem on SE(2).
In preparation.
\end{thebibliography}
\end{document}
